\chapter{Stochastic and Financial Foundations} \label{ch:cap2} 

% \section{Preliminaries}

% We shall assume some prior knowledge on Probability Theory, which can be found in any book on the matter. As for a personal preference, Billingsley's oeuvre (\textbf{reference}).

% Let (Omega, F, P) be a probability space. It is known that we can assume the probability space to be \textit{complete}, that is, if $A \in F$ and $P(A) = 0$, then every $B \subset A$ also belongs to $F$ and $P(B) = 0.$ We will always assume that our probability space is complete.


% % Throughout the text, we will mainly work with continuous stochastic processes (appropriately defined in \textbf{REFERENCE}). However, in order to give a wide setting for these processes, let $\mathbb{T}$ be the time that indexes a process.


% \section{Elementary Processes}

% The first key concept we need to address is that of stochastic process.

% \begin{definition}
%     Let (Omega, F, P) a probability space. Let $\mathbb{T} \in \{\mathbb{N}, \mathbb{Z}, [0,T], [0, \infty) , \mathbb{R}\}$ for $T > 0.$ We say $X(t,\omega)$ is a stochastic process ...
% \end{definition}







We assume the reader is familiar with the foundations of Probability Theory. 
Standard references include, for example, Billingsley's classic text 
\cite{Biggy}, as well as the treatments in \cite{durrett2019probability} and \cite{kallenberg2021foundations}. 
Nevertheless, we briefly recall the basic probabilistic framework in order to fix 
notation and ensure clarity.

\medskip

Throughout, we work on a probability space 
\((\Omega, \mathcal{F}, \mathbb{P})\), where:
\begin{itemize}
    \item \(\Omega\) is the \emph{sample space}, representing the collection of 
    all possible outcomes of the experiment under consideration.
    \item \(\mathcal{F}\) is a \(\sigma\)-algebra of subsets of \(\Omega\), whose 
    elements are called \emph{events}. By definition, \(\mathcal{F}\) is closed 
    under countable unions, countable intersections, and complements, thereby 
    providing the natural domain for the probability measure.
    \item \(\mathbb{P}\) is a probability measure on \((\Omega, \mathcal{F})\). The measure-theoretic formulation of 
    probability ensures that both discrete and continuous models can be treated 
    within a unified framework.
\end{itemize}

Without loss of generality, we shall assume that the probability space is 
\emph{complete}. That is, if \(A \in \mathcal{F}\) satisfies \(\mathbb{P}(A) = 0\), 
then every subset \(B \subseteq A\) also belongs to \(\mathcal{F}\), and moreover 
\(\mathbb{P}(B) = 0\). This technical assumption rules out ambiguities caused by 
“missing” null sets, and will be taken for granted throughout.

\medskip

\begin{definition}[Random Variable] \label{def:random-variable}
A \textbf{random variable} is a measurable function 
\[
X : (\Omega, \mathcal{F}) \to (E, \mathcal{E}),
\]
where \((E, \mathcal{E})\) is a measurable space. 
In most of our applications, the target space will be 
\((\mathbb{R}, \mathcal{B}(\mathbb{R}))\), where 
\(\mathcal{B}(\mathbb{R})\) denotes the Borel \(\sigma\)-algebra. 
\end{definition}

Intuitively, a random variable associates to each elementary outcome 
\(\omega \in \Omega\) a numerical value \(X(\omega)\). 
Formally, the measurability condition ensures that events of the form 
\(\{ \omega \in \Omega : X(\omega) \in B \}\) belong to \(\mathcal{F}\) 
whenever \(B \in \mathcal{B}(\mathbb{R})\), so that their probabilities 
are well-defined.

\medskip

This general framework provides the foundation for the study of stochastic processes, 
which may be regarded as families of random variables indexed by time or some other 
parameter. Before turning to them, we emphasize that the measure-theoretic language 
is not merely a technical convenience: it is indispensable for developing modern 
stochastic calculus, where delicate issues such as path regularity, filtrations, and 
martingale properties require the full power of measure theory.
\section{Probabilistic Preliminaries}

The notion of \emph{stochastic process} generalizes that of a single random variable by allowing evolution in an external parameter, usually interpreted as time. To accommodate both discrete-time and continuous-time settings in a unified framework, we let \(\mathbb{T}\) denote the \textbf{index set}, representing time. Typically, \(\mathbb{T}\) is a subset of \(\mathbb{R}\), such as:
\[
\mathbb{T} \in \{ \mathbb{N}, \mathbb{N}_0, \mathbb{Z}, [0, T], [0, \infty), \mathbb{R} \},
\]
depending on the application.

\begin{definition}
Let \((\Omega, \mathcal{F}, \mathbb{P})\) be a complete probability space, and let \(\mathbb{T} \) be a time index set. A \textbf{stochastic process} is a collection of random variables \(\{ X_t : t \in \mathbb{T} \}\) defined on \((\Omega, \mathcal{F}, \mathbb{P})\) with values in a measurable space \((E, \mathcal{E})\), such that the map
\[
X : \mathbb{T} \times \Omega \to E, \quad (t, \omega) \mapsto X_t(\omega)
\]
is jointly measurable with respect to \(\mathcal{B}(\mathbb{T}) \otimes \mathcal{F}\). In particular:
\begin{itemize}
    \item For each fixed \(t \in \mathbb{T}\), the function \(\omega \mapsto X_t(\omega)\) is a random variable on \((\Omega, \mathcal{F}, \mathbb{P})\).
    \item For each fixed \(\omega \in \Omega\), the function \(t \mapsto X_t(\omega)\) is called the \emph{sample path} or \emph{trajectory} of the process.
\end{itemize}
We may denote the process by \(X = \{X_t\}_{t \in \mathbb{T}}\), and abbreviate \(X_t(\omega)\) as \(X(t)\) or simply \(X_t\) when context allows. We say $X$ is \textit{discrete} whenever $\mathbb{T}$ is $\mathbb{N}$ or $\mathbb{Z}$ and \textit{continuous} if $\mathbb{T}$ is $[0,T]$ for some $T>0$ or $[0,\infty).$
\end{definition}

This formulation includes many important examples such as Markov chains, Poisson processes, and Brownian motion. Note that the nature of the index set \(\mathbb{T}\) often dictates the regularity assumptions or path properties we may require (e.g., right-continuity, càdlàg, continuity in probability, etc.).  Analyzing the dynamic evolution of these processes requires us to quantify their expected future behavior given currently available information. This idea is contained in the \textbf{conditional expectation}. Its existence relies on the following profound result:

\begin{theorem}[Radon--Nikodym]\label{thm:radon-nikodym}
    Let \(\nu\) and \(\lambda\) be two \(\sigma\)-finite measures on the same measurable space \((\Omega, \mathcal{F})\). Suppose that \(\nu\) is \emph{absolutely continuous} with respect to \(\lambda\), denoted \(\nu \ll \lambda\); that is,
\begin{equation} \label{eq:abscont}
\lambda(A) = 0 \quad \Rightarrow \quad \nu(A) = 0, \quad \text{for all } A \in \mathcal{F}.
\end{equation}

\noindent Then there exists a measurable function \(f : \Omega \to [0, \infty)\), unique up to $\lambda$-null sets, such that for all \(A \in \mathcal{F}\),
\begin{equation} \label{eq:radonnikodym}
\nu(A) = \int_A f \, d\lambda.
\end{equation}
This function \(f\) is called the \emph{Radon--Nikodym derivative} of \(\nu\) with respect to \(\lambda\), and is denoted by $f = d\nu/d\lambda$. \hfill $\blacksquare$
\end{theorem}

The letter $\mu$ is deliberately avoided here. From Section \ref{sec:stochasticcalculus} onwards it will denote a drift coefficient and nothing else; abstract measures are $\nu$ and $\lambda$ throughout.

We shall need the statement in slightly greater generality than above, since the measure we are about to construct is not positive.

\begin{corollary}[Signed version]\label{cor:radon-nikodym-signed}
Let $\nu$ be a finite \emph{signed} measure on $(\Omega,\mathcal{F})$ and $\lambda$ a $\sigma$-finite positive measure with $\nu \ll \lambda$. Then there exists $f \in L^1(\Omega,\mathcal{F},\lambda)$, real-valued and unique up to $\lambda$-null sets, satisfying \eqref{eq:radonnikodym}.
\end{corollary}

\textbf{\textit{Proof:}} Apply the Jordan decomposition $\nu = \nu^+ - \nu^-$ into mutually singular finite positive measures. Both $\nu^{\pm} \ll \lambda$, so Theorem \ref{thm:radon-nikodym} yields non-negative densities $f^{\pm}$, and $f := f^+ - f^-$ works; integrability follows from $\int (f^+ + f^-)\,d\lambda = |\nu|(\Omega) < \infty$. \hfill $\blacksquare$

As a consequence, if \(X \in L^1(\Omega, \mathcal{F}, \mathbb{P})\) (that is, $\int |X| \, d\mathbb{P}$ is finite) and letting \(\mathcal{G} \subseteq \mathcal{F}\) be a sub-\(\sigma\)-algebra, the \emph{conditional expectation} of \(X\) given \(\mathcal{G}\), denoted \(\mathbb{E}[X \mid \mathcal{G}]\), is the unique (up to \(\mathbb{P}\)-a.e. equality) \(\mathcal{G}\)-measurable function satisfying:
\begin{equation}
\int_G \mathbb{E}[X \mid \mathcal{G}] \, d\mathbb{P} = \int_G X \, d\mathbb{P}, \quad \text{for all } G \in \mathcal{G}.
\end{equation}

This definition is structurally equivalent to the Radon--Nikodym theorem. Define the set function \(\nu : \mathcal{G} \to \mathbb{R}\) by
\begin{equation}
\nu(G) := \int_G X \, d\mathbb{P}, \quad \text{for } G \in \mathcal{G}.
\end{equation}
Since \(X\) is integrable, \(\nu\) is a finite \emph{signed} measure on \((\Omega, \mathcal{G})\), and it is absolutely continuous with respect to \(\mathbb{P}|_{\mathcal{G}}\), the restriction of \(\mathbb{P}\) to \(\mathcal{G}\). It is Corollary \ref{cor:radon-nikodym-signed}, rather than Theorem \ref{thm:radon-nikodym} itself, that applies here: $X$ need not be non-negative, so neither is $\nu$, and the density produced must be allowed to change sign. With that proviso (see \cite{Biggy} or \cite{RoydenFitzpatrick2010}), there exists a \(\mathcal{G}\)-measurable function, namely \(\mathbb{E}[X \mid \mathcal{G}]\), such that:
\begin{equation}
\int_G \mathbb{E}[X \mid \mathcal{G}] \, d\mathbb{P} = \int_G X \, d\mathbb{P}, \quad \text{for all } G \in \mathcal{G}.
\end{equation}

\begin{proposition}[Properties of conditional expectation]\label{def:conditional-expectation}
Let $(\Omega, \mathcal{F}, \mathbb{P})$ be a probability space and let $\mathcal{G}$ be a sub-$\sigma$-field of $\mathcal{F}$. Then, the following hold almost surely:

\begin{itemize}
    \item[1.] $\mathbb{E}\!\left[\,\mathbb{E}[X \mid \mathcal{G}]\,\right] = \mathbb{E}[X]$. Hence the conditional expectation $\mathbb{E}[X \mid \mathcal{G}]$ and $X$ have the same expectation.

    \item[2.] If $X$ is $\mathcal{G}$-measurable, then $\mathbb{E}[X \mid \mathcal{G}] = X$.

    \item[3.] If $X$ and $\mathcal{G}$ are independent, then $\mathbb{E}[X \mid \mathcal{G}] = \mathbb{E}[X]$   ($\{X \in U\}$ and $A$ are independent events for any Borel subset $U \subseteq \mathbb{R}$ and $A \in \mathcal{G}$).


    \item[4.] (\textbf{Taking out what is known}) If $X \in L^1(\Omega)$, $Y$ is $\mathcal{G}$-measurable and $\mathbb{E}[\,|XY|\,] < \infty$, then
    $
        \mathbb{E}[XY \mid \mathcal{G}] = Y \, \mathbb{E}[X \mid \mathcal{G}]. $
    

    \item[5.] If $\mathcal{H}$ is a sub-$\sigma$-field of $\mathcal{G}$, then
    $
        \mathbb{E}[X \mid \mathcal{H}] = \mathbb{E}\!\left[\,\mathbb{E}[X \mid \mathcal{G}] \,\middle|\, \mathcal{H}\,\right]. $
    


    \item[6.] If $X, Y \in L^{1}(\Omega)$ and $X \leq Y$, then
    $
        \mathbb{E}[X \mid \mathcal{G}] \leq \mathbb{E}[Y \mid \mathcal{G}]$ and, therefore, $\left| \mathbb{E}[X \mid \mathcal{G}] \right| \leq \mathbb{E}[\,|X| \mid \mathcal{G}].$


    \item[7.] For all $a,b \in \mathbb{R}$ and $X,Y \in L^{1}(\Omega)$,
    \[
        \mathbb{E}[aX + bY \mid \mathcal{G}] = a \, \mathbb{E}[X \mid \mathcal{G}] + b \, \mathbb{E}[Y \mid \mathcal{G}].
    \]
    \textit{Remark:} It follows that the conditional expectation $\mathbb{E}[\cdot \mid \mathcal{G}]$ is a bounded linear operator from $L^{1}(\Omega,\mathcal{F})$ into $L^{1}(\Omega,\mathcal{G})$.

    \item[8.] (\textbf{Conditional Jensen's inequality}) If $\phi : \mathbb{R} \to \mathbb{R}$ is convex and both $X$ and $\phi(X)$ are integrable, then
    \[
        \phi\bigl(\mathbb{E}[X \mid \mathcal{G}]\bigr) \leq \mathbb{E}[\phi(X) \mid \mathcal{G}].
    \]

    \item[9.] (\textbf{$L^2$ projection}) If $X \in L^2(\Omega,\mathcal{F},\mathbb{P})$, then $\mathbb{E}[X \mid \mathcal{G}]$ is the orthogonal projection of $X$ onto the closed subspace $L^2(\Omega,\mathcal{G},\mathbb{P})$. Equivalently, it is the unique $\mathcal{G}$-measurable $Y$ minimising $\mathbb{E}[(X-Y)^2]$, and
    \[
        \mathbb{E}\bigl[(X - \mathbb{E}[X\mid\mathcal{G}])\,Z\bigr] = 0 \quad \text{for every } Z \in L^2(\Omega,\mathcal{G},\mathbb{P}).
    \] \hfill $\blacksquare$
\end{itemize}
\end{proposition}

The conditional versions of Fatou's lemma and of the monotone and dominated convergence theorems also hold, in exactly the form one would expect; since we shall use them only in passing, they are collected in Appendix \ref{app:auxiliary}.

Two of the items above deserve emphasis, since they will be used more than all the others combined. Property 8 is what guarantees that $X^2$ is a submartingale whenever $X$ is a martingale --- a fact invoked twice in what follows, in Subsection \ref{subsec:martingales} and again in the Doob-Meyer decomposition --- and it is the reason convexity is such a persistent theme in option pricing.

Property 9 deserves more than emphasis. It says that conditioning is nothing more exotic than orthogonal projection in a Hilbert space, and once this is seen, a good deal of what follows stops looking like a series of coincidences. The Itô isometry of Theorem \ref{thm:itoisometry} will be an isometry between two Hilbert spaces; the density argument that extends the stochastic integral from step processes to all of $L^2_{\mathrm{ad}}$ is the standard extension of a bounded operator from a dense subspace; and the Martingale Representation Theorem of Subsection \ref{subsec:measurechanges} will say that a certain family of stochastic integrals is not merely orthogonal but \emph{complete}. The geometry is the same throughout.


There is one final structural component that needs to be described. In the study of stochastic processes, it is essential to formalize how information accumulates over time. This is done through the concept of a \textbf{filtration}, which models the \emph{evolution of knowledge} or \emph{information available} up to each time \(t\).



\begin{definition}[Filtration]
Let \((\Omega, \mathcal{F}, \mathbb{P})\) be a probability space and let \(\mathbb{T}\) be a time index set. A \textbf{filtration} is a family \(\{\mathcal{F}_t\}_{t \in \mathbb{T}}\) of sub-\(\sigma\)-algebras of \(\mathcal{F}\) satisfying:
\begin{equation}
\mathcal{F}_s \subseteq \mathcal{F}_t \subseteq \mathcal{F} \quad \text{for all } s \leq t.
\end{equation}
That is, the filtration is \emph{increasing} in \(t\). We interpret \(\mathcal{F}_t\) as the collection of events whose outcomes are known or observable up to  \(t\). A stochastic process \(X = \{X_t\}_{t \in \mathbb{T}}\) is said to be \textbf{adapted} to the filtration \(\{\mathcal{F}_t\}\) if, for each \(t \in \mathbb{T}\), the random variable \(X_t\) is measurable with respect to \(\mathcal{F}_t\).  We say that $(\Omega, \mathcal F, \{\mathcal F_t\}, \mathbb P)$ is a \textbf{filtered probability space}.
\end{definition}

Bare adaptedness turns out to be slightly too weak for comfort. Two further requirements are so universally imposed that they have acquired a name of their own, and we impose them here once and for all rather than invoking them piecemeal later.

\begin{definition}[The usual conditions]\label{def:usualconditions}
A filtered probability space $\FPSpace$ is said to satisfy the \textbf{usual conditions} if
\begin{itemize}
    \item[i)] $\mathcal{F}_0$ contains every $\mathbb{P}$-null set of $\mathcal{F}$ (\emph{completeness}); and
    \item[ii)] the filtration is \emph{right-continuous}, that is,
    \begin{equation}\label{eq:rightcontinuous}
    \mathcal{F}_t = \mathcal{F}_{t^+} := \bigcap_{n\geq 1}\mathcal{F}_{t + 1/n}, \qquad t \in [0,T).
    \end{equation}
\end{itemize}
\end{definition}

\begin{convention}\label{conv:usualconditions}
\textbf{From this point onwards, every filtered probability space is assumed to satisfy the usual conditions,} and this will not be restated.
\end{convention}

Neither requirement is cosmetic. Completeness is what allows us to modify a process on a null set --- as we shall do whenever we choose a continuous version of an object defined only up to indistinguishability --- without leaving the filtration. Right-continuity is what makes first hitting times of open sets stopping times, and it is used explicitly in Proposition \ref{prop:stoppingtimechar} below and again in the Doob-Meyer decomposition of Subsection \ref{subsec:measurechanges}. Neither is automatic for a natural filtration $\mathcal{F}^X_t$, but both can always be arranged: one replaces $\mathcal{F}^X_t$ by the $\mathbb{P}$-completion of $\mathcal{F}^X_{t^+}$, the so-called \emph{augmented} filtration, and for the filtrations generated by the processes we shall study this augmentation changes nothing of substance.

\begin{remark}[Modification versus indistinguishability]\label{rem:modification}
A related piece of hygiene, which we shall need whenever a theorem produces ``a continuous version'' of a process. Two processes $X$ and $Y$ on the same space are said to be \emph{modifications} of each other if $\mathbb{P}(X_t = Y_t) = 1$ for every fixed $t$, and \emph{indistinguishable} if $\mathbb{P}(X_t = Y_t \text{ for all } t) = 1$. The second is strictly stronger: the exceptional null set is allowed to depend on $t$ in the first case, and a union of uncountably many null sets need not be null. The two notions do coincide when both processes have right-continuous paths, which is the situation we shall almost always be in --- but the coincidence is a small theorem, not a definition, and we shall be careful to say which of the two we mean.
\end{remark}


Recall that if \(\chi\) is a real-valued random variable on a probability space \((\Omega, \mathcal{F}, \mathbb{P})\), then the \(\sigma\)-algebra generated by \(\chi\) is denoted \(\sigma(\chi)\) and defined as:
\begin{equation}
\sigma(\chi) := \{ A \in \mathcal{F} : A = \chi^{-1}(B) \text{ for some } B \in \mathcal{B}(\mathbb{R}) \}.
\end{equation}


\begin{example}[Natural Filtration]
Given a stochastic process \(X = \{X_t\}_{t \in \mathbb{T}}\), the \textbf{natural filtration} generated by \(X\) is defined by:
\begin{equation}
\mathcal{F}_t^X := \sigma(X_s : s \leq t),
\end{equation}
the smallest \(\sigma\)-algebra with respect to which all variables \(X_s\) for \(s \leq t\) are measurable. This filtration encodes exactly the information revealed by observing the process up to time \(t\).
\end{example}

\begin{definition}[Stopping time]\label{stoppingTime}
Let $(\Omega,\mathcal{F}, \mathbb{P})$ a probability space and a filtration $\{\mathcal{F}_t\}.$ A mapping $\tau:\Omega \rightarrow [0,T]$ is  a \textit{stopping time} if, for each $t \in [0,T] (\textit{respect. } t\in[0,\infty)])$, $\{\omega  :  \tau(\omega) \leq t\} \in \mathcal{F}_t.$ 
\end{definition}


\begin{proposition}\label{prop:stoppingtimechar}
 Let \(\Filt\) be a right-continuous filtration. A random variable \(\tau : \Omega \to [0, T]\) is a stopping time with respect to \(\Filt\) if and only if  
\[
\{ \, \omega \in \Omega : \tau(\omega) < t \, \} \in \mathcal{F}_t, 
\quad \text{for all } t \in [0, T].
\]
\end{proposition}

The implication from left to right requires nothing, since $\{\tau < t\} = \bigcup_n \{\tau \leq t - 1/n\}$; it is the converse that consumes the right-continuity of the filtration, via $\{\tau \leq t\} = \bigcap_n \{\tau < t + 1/n\} \in \bigcap_n \mathcal{F}_{t+1/n} = \mathcal{F}_t$. This is the first of several places where the usual conditions of Definition \ref{def:usualconditions} do real work rather than merely tidy up.

\textbf{\textit{Proof:}} (...) \hfill $\blacksquare$

\begin{example}\label{example_stoppingTimes}
    Let $X$ be a measurable adapted process on $[0,T]$ with $\int_0^T |X_s|^2\,ds < \infty$ almost surely. For $n \in \mathbb{N}$ set
    \[
    A_n(\omega) =  \Bigl\{ t \in [0,T] : \int_0^t |X_s(\omega)|^2\,ds > n\Bigr\},
    \]
    and let $\tau_n(\omega) = \inf A_n(\omega)$ if $A_n(\omega) \neq \emptyset$, and $\tau_n(\omega) = T$ otherwise. Since $t \mapsto \int_0^t |X_s|^2 ds$ is continuous and non-decreasing, the threshold $n$ is exceeded strictly before $t$ if and only if it is exceeded at some rational time strictly before $t$, so that for $t\in (0,T]$
    \begin{equation}
        \{\tau_n < t\} = \bigcup_{\substack{0<r<t \\ r \in \mathbb{Q}}} \Bigl\{\int_0^r |X_s|^2\,ds > n\Bigr\} \in \mathcal{F}_t,
    \end{equation}
    whereas $\{\tau_n < 0\} = \emptyset \in \mathcal{F}_0$. By Proposition \ref{prop:stoppingtimechar}, $\tau_n$ is a stopping time.

    This sequence is not an idle illustration. It is exactly the localising sequence that will allow us, in Subsection \ref{subsec:itointegral}, to define the stochastic integral for integrands that are square-integrable pathwise but not in expectation, and it is the prototype of every localisation argument in this chapter.
\end{example}

\begin{example}[First hitting time]
    Let $X_t$ be a right continuous process adapted to $\{\mathcal{F}_t\}.$ If $U$ is an open set of $\mathbb{R}$, then the \textit{first hitting time}, $\tau_U$ is defined as
    \begin{equation}
        \tau_U(\omega) = \inf\{t>0 : X_t(\omega) \in U\}
    \end{equation}

     is a stopping time.
    
    % since

    % \[
    % \{\tau_U \leq t\} = \cap_n \{\tau_U < t + \frac{1}{n}\} = \cap_n \cup_{r\in\mathbb{Q} \cap (-\infty,t+1/n)} \{X_r \in U\} \in
    % \]
\end{example}

%     \subsection{Fundamental Distributions}


% Before presenting some cornerstone distributions, recall that the \textbf{characteristic function} of a (real) random variable $X$ is defined by
% \begin{equation}
% \varphi_X(u) := \mathbb{E}[e^{i u X}], \quad u \in \mathbb{R}.
% \end{equation}
% It is the Fourier transform of the probability distribution of $X$, \hl{and it uniquely determines the law of $X$. }

% Some basic properties of this new variable are:
% \begin{itemize}
%   \item $\varphi_X(0) = 1$.
%   \item $|\varphi_X(u)| \leq 1$ for all $u \in \mathbb{R}$.
%   \item $\varphi_X(-u) = \overline{\varphi_X(u)}$ (Hermitian symmetry in the complex case).
%   \item $\varphi_{aX + b}(u) = e^{i u b} \varphi_X(a u)$ for all $a, b \in \mathbb{R}$.
%   \item If $X$ and $Y$ are independent, then $\varphi_{X + Y}(u) = \varphi_X(u) \varphi_Y(u)$.
%   %\item $\varphi_X$ is uniformly continuous on $\mathbb{R}$.
% \end{itemize}

% If $X$ has finite $n$-th moment, i.e.\ $\mathbb{E}[|X|^n] < \infty$, then $\varphi_X$ is $n$ times differentiable and
% \[
% \varphi_X^{(k)}(0) = i^k \, \mathbb{E}[X^k], \quad \text{for } 1 \leq k \leq n.
% \]
% Conversely, the existence of continuous derivatives of $\varphi_X$ on a nieghborhood of $0$ imply that the corresponding moments exist.

% \begin{theorem}
% Let $X$ be a random variable with distribution function $F_X$. If $\varphi_X$ is integrable, then $F_X$ admits a continuous density $f_X$ given by
% \[
% f_X(x) = \frac{1}{2\pi} \int_{\mathbb{R}} e^{-i u x} \varphi_X(u) \, du.
% \]
% \end{theorem}

% In the general case, without assuming a density, the distribution function $F_X$ can still be recovered via the \emph{Lévy inversion formula}:
% \begin{equation}\label{levyInversionFormula}
% F_X(b) - F_X(a) = \lim_{T \to \infty} \frac{1}{2\pi} \int_{-T}^T \frac{e^{-i a u} - e^{-i b u}}{i u} \, \varphi_X(u) \, du.
% \end{equation}



% We deduce that if two random variables $X$ and $Y$ have the same characteristic function, i.e.\ $\varphi_X(u) = \varphi_Y(u)$ for all $u \in \mathbb{R}$, then $X$ and $Y$ have the same distribution. That is,
% \[
% \varphi_X = \varphi_Y \quad \Longleftrightarrow \quad \mathbb{P}_X = \mathbb{P}_Y.
% \]
% \begin{theorem}[Continuity Theorem, Lévy]\label{theorem:LevyContinuity}
% Let $\{X_n\}_{n \in \mathbb{N}}$ be a sequence of random variables with characteristic functions $\varphi_{X_n}$, and let $X$ be a random variable with characteristic function $\varphi_X$. Then the following are equivalent:
% \begin{itemize}
%   \item $X_n \xrightarrow{d} X$ (convergence in distribution),
%   \item $\varphi_{X_n}(u) \to \varphi_X(u)$ for all $u \in \mathbb{R}$, and $\varphi_X$ is continuous at $u = 0$.
% \end{itemize}
% \end{theorem}

% We now give some description of certain processes  with widespread name and usefulness. The first place must be given to the normal distribution.


% \medskip

% \begin{example}[Normal Distribution]
% Let $\mu \in \mathbb{R}$ and $\sigma > 0$, and define the real-valued random variable $X \colon \Omega \to \mathbb{R}$ with distribution $X \sim \mathcal{N}(\mu, \sigma^2),$
% that is, $X$ has the probability density function
% \[
% f_X(x) = \frac{1}{\sqrt{2\pi \sigma^2}} \exp\left( -\frac{(x - \mu)^2}{2\sigma^2} \right), \quad x \in \mathbb{R}.
% \]
% Then $X$ is said to be normally distributed with mean $\mu$ and variance $\sigma^2$, so that a normal distribution is characterized by its two first moments.
% \end{example}

% The normal distribution is fundamental in probability theory and statistics. It appears as the limiting distribution in the Central Limit Theorem and models many naturally occurring phenomena. Its characteristic function adopts the form $\varphi_X(u) = \exp\left( i u \mu - \frac{1}{2} \sigma^2 u^2 \right).$ Its resemblance to the \textit{heat kernel} (see \cite{Evans2010}) is not a coincidence and will be explored in the next section.

% Let $X \sim \mathcal{N}(\mu_X, \sigma_X^2)$, $Y \sim \mathcal{N}(\mu_Y, \sigma_Y^2)$ be two normal random variables. The following are easily verified through Theorem \ref{theorem:LevyContinuity}:

% \begin{itemize}
%     \item For $a, b \in \mathbb{R}, aX + bY \sim \mathcal{N}(a\mu_X+b\mu_Y, a^2\sigma_X^2+b^2 \sigma_Y^2)$. For appropriateness when working with constant variables, we may write $ Z \sim \mathcal{N}(\mu, 0)$ for a constant random variable that only takes the value $\mu$. With this notation, the former idendity is still fulfilled when constant and normal variables are put together.
%     \item Its characteristic function is entire analytic.
%     \item Normal distribution maximizes \textbf{entropy} among distributions for fixed first two moments. This makes it the \textit{least informative} among these distributions.
%     \item If $\{X _n \sim \mathcal{N}(\mu_n, \sigma_n^2)\}_{n=1}^\infty$ is a sequence of normal random variables, then $X_n \rightarrow^d X \sim \mathcal{N}(\mu, \sigma^2)$, with $\sigma > 0$, if, and only if, $\mu_n \rightarrow^{n\rightarrow \infty}$ and $\sigma_n^2 \rightarrow ^{n\rightarrow \infty} \sigma^2.$
%     \item In general, two independent random variables have null correlation. In the special case of normal variables, their correlation being zero implies their independence.
    
    
% \end{itemize}

% However, the most sensible reason to why normal distribution is so important relies on the following fundamental result, in a relatively weak form:

% \begin{theorem}[Central Limit Theorem]
%     Let $\{X_n\}_{n=1}^\infty$ a sequence of independent, identically distributed random variables (hereafter, \emph{iidrv}) with finite first two moments, $\mu$ and $\sigma$ respectively. Then,

%     $$
%     \frac{\sum_{k=1}^nX_k-\mu}{\sqrt{n}} \rightarrow^d \mathcal{N}(\sigma^2).
%     $$
% \end{theorem}

% It is hard to stress enough the importance and implications of the normal distribution throughout mathematics. Some insight on the multivariate case shall also be provided.

\subsection{Fundamental Distributions}

Before presenting some cornerstone distributions, recall that the \textbf{characteristic function} of a (real) random variable $X$ is defined by
\begin{equation}
\varphi_X(u) := \mathbb{E}[e^{i u X}], \quad u \in \mathbb{R}.
\end{equation}
It is the Fourier transform of the probability distribution of $X$, and it uniquely determines the law of $X$.

Some basic properties of the characteristic function are:
\begin{itemize}
  \item $\varphi_X(0) = 1$.
  \item $|\varphi_X(u)| \leq 1$ for all $u \in \mathbb{R}$.
  \item $\varphi_X(-u) = \overline{\varphi_X(u)}$ (Hermitian symmetry in the complex case).
  \item $\varphi_{aX + b}(u) = e^{i u b} \varphi_X(a u)$ for all $a, b \in \mathbb{R}$.
  \item If $X$ and $Y$ are independent, then $\varphi_{X + Y}(u) = \varphi_X(u) \varphi_Y(u)$.
\end{itemize}

If $X$ has finite $n$-th moment, i.e.\ $\mathbb{E}[|X|^n] < \infty$, then $\varphi_X$ is $n$ times differentiable and
\begin{equation}
\varphi_X^{(k)}(0) = i^k \, \mathbb{E}[X^k], \quad \text{for } 1 \leq k \leq n.
\end{equation}
Conversely, the existence of continuous derivatives of $\varphi_X$ on a neighborhood of $0$ implies that the corresponding moments exist.

\begin{theorem}
Let $X$ be a random variable with distribution function $F_X$. If $\varphi_X$ is integrable, then $F_X$ admits a continuous density $f_X$ given by
\begin{equation}
f_X(x) = \frac{1}{2\pi} \int_{\mathbb{R}} e^{-i u x} \varphi_X(u) \, du.
\end{equation}
\end{theorem}

In the general case, without assuming a density, the distribution function $F_X$ can still be recovered via the \emph{Lévy inversion formula}: for any $\alpha < \beta$ that are \emph{continuity points} of $F_X$,
\begin{equation}\label{levyInversionFormula}
F_X(\beta) - F_X(\alpha) = \lim_{R \to \infty} \frac{1}{2\pi} \int_{-R}^R \frac{e^{-i \alpha u} - e^{-i \beta u}}{i u} \, \varphi_X(u) \, du.
\end{equation}
The restriction to continuity points is necessary and not a technicality: at a jump of $F_X$ the right-hand side converges to the average of the left and right limits rather than to $F_X$ itself. Since $F_X$ is monotone it has at most countably many discontinuities, so continuity points are dense and the formula determines $F_X$ everywhere by right-continuity. Note also that the truncation parameter $R$ is unrelated to the time horizon $T$ fixed in the notation table; the symbol has been changed to avoid the collision.

We deduce that if two random variables $X$ and $Y$ have the same characteristic function, i.e.\ $\varphi_X(u) = \varphi_Y(u)$ for all $u \in \mathbb{R}$, then $X$ and $Y$ have the same distribution. \hfill $\blacksquare$

\begin{theorem}[Lévy's Continuity Theorem]\label{theorem:LevyContinuity}
Let $\{X_n\}_{n \in \mathbb{N}}$ be a sequence of random variables with characteristic functions $\varphi_{X_n}$.
\begin{itemize}
  \item If $X_n \xrightarrow{d} X$ for some random variable $X$, then $\varphi_{X_n}(u) \to \varphi_X(u)$ for every $u \in \mathbb{R}$.
  \item Conversely, suppose $\varphi_{X_n}(u) \to \varphi(u)$ for every $u \in \mathbb{R}$, where $\varphi : \mathbb{R}\to\mathbb{C}$ is \emph{some} function continuous at $u=0$. Then $\varphi$ is the characteristic function of a random variable $X$, and $X_n \xrightarrow{d} X$. \hfill $\blacksquare$
\end{itemize}
\end{theorem}

The asymmetry between the two halves is deliberate and is where the theorem earns its keep. In the converse direction one is \emph{not} allowed to assume that the limit $\varphi$ is a characteristic function --- that is the conclusion, not a hypothesis. Continuity at the origin is what rules out mass escaping to infinity, and it is a genuine restriction on $\varphi$; note that if one already knew $\varphi$ to be a characteristic function the requirement would be vacuous, since characteristic functions are automatically uniformly continuous. It is precisely this form of the statement that makes the Central Limit Theorem below a two-line computation: one expands $\varphi_{X_n}$, observes that the limit is $e^{-u^2/2}$, and only then concludes that the limit is Gaussian.

We now describe certain distributions of widespread importance and usefulness. The first place must be given to the normal distribution.

\medskip

\begin{example}[Normal Distribution]
Let $\mu \in \mathbb{R}$ and $\sigma > 0$, and define the real-valued random variable $X \colon \Omega \to \mathbb{R}$ with distribution $X \sim \mathcal{N}(\mu, \sigma^2),$
that is, $X$ has the probability density function
\begin{equation}
f_X(x) = \frac{1}{\sqrt{2\pi \sigma^2}} \exp\left( -\frac{(x - \mu)^2}{2\sigma^2} \right), \quad x \in \mathbb{R}.
\end{equation}
Then $X$ is said to be normally distributed with mean $\mu$ and variance $\sigma^2$; thus, a normal distribution is characterized by its first two moments.
\end{example}

The normal distribution is fundamental in probability theory and statistics. It appears as the limiting distribution in the Central Limit Theorem and models many naturally occurring phenomena. Its characteristic function adopts the form
\begin{equation}
\varphi_X(u) = \exp\left( i u \mu - \tfrac{1}{2} \sigma^2 u^2 \right).
\end{equation}
Its resemblance to the \textit{heat kernel} (see \cite{Evans2010}) is not a coincidence and will be explored in the next section.

Let $X \sim \mathcal{N}(\mu_X, \sigma_X^2)$ and $Y \sim \mathcal{N}(\mu_Y, \sigma_Y^2)$ be two \emph{independent} normal random variables. The following are easily verified through Theorem \ref{theorem:LevyContinuity}:

\begin{itemize}
    \item For $a, b \in \mathbb{R}$, $aX + bY \sim \mathcal{N}(a\mu_X+b\mu_Y,\, a^2\sigma_X^2+b^2 \sigma_Y^2)$. The independence hypothesis is essential here and cannot be weakened to ``both marginals are normal'': see Remark \ref{rem:gaussiantraps} below. If $X$ and $Y$ are jointly normal with correlation $\rho$, the same statement holds with variance $a^2\sigma_X^2+b^2\sigma_Y^2+2ab\rho\,\sigma_X\sigma_Y$. For convenience when working with constant random variables, we may write $Z \sim \mathcal{N}(\mu, 0)$ for a constant random variable that only takes the value $\mu$. With this notation, the former identity is still fulfilled when constant and normal variables are combined.
    \item Its characteristic function is entire analytic.
    \item The normal distribution maximizes \textbf{entropy} among distributions with fixed first two moments. This makes it the \textit{least informative} among these distributions.
    \item If $\{X _n \sim \mathcal{N}(\mu_n, \sigma_n^2)\}_{n=1}^\infty$ is a sequence of normal random variables, then $X_n \xrightarrow{d} X \sim \mathcal{N}(\mu, \sigma^2)$, with $\sigma > 0$, if and only if $\mu_n \to \mu$ and $\sigma_n^2 \to \sigma^2$ as $n \to \infty$.
    \item In general, two independent random variables have zero correlation. For normal variables the converse holds too, \emph{provided the pair is jointly normal}; marginal normality alone is not enough.
\end{itemize}

Both hypotheses above deserve a word, because the two statements they guard are among the most frequently misquoted in the subject, and a single counterexample disposes of both at once.

\begin{remark}\label{rem:gaussiantraps}
Let $X \sim \mathcal{N}(0,1)$ and let $\varepsilon$ be independent of $X$ with $\mathbb{P}(\varepsilon = 1) = \mathbb{P}(\varepsilon = -1) = \tfrac12$. Put $Y := \varepsilon X$. Conditioning on $\varepsilon$ shows at once that $Y \sim \mathcal{N}(0,1)$, and
\[
\mathbb{E}[XY] = \mathbb{E}[\varepsilon]\,\mathbb{E}[X^2] = 0,
\]
so $X$ and $Y$ are uncorrelated standard normal variables. They are very far from independent: $|X| = |Y|$ almost surely. Moreover $X + Y = (1+\varepsilon)X$ vanishes with probability $\tfrac12$, hence is not normal at all.

The moral is that ``normal'' is a property of a \emph{law on $\mathbb{R}^d$}, not a property of a list of marginals. The pair $(X,Y)$ above has normal marginals but its joint law is not multivariate normal in the sense of Example \ref{example:multivariatenormal}, and every statement in the list above fails for it. Whenever we invoke Gaussianity for several variables at once --- as we shall constantly do for the increments of Brownian motion, and for the driving noises of the multi-factor models of Section \ref{sec:othermodels} --- it is joint normality that is meant.
\end{remark}

However, the most compelling reason why the normal distribution is so important rests on the following fundamental result, in a relatively weak form:

\begin{theorem}[Central Limit Theorem]
    Let $\{X_n\}_{n=1}^\infty$ be a sequence of independent, identically distributed (hereafter, \emph{i.i.d.}) random variables  with finite mean $\mu$ and variance $\sigma^2 > 0$. Then
    \begin{equation}
    \frac{\sum_{k=1}^n X_k - n\mu}{\sigma \sqrt{n}} \xrightarrow{d} \mathcal{N}(0,1). \text{\hspace{0.6cm}} \blacksquare
    \end{equation}
\end{theorem}

It is hard to stress enough the importance and implications of the normal distribution throughout mathematics. Some insight on the multivariate case shall also be provided.


% \begin{example}[Multivariate Normal Distribution]
% Let \( X = (X_1, \dots, X_d)^\top \) be a \( \mathbb{R}^d \)-valued random vector defined on a probability space \( (\Omega, \mathcal{F}, \mathbb{P}) \). The vector \( X \) is said to follow a \textbf{multivariate normal distribution} with mean vector \( \mu \in \mathbb{R}^d \) and covariance matrix \( \Sigma \in \mathbb{R}^{d \times d} \), written as
% \[
% X \sim \mathcal{N}_d(\mu, \Sigma),
% \]
% if for every \( a \in \mathbb{R}^d \), the real-valued random variable \( a^\top X \) is normally distributed in \( \mathbb{R} \), i.e.,
% \[
% a^\top X \sim \mathcal{N}(a^\top \mu, a^\top \Sigma a).
% \]


% \end{example}

% Note that $\Sigma$ is symmetric. If \( \Sigma \) is invertible (that is, positive definite), the distribution is absolutely continuous with full support on \( \mathbb{R}^d \). If \( \Sigma \) is only positive semi-definite and not full-rank, then \( X \) is supported on a lower-dimensional affine subspace of \( \mathbb{R}^d \). In the former case, \( X \) has the following density with respect to Lebesgue measure on \( \mathbb{R}^d \):
% \[
% f_X(x) = \frac{1}{(2\pi)^{d/2} \det(\Sigma)^{1/2}} \exp\left( -\frac{1}{2}(x - \mu)^\top \Sigma^{-1} (x - \mu) \right).
% \]

% The characteristic function of a multivariate normal random vector \( X \sim \mathcal{N}_d(\mu, \Sigma) \) is $\varphi_X(u) = \mathbb{E}[e^{i u^\top X}] = \exp\left( i u^\top \mu - \frac{1}{2} u^\top \Sigma u \right), \quad \text{for all } u \in \mathbb{R}^d.$ As in the real case, these variables are determined completely by a mean vector and a covariance matrix.




% Similarly to real normal variables,  for multivariate normals uncorrelatedness implies independence. 


% \begin{example}[Log-Normal Distribution]

% A random variable \( X \) is said to follow a \textbf{log-normal distribution} with parameters \( \mu \in \mathbb{R} \) and \( \sigma > 0 \), denoted
% \[
% X \sim \text{Log}\mathcal{N}(\mu, \sigma^2),
% \]
% if the logarithm of \( X \) is normally distributed, i.e.,
% \[
% \log X \sim \mathcal{N}(\mu, \sigma^2).
% \]
% \end{example}

% The probability density function of \( X \sim \text{Log}\mathcal{N}(\mu, \sigma^2) \), after a minor change of variable calculation, is
% \[
% f_X(x) = \frac{1}{x \sigma \sqrt{2\pi}} \exp\left( -\frac{(\log x - \mu)^2}{2\sigma^2} \right), \quad x > 0.
% \]

% All positive moments of a log-normal variable exist and are explicitly given by:
% \[
% \mathbb{E}[X^k] = \exp\left( k\mu + \frac{1}{2}k^2\sigma^2 \right), \quad \text{for all } k > 0.
% \]
% In particular,
% \[
% \mathbb{E}[X] = \exp\left( \mu + \frac{\sigma^2}{2} \right), \qquad
% \mathbb{V}[X] = \left( e^{\sigma^2} - 1 \right) e^{2\mu + \sigma^2}.
% \]


% The characteristic function of a log-normal variable has no closed-form expression. Moreover, it is not integrable, and \( \varphi_X(u) \) is not defined for all complex values of \( u \). Despite having finite moments, this distribution is not uniquely determined by them. The key assumption that fails for taking advantace of Lévy inversion formula \eqref{levyInversionFormula} is the existence of the limit when $T \rightarrow \infty.$

    
% The usefulness of the log-normal distribution in finance is something that will be presented in Section  \hl{XX}.












% \begin{example}[Random Walk]
% Let \(\mathbb{T} = \mathbb{N}_0\) and consider the probability space \((\Omega, \mathcal{F}, \mathbb{P})\) where we define a sequence of i.i.d. random variables \(\{\xi_n\}_{n \in \mathbb{N}}\) such that \(\mathbb{P}(\xi_n = 1) = \mathbb{P}(\xi_n = -1) = \frac{1}{2}\). The \textbf{simple symmetric random walk} is the stochastic process \(X = \{X_n\}_{n \in \mathbb{N}_0}\) defined by:
% \[
% X_0 := 0, \quad X_n := \sum_{k=1}^n \xi_k \quad \text{for } n \geq 1.
% \]
% This process models a discrete-time evolution in which the system moves at each step by \(\pm1\) with equal probability. Each sample path is piecewise constant and consists of successive steps forming a path on \(\mathbb{Z}\). It is a classic example of a discrete-time, integer-valued stochastic process with stationary and independent increments.
% \end{example}

% It is clear that $\varphi_\xi(u) = \frac{e^{iu}+e^{-iu}}{2}=\cos(u)$ and, due to independence, $\varphi_{X_n}(u) =  \cos^n(u).$ We will return to this process, where the size of the jump and the time between jumps are put into agreement to allow a new process to take place.

% \begin{example}[Poisson Process]
% Let \(\mathbb{T} = [0, \infty)\) and fix a rate \(\lambda > 0\). On a suitable probability space \((\Omega, \mathcal{F}, \mathbb{P})\), define a stochastic process \(N = \{N_t\}_{t \geq 0}\) with values in \(\mathbb{N}_0\), where \(N_t\) counts the number of events that have occurred up to time \(t\). The process \(N\) is called a \textbf{Poisson process with rate \(\lambda\)} if it satisfies:
% \begin{itemize}
%     \item \(N_0 = 0\) almost surely,
%     \item \(N\) has independent increments,
%     \item The number of events in any interval of length \(t\) follows a Poisson distribution with mean \(\lambda t\), i.e.,
%     \[
%     \mathbb{P}(N_{t+s} - N_s = k) = \frac{(\lambda t)^k}{k!} e^{-\lambda t}, \quad \text{for all } k \in \mathbb{N}_0, \, s,t \geq 0.
%     \]
% \end{itemize}
% Sample paths of the Poisson process are piecewise constant and non-decreasing, with jumps of size \(+1\) occurring at random times. It is a fundamental model for counting processes and illustrates a continuous-time stochastic process with jump discontinuities.
% \end{example}

% After some  calculations, $\mathbb{E}[N_t] = \mathbb{V}[N_t] = \lambda t.$ We may also compute its characteristic function, arriving to $\varphi_{N_t}(u) = e^{\lambda t (e^{iu}-1)}$

% Poisson processes are useful when trying to count the amount of independent events with exponential decay that may take place during a certain amount of time. A well-known application relies on counting the amount of calls that a "centralita" receives during a normal day. It is also being used to model electricity curves, where abrupt jumps are customary. 


% Note that the Poisson process is \emph{continuous in time} because \(\mathbb{T} = [0,\infty)\), even though its sample paths exhibit discontinuities. This distinction between the continuity of the index set and the continuity of sample paths shall not lead to confusion.



% \begin{remark}
%     \emph{Despite the importance of discrete processes, \textbf{for most of the dissertation we will assume that the process is continuous}, unless otherwise stated.}
% \end{remark}

\begin{example}[Multivariate Normal Distribution]\label{example:multivariatenormal}
Let \( X = (X_1, \dots, X_d)^\top \) be a \( \mathbb{R}^d \)-valued random vector defined on a probability space \( (\Omega, \mathcal{F}, \mathbb{P}) \). The vector \( X \) is said to follow a \textbf{multivariate normal distribution} with mean vector \( \mu \in \mathbb{R}^d \) and covariance matrix \( \Sigma \in \mathbb{R}^{d \times d} \), written as
\begin{equation}
X \sim \mathcal{N}_d(\mu, \Sigma),
\end{equation}
if for every \( a \in \mathbb{R}^d \), the real-valued random variable \( a^\top X \) is normally distributed in \( \mathbb{R} \), i.e.,
\begin{equation}
a^\top X \sim \mathcal{N}(a^\top \mu, a^\top \Sigma a).
\end{equation}
\end{example}

Note that $\Sigma$ is symmetric and positive semi-definite. If \( \Sigma \) is invertible (that is, positive definite), the distribution is absolutely continuous with full support on \( \mathbb{R}^d \). If \( \Sigma \) is only positive semi-definite and not full-rank, then \( X \) is supported on a lower-dimensional affine subspace of \( \mathbb{R}^d \). In the former case, \( X \) admits the following density with respect to Lebesgue measure on \( \mathbb{R}^d \):
\begin{equation}
f_X(x) = \frac{1}{(2\pi)^{d/2} \det(\Sigma)^{1/2}} \exp\left( -\frac{1}{2}(x - \mu)^\top \Sigma^{-1} (x - \mu) \right).
\end{equation}

The characteristic function of a multivariate normal random vector \( X \sim \mathcal{N}_d(\mu, \Sigma) \) is
\begin{equation}\label{eq:mvncharfn}
\varphi_X(u) = \mathbb{E}[e^{i u^\top X}] = \exp\left( i u^\top \mu - \tfrac{1}{2} u^\top \Sigma u \right), \quad u \in \mathbb{R}^d.
\end{equation}
For a multivariate normal vector --- and, as Remark \ref{rem:gaussiantraps} showed, only then --- uncorrelatedness does imply independence. Indeed, if $\Sigma$ is block-diagonal then \eqref{eq:mvncharfn} factorises as a product of the characteristic functions of the corresponding sub-vectors, which is exactly independence.

\begin{example}[Log-Normal Distribution]
A random variable \( X \) is said to follow a \textbf{log-normal distribution} with parameters \( \mu \in \mathbb{R} \) and \( \sigma > 0 \), denoted
\begin{equation}
X \sim \text{Log}\mathcal{N}(\mu, \sigma^2),
\end{equation}
if the logarithm of \( X \) is normally distributed, i.e.,
\begin{equation}
\log X \sim \mathcal{N}(\mu, \sigma^2).
\end{equation}
\end{example}

The probability density function of \( X \sim \text{Log}\mathcal{N}(\mu, \sigma^2) \), after a straightforward change-of-variable calculation, is
\begin{equation}
f_X(x) = \frac{1}{x \sigma \sqrt{2\pi}} \exp\left( -\frac{(\log x - \mu)^2}{2\sigma^2} \right), \quad x > 0.
\end{equation}

All positive moments of a log-normal variable exist and are explicitly given by
\begin{equation}
\mathbb{E}[X^k] = \exp\left( k\mu + \tfrac{1}{2}k^2\sigma^2 \right), \quad k > 0.
\end{equation}
In particular,
\begin{equation}
\mathbb{E}[X] = \exp\left( \mu + \tfrac{\sigma^2}{2} \right), \qquad
\mathbb{V}[X] = \left( e^{\sigma^2} - 1 \right) e^{2\mu + \sigma^2}.
\end{equation}

All real moments therefore exist, and yet the log-normal distribution has two unpleasant analytical features that will matter later.

The first is that its characteristic function admits no closed-form expression. This is a genuine inconvenience --- it is precisely why the Black-Scholes model, despite its explicit pricing formula, is not amenable to the Fourier techniques that make the Heston model of Subsection \ref{subsec:stochvol} tractable --- but it is not a pathology. Note that $\varphi_X$ is defined and bounded by $1$ for every \emph{real} $u$, as it is for any random variable; what fails is the \emph{moment generating function} $\mathbb{E}[e^{uX}]$, which is infinite for every $u > 0$. The two objects are easily confused and we shall keep them apart.

The second is more surprising. Despite having moments of all orders, \textbf{the log-normal distribution is not determined by them}: there exists an explicit one-parameter family of mutually distinct densities on $(0,\infty)$ sharing every moment with it. Writing $f_X$ for the density above with $\mu = 0$, $\sigma = 1$, the family
\begin{equation}\label{eq:heyde}
f_a(x) = f_X(x)\bigl(1 + a \sin(2\pi \log x)\bigr), \qquad a \in [-1,1],
\end{equation}
consists of probability densities, and a direct computation of $\int_0^\infty x^k f_X(x)\sin(2\pi \log x)\,dx$ --- substitute $\log x = t + k$ and use that $\sin$ is odd --- shows that this integral vanishes for every $k \in \mathbb{N}$. Hence all members of \eqref{eq:heyde} have identical moments. This is Heyde's classical example \cite{heyde1963}.

It is worth being precise about what does \emph{not} fail here, since the point is easy to garble. Lévy's inversion formula \eqref{levyInversionFormula} holds for every distribution without exception, and the characteristic function does determine the log-normal law uniquely. What fails is the passage from the \emph{moment sequence} to the characteristic function: the moments grow as $e^{k^2\sigma^2/2}$, fast enough that Carleman's condition $\sum_k m_{2k}^{-1/2k} = \infty$ is violated, and with it any hope of recovering the law from its moments. The lesson is one we shall meet again in Chapter \ref{ch:cap5}: knowing every moment of an object is strictly weaker than knowing the object, and the signature transform will be interesting precisely because, for paths, the analogous statement can be made to go the right way.



\begin{example}[Random Walk]\label{example:randomwalk}
Let \(\mathbb{T} = \mathbb{N}_0\) and consider the probability space \((\Omega, \mathcal{F}, \mathbb{P})\) where we define a sequence of  independent, identically distributed random variables \(\{\xi_n\}_{n \in \mathbb{N}}\) such that \(\mathbb{P}(\xi_n = 1) = \mathbb{P}(\xi_n = -1) = \tfrac{1}{2}\). The \textbf{simple symmetric random walk} is the stochastic process \(X = \{X_n\}_{n \in \mathbb{N}_0}\) defined by
\begin{equation}
X_0 := 0, \quad X_n := \sum_{k=1}^n \xi_k, \quad n \geq 1.
\end{equation}
This process models a discrete-time evolution in which the system moves at each step by \(\pm1\) with equal probability. Each sample path is piecewise constant and consists of successive steps forming a path on \(\mathbb{Z}\). It is a classic example of a discrete-time, integer-valued stochastic process with stationary and independent increments.
\end{example}

It is clear that
\begin{equation}
\varphi_\xi(u) = \frac{e^{iu}+e^{-iu}}{2}=\cos(u),
\end{equation}
and, due to independence,
\begin{equation}
\varphi_{X_n}(u) =  \cos^n(u).
\end{equation}
We shall return to this process in Theorem \ref{thm:donsker}, where the size of the jump and the time between jumps are adjusted simultaneously, and a new process emerges in the limit.

\begin{example}[Poisson Process]\label{example:poissonprocess}
Let \(\mathbb{T} = [0, \infty)\) and fix a rate \(\lambda > 0\). On a suitable probability space \((\Omega, \mathcal{F}, \mathbb{P})\), define a stochastic process \(N = \{N_t\}_{t \geq 0}\) with values in \(\mathbb{N}_0\), where \(N_t\) counts the number of events that have occurred up to time \(t\). The process \(N\) is called a \textbf{Poisson process with rate \(\lambda\)} if it satisfies:
\begin{itemize}
    \item \(N_0 = 0\) almost surely,
    \item \(N\) has independent increments,
    \item The number of events in any interval of length \(t\) follows a Poisson distribution with mean \(\lambda t\), i.e.,
    \begin{equation}
    \mathbb{P}(N_{t+s} - N_s = k) = \frac{(\lambda t)^k}{k!} e^{-\lambda t}, \quad k \in \mathbb{N}_0, \, s,t \geq 0.
    \end{equation}
\end{itemize}
Sample paths of the Poisson process are piecewise constant and non-decreasing, with jumps occurring at random times. It is a fundamental model for counting processes and illustrates a continuous-time stochastic process with jump discontinuities.
\end{example}

After some calculations,
\begin{equation}
\mathbb{E}[N_t] = \mathbb{V}[N_t] = \lambda t.
\end{equation}
We may also compute its characteristic function, leading to
\begin{equation}
\varphi_{N_t}(u) = e^{\lambda t (e^{iu}-1)}.
\end{equation}

Poisson processes are useful when trying to count the number of independent events occurring randomly over time, with exponentially distributed inter-arrival times. A well-known application is counting the number of calls that a switchboard receives during a typical day. Poisson processes are also used to model electricity demand curves, where abrupt jumps are customary.

Note that the Poisson process is \emph{continuous in time} because \(\mathbb{T} = [0,\infty)\), even though its sample paths exhibit discontinuities. This distinction between the continuity of the index set and the continuity of sample paths should not cause confusion.

\begin{remark}
    \emph{Despite the importance of discrete processes, \textbf{for most of this dissertation we will assume that the process is continuous}, unless otherwise stated.}
\end{remark}














\subsection{Markov Processes}

Markov processes are characterized by a specific form of \emph{memorylessness}. They are particularly suited to describe the evolution of systems where the future depends only on the present state and not on the past history. Some of these processes arise as \textit{diffusions} and in that case, a lot can be said about the distributions.




\begin{definition}[Markov Process]
Let \((\Omega, \mathcal{F}, \mathbb{P})\) be a probability space, and let \(X = \{X_t\}_{t \in \mathbb{T}}\) be a stochastic process taking values in a measurable space \((E, \mathcal{E})\). We say that \(X\) is a \textbf{Markov process} if for every $s < t$ and every measurable set $A \in \mathcal{E}$,
\end{definition}

\begin{equation}
\mathbb{P}(X_t \in A \mid \mathcal{F}_s) = \mathbb{P}(X_t \in A \mid X_s) \quad \text{almost surely}, 
\end{equation}

where $\mathcal{F}_s = \sigma(X_r : r \leq s)$ is the $\sigma$-algebra generated by the history of the process up to time $s$. Let us define the \textbf{transition probability} of  Markov process by $  \mathbb{P}(X_{t} \in A \mid X_0 = x)$. In order to integrate against the measure given by the transition probability, we can write $P_{s,x}(t,dy) = \mathbb{P}(X_t \in dy | X_s = x).$


There is a rather sharp distinction between finite-state Markov processes and those with infinitely many states. In the finite case, the theory is built around transition matrices: the probability of moving from one state to another in a single step is recorded in a matrix \(P_t\), so that the evolution can either be understood as $\prod_{i=1}^n P_i$. Long-run behaviour can then often be analysed by classical linear-algebraic tools, specially in the stationary case, that is, when the transition probability is independent of $t$.

In the infinite-state setting,  the dynamics are described by operators acting on spaces of functions or measures. The study of these processes therefore relies on functional analysis, differential equations and probabilistic techniques, showing that these topics are deeply linked. The evolution can be described by a family of linear operators $(P_t)_{t \geq 0}$ acting on a suitable class of functions $f : S \to \mathbb{R}$, following:


\begin{equation}
P_t f(x) := \mathbb{E}^x[f(X_t)] = \int_S f(y) P_t(x, dy), 
\end{equation}

where $\mathbb{E}^x[\cdot]$ denotes expectation conditioned on $X_0 = x$. The family $(P_t)_{t \geq 0}$ forms a \textbf{semigroup of operators} in the sense that: $P_0 = I$ and $P_{t+s} = P_t \circ P_s \quad \forall t, s \geq 0 $. Moreover, under suitable regularity conditions, we have:

\[
\lim_{t \to 0^+} P_t f = f.
\]


The \textbf{infinitesimal generator} $A$ of the semigroup $(P_t)$ is the linear operator defined by:

\begin{equation}
Af := \lim_{t \to 0^+} \frac{P_t f - f}{t}, 
\end{equation}

for those functions $f$ for which the limit exists in an appropriate sense (e.g., uniform limit if $f$ is bounded and continuous). The domain of the generator is denoted by

\[
\mathcal{D}(A) := \left\{ f \in \mathcal{B}(S) : \lim_{t \to 0^+} \frac{P_t f - f}{t} \text{ exists} \right\}.
\]

The generator captures the infinitesimal behavior of the process. This generator is of seminal importance because it allows to connect Markov processes with PDEs and functional analysis, as will be shown in Subsection \ref{section:itoDiff}.














\subsection{Martingales, semimartingales and local martingales}\label{subsec:martingales}

Martingales are processes for which the expected value of a future instant is precisely the state of the process up to present time. Martingales are absolutely key in the study of stochastic processes. Some generalizations are needed as well. We keep the notation used in the previous section.


\begin{definition}{(Martingale)} 
    Let $X_t$ be a stochastic process, adapted to a given filtration $\mathcal{F}_t$. We say it is a \textbf{martingale} provided:
    \begin{itemize}
        \item $\mathbb{E}[|X_t|]<\infty$ for almost all $t.$
        \item $X_s = \mathbb{E}[X_t | \mathcal{F}_s]$ for $s<t.$
    \end{itemize}
    We say that $X_t$ is a \textbf{supermartinagle} if the second condition is replaced by $X_s  \geq \mathbb{E}[X_t | \mathcal{F}_s]$ and a \textbf{submartingale} if $X_s\leq \mathbb{E}[X_t | \mathcal{F}_s].$
\end{definition}

The martingale property is stated at deterministic times, but a great deal of what we shall do --- 
localisation in Section \ref{sec:stochasticcalculus}, and the early exercise of American options in 
Subsection \ref{subsec:american} --- involves evaluating processes at \emph{random} times. For this 
to make sense we must first say what information is available at a stopping time.

\begin{definition}[The $\sigma$-algebra $\mathcal{F}_\tau$]\label{def:Ftau}
Let $\tau$ be a stopping time with respect to $\Filt$. The $\sigma$-algebra of events observable up 
to $\tau$ is
\begin{equation}
\mathcal{F}_\tau := \bigl\{ A \in \mathcal{F} \;:\; A \cap \{\tau \leq t\} \in \mathcal{F}_t \ \text{ for every } t \in [0,T] \bigr\}.
\end{equation}
\end{definition}

The definition reads oddly at first but says something simple: an event belongs to $\mathcal{F}_\tau$ 
if, whenever $\tau$ has occurred by time $t$, an observer at time $t$ can already tell whether the 
event happened. One checks readily that $\mathcal{F}_\tau$ is indeed a $\sigma$-algebra, that 
$\tau$ is $\mathcal{F}_\tau$-measurable, that $\mathcal{F}_\tau = \mathcal{F}_t$ when 
$\tau \equiv t$ is constant, and that $\sigma \leq \tau$ implies 
$\mathcal{F}_\sigma \subseteq \mathcal{F}_\tau$. If $X$ is right-continuous and adapted, then 
$X_\tau$ is $\mathcal{F}_\tau$-measurable on $\{\tau < \infty\}$.

\begin{theorem}[Doob's optional stopping theorem]\label{thm:optionalstopping}
Let $X$ be a right-continuous martingale with respect to $\Filt$, and let $\sigma \leq \tau$ be two 
stopping times. If any one of the following holds --- 
\begin{itemize}
    \item[i)] $\tau$ is bounded, i.e.\ $\tau \leq T$ almost surely for the finite horizon $T$;
    \item[ii)] $X$ is uniformly integrable; or
    \item[iii)] there exists $Y \in L^1(\Omega)$ with $|X_{t}| \leq Y$ for all $t$,
\end{itemize}
--- then $X_\sigma$ and $X_\tau$ are integrable and
\begin{equation}\label{eq:optionalstopping}
\mathbb{E}\bigl[ X_\tau \mid \mathcal{F}_\sigma \bigr] = X_\sigma \qquad \text{almost surely.}
\end{equation}
In particular $\mathbb{E}[X_\tau] = \mathbb{E}[X_0]$. The corresponding statements with $=$ replaced 
by $\leq$ or $\geq$ hold for sub- and supermartingales respectively. \hfill $\blacksquare$
\end{theorem}

Equation \eqref{eq:optionalstopping} is the assertion that a martingale remains a martingale when 
observed along a random clock, and it is one of those results whose hypotheses are more instructive 
than its conclusion. Some hypothesis is indispensable: let $W$ be a Brownian motion and 
$\tau = \inf\{t : W_t = 1\}$, which is finite almost surely. Then $W_\tau = 1$ identically, so 
$\mathbb{E}[W_\tau] = 1 \neq 0 = \mathbb{E}[W_0]$, and \eqref{eq:optionalstopping} fails. Nothing is 
wrong with the theorem; the stopping time is simply unbounded and $W_{t\wedge\tau}$ is not uniformly 
integrable. The financial reading of this example is exact and worth stating: a strategy that waits 
however long is necessary to be up by one unit does indeed win with certainty, and it is not an 
arbitrage only because executing it may require unbounded losses along the way. That is precisely 
what the admissibility constraint of Definition \ref{def:admissible} will forbid.

Because $\{X_{t\wedge\tau}\}$ is a martingale whenever $X$ is --- which is case i) applied at each 
$t$ --- optional stopping is also the mechanism behind localisation. A local martingale is by 
Definition \ref{def:localmartingale} a process which becomes a martingale after stopping; the content 
of that definition is that all the martingale machinery, including \eqref{eq:optionalstopping} and 
Proposition \ref{prop:doobmaximal}, remains available on each piece.

The following result will be useful when constructing stochastic processes from Itô integrals. Notice that if $X_t$ is a martingale, then by Jensen's inequality (Proposition \ref{def:conditional-expectation}, item 8) $X_t^2$ is a submartingale.
 
\begin{proposition}[Doob's submartingale inequality]\label{prop:doobmaximal}
    Let $Z = \{Z_t\}_{t\in[0,T]}$ be a right-continuous submartingale. Then, for any $\varepsilon > 0$,
    $$
    \mathbb{P}\Bigl(\sup_{0\leq t\leq T} Z_t \geq \varepsilon\Bigr) \leq \frac{\mathbb{E}[\max\{Z_T,0\}]}{\varepsilon}.
    $$
    In particular, for a right-continuous martingale $X$, $\mathbb{P}(\sup_{0\leq t\leq T} |X_t| \geq \varepsilon) \leq \mathbb{E}[|X_T|]/\varepsilon.$
\end{proposition}
\textbf{\textit{Proof:}} (...) \hfill $\blacksquare$

The bound is in terms of the \emph{terminal} value $X_T$, not of a running $X_t$; this is the whole content of the inequality, since it converts information at a single time into control of the entire path. For the $L^2$ theory of the Itô integral we shall need the following stronger, $L^p$ form, which is proved by integrating the tail estimate above.

\begin{proposition}[Doob's $L^p$ maximal inequality]\label{prop:dooblp}
Let $X$ be a right-continuous martingale and let $p > 1$. Then
\[
\mathbb{E}\Bigl[\sup_{0 \leq t \leq T}|X_t|^p\Bigr] \leq \left(\frac{p}{p-1}\right)^{p} \mathbb{E}\bigl[|X_T|^p\bigr].
\]
In particular, for $p=2$, $\mathbb{E}\bigl[\sup_{t\leq T}|X_t|^2\bigr] \leq 4\,\mathbb{E}[|X_T|^2]$. \hfill $\blacksquare$
\end{proposition}

The case $p=2$ is the one we shall use repeatedly. It says that for a square-integrable martingale, controlling the terminal second moment controls the second moment of the whole trajectory --- which is exactly what is needed to upgrade convergence of the Itô integral in $L^2(\Omega)$, at each fixed time, to uniform convergence along a subsequence, and hence to continuity of sample paths (Theorem \ref{thm:itocontinuity}).





\noindent
The notion of a martingale is often too restrictive, particularly when considering processes such as stochastic exponentials or certain diffusion models. A relaxation is therefore introduced.

\begin{definition}[Local Martingale]\label{def:localmartingale}
A process $X = \{X_t\}_{t \geq 0}$ is a \textbf{local martingale} if there exists an increasing sequence of stopping times $(\tau_n)_{n \geq 1}$ with $\tau_n \uparrow \infty$ almost surely, such that each stopped process $X^{\tau_n}_t := X_{t \wedge \tau_n}$ is a martingale.
\end{definition}

The idea is that $X$ behaves like a martingale, but only up to random horizons. An analogy could drawn from ODEs: there are solutions that explode to infinity in finite time, and those that can be extented to all times. Here, the integrability condition can be somewhat replaced if we work with sequences of stopping times. This mathematical object deserves its own study but lies beyond the scope of this dissertation.

\noindent
Finally, there is another related concept, broad enough to serve as a class of \textit{good integrators}. In particular, it extends the role played by Brownian motion as a canonical integrator, while retaining enough structure to allow a general theory of stochastic calculus.


\begin{definition}[Semimartingale]\label{def:semimartingale}
A process $X$ is called a \textbf{semimartingale} if it can be decomposed as
\[
X_t = M_t + A_t,
\]
where $M_t$ is a local martingale and $A_t$ is  \textit{continue à droite, avec limite à gauche} (càdlàg) and of locally bounded variation.
\end{definition}

This last definition is the correct level of generality for stochastic integration, and it is 
worth understanding why the decomposition takes the shape it does. The two summands are exactly the 
two mechanisms by which an integrator can be usable: $A$ has bounded variation, so 
$\int g\,dA$ exists pathwise as a Riemann--Stieltjes integral in the classical sense; and $M$ is a 
local martingale, so $\int g\,dM$ exists in the Itô sense, by the $L^2$ construction of Section 
\ref{sec:stochasticcalculus} together with localisation. A semimartingale is precisely a process 
that can be integrated against by splitting the work between deterministic analysis and 
probability, and the Bichteler--Dellacherie theorem makes this converse precise: semimartingales 
are, in a suitable topological sense, the \emph{only} processes admitting a reasonable theory of 
stochastic integration; see \cite{protter2005}.

We shall not develop that general theory --- every integrator in Part I is continuous, and for 
continuous processes the decomposition is unique --- but the definition is recorded because of what 
happens when it fails. Not every process of interest is a semimartingale: fractional Brownian 
motion with Hurst parameter $H \neq \tfrac12$ is not one, a result of Rogers \cite{rogers1997}, and consequently 
neither the Itô integral nor any of the machinery built on it is available for it. Since such 
processes are exactly what the empirical evidence on volatility will point towards, the reader 
should register Definition \ref{def:semimartingale} less as a technical generalisation than as a 
boundary: it marks the outer limit of the theory developed in this chapter, and Part 
\ref{part:part2} begins where it ends.








\subsection{Brownian Motion}\label{subsec:brownianmotion}

Brownian motion is a fundamental stochastic process whose importance has grown 
throughout many areas of mathematics. It is complex enough to generate a new class 
of \textit{stochastic differential equations} and to address numerous applied problems 
involving randomness, yet simple enough to allow explicit calculations. 
Its properties also enable efficient simulation, making it a cornerstone of Monte Carlo 
methods across disciplines.

\begin{definition}[Brownian Motion] \label{def:Brownian-motion}
A stochastic process \(W : \mathbb{R}^+ \times \Omega \to \mathbb{R}\) is called a 
\textbf{Brownian motion} (or \textbf{Wiener process}) if the following conditions hold:
\begin{itemize}
    \item[i)] \begin{equation}
        \mathbb{P}(W_0 = 0) = 1.
    \end{equation}
    \item[ii)] \begin{equation}
        \mathbb{P}\big( \{ \omega : t \mapsto W_t(\omega) 
        \text{ is continuous} \} \big) = 1.
    \end{equation}
    \item[iii)] For \(0 \leq s < t\), 
    \begin{equation}
        W_t - W_s \sim \mathcal{N}(0, t-s).
    \end{equation}
    \item[iv)] The process has independent increments: for 
    \(0 \leq t_1 < t_2 < \dots < t_n\), the random variables
    \begin{equation}
        \{ W_{t_1}, \ W_{t_2}-W_{t_1}, \ W_{t_3}-W_{t_2}, \ \dots, \ W_{t_n}-W_{t_{n-1}} \}
    \end{equation}
    are independent.
\end{itemize}
\end{definition}

Some properties of Brownian motion follow immediately from the definition 
and are presented below. The last one illustrates the property known as 
\emph{stochastic self-similarity}. 

\medskip

In order to work with the broadest possible class of stochastic processes, 
we introduce the following general result.

\begin{proposition} \label{prop:independent-increments}
Let \(S(t), \, t \geq 0\), be a process adapted to the filtration 
\((\mathcal{F}_t)_{t \geq 0}\) such that for any \(0 \leq s < t\),
the increment \(S(t)-S(s)\) is independent of \(\mathcal{F}_s\). 
Then \(S(t)\) has independent increments.
\end{proposition}

\textit{\textbf{Proof}.} (...) \hfill $\blacksquare$

\medskip

Therefore, a process \(B(t)\) satisfying the conditions of 
Definition~\ref{def:Brownian-motion}, together with a filtration satisfying 
Proposition~\ref{prop:independent-increments}, is a Brownian motion.\footnote{
Analogous considerations hold for a Poisson process; see \cite{durrett2019probability}.} 
Throughout, we will work with processes adapted to their natural filtration, i.e. 
\(B_t\) is \(\mathcal{F}_t\)-measurable and increments \(B(t)-B(s)\) are independent 
of \(\mathcal{F}_s\) whenever \(s \leq t\).

The next definition and proposition highlight some of the structural richness 
of such processes. In particular, the last property quantifies the 
\emph{roughness} of sample paths, a key feature of Brownian motion.

\begin{definition}[Quadratic Variation] \label{def:quadratic-variation}
Let \(f : [0,T] \to \mathbb{R}\), and for a partition \(\mathcal{P} = \{0=t_0 < t_1 < \dots < t_n=t\}\) of \([0,t]\) write
\begin{equation}
\lVert \mathcal{P} \rVert = \max_{1 \leq k \leq n} (t_k - t_{k-1}).
\end{equation}
If the limit
\begin{equation}\label{eq:qvdef}
[f]_t := \lim_{\lVert \mathcal{P} \rVert \to 0} 
\sum_{k=1}^n \big( f(t_k) - f(t_{k-1}) \big)^2
\end{equation}
exists, it is called the \textbf{quadratic variation} of \(f\) on \([0,t]\).
\end{definition}

Two cautions are in order. First, the limit in \eqref{eq:qvdef} need not exist for an arbitrary $f$, which is why the definition is phrased conditionally. Second --- and this is the point that will matter --- for a \emph{stochastic} process the sums in \eqref{eq:qvdef} are random variables, and the limit must be specified as a limit of random variables. We shall always understand it in $L^2(\Omega)$, which for our purposes is both the easiest to obtain and strong enough to imply convergence in probability, and hence almost sure convergence along a subsequence.

It is worth pausing on how strange \eqref{eq:qvdef} is from the point of view of ordinary calculus. If $f$ is continuously differentiable, then $|f(t_k)-f(t_{k-1})| \leq C\,(t_k - t_{k-1})$ and the sum is bounded by $C^2 \lVert\mathcal{P}\rVert \, t \to 0$: smooth functions have zero quadratic variation. The same computation applies to any function of bounded variation that is continuous. A non-zero quadratic variation is therefore a certificate of irregularity, and the following proposition asserts that Brownian motion carries exactly such a certificate.

\begin{proposition} \label{prop:bm-properties}
Let \((W_t)_{t \geq 0}\) be a Brownian motion. Then:
\begin{itemize}
    \item[a)] For \(t > 0\), \(W_t \sim \mathcal{N}(0,t)\). Moreover, for all 
    \(s,t \geq 0\),
    \begin{equation}
        \mathbb{E}[W_s W_t] = \min\{s,t\}.
    \end{equation}
    \item[b)] For \(t_0 \geq 0\), the process
    \begin{equation}
        Y(t) = W(t+t_0) - W(t_0)
    \end{equation}
    is also a Brownian motion.
    \item[c)] For any \(\lambda > 0\), the process
    \begin{equation}
        Z(t) = \frac{1}{\sqrt{\lambda}} W(\lambda t)
    \end{equation}
    is also a Brownian motion.
    \item[d)] Brownian motion is both a Markov process and a martingale. 
    \item[e)] The quadratic variation of \(W\) over \([0,t]\) exists and equals
    \begin{equation}\label{eq:bmqv}
        [W]_t = t \qquad \text{the limit being taken in } L^2(\Omega).
    \end{equation}
    \item[f)] The process $W_t^2 - t$ is a martingale.
\end{itemize}
Properties b) and c) are called \emph{translation invariance} 
and \emph{scaling invariance}, respectively.
\end{proposition}

\textit{\textbf{Proof}.} (...) \hfill $\blacksquare$

\medskip

One might wonder whether such a process actually exists, given the strong 
set of requirements in Definition~\ref{def:Brownian-motion}. The question is not idle: 
we have demanded independent Gaussian increments \emph{and} continuity of every sample path, 
and there is no reason a priori why the two should be compatible. A full construction would 
take us too far afield, but it would be unsatisfying to leave the matter entirely unaddressed, 
so we indicate the two results which settle it --- both of which we shall have occasion to use 
again for their own sake.

The first supplies the process. Recall the simple symmetric random walk of Example 
\ref{example:randomwalk}, where we promised to return once the jump size and the waiting time 
between jumps had been adjusted. This is that return.

\begin{theorem}[Donsker's invariance principle]\label{thm:donsker}
Let $\{\xi_k\}_{k\geq1}$ be i.i.d.\ with $\mathbb{E}[\xi_1]=0$ and $\mathbb{V}[\xi_1]=1$, and let 
$X_n = \sum_{k=1}^n \xi_k$ be the associated random walk. Define the rescaled, linearly 
interpolated process
\begin{equation}\label{eq:donskerscaling}
W^{(n)}_t := \frac{1}{\sqrt{n}}\Bigl( X_{\lfloor nt \rfloor} + (nt - \lfloor nt \rfloor)\,\xi_{\lfloor nt\rfloor +1}\Bigr), \qquad t \in [0,T].
\end{equation}
Then $W^{(n)} \Rightarrow W$ as $n \to \infty$, where the convergence is in distribution on the 
space $C[0,T]$ of continuous functions with the uniform topology, and $W$ is a Brownian motion. 
\hfill $\blacksquare$
\end{theorem}

The scaling in \eqref{eq:donskerscaling} is the whole content of the theorem: space is compressed 
by $\sqrt{n}$ while time is compressed by $n$. That is exactly the relation $\text{space} \sim 
\sqrt{\text{time}}$ which Definition \ref{def:Brownian-motion} imposes through 
$W_t - W_s \sim \mathcal{N}(0,t-s)$, and which reappears as the scaling invariance of Proposition 
\ref{prop:bm-properties}(c). Donsker's theorem also explains why Brownian motion is 
\emph{universal} in the same sense that the Gaussian distribution is: the limit does not depend on 
the law of $\xi_1$ beyond its first two moments, so any accumulation of small, independent, 
finite-variance shocks produces the same object. This is the honest reason Brownian motion is 
ubiquitous in modelling, and it is worth keeping in mind that the same universality argument gives 
no warrant at all for the shocks being independent --- a point the empirical evidence of Part 
\ref{part:part2} will press hard.

The second result supplies the continuity, and does so from a moment condition alone.

\begin{theorem}[Kolmogorov's continuity criterion]\label{thm:kolmogorovcontinuity}
Let $X = \{X_t\}_{t\in[0,T]}$ be a stochastic process for which there exist constants 
$\alpha, \beta, C > 0$ with
\begin{equation}\label{eq:kolmogorovmoment}
\mathbb{E}\bigl[\,|X_t - X_s|^{\alpha}\,\bigr] \leq C\,|t-s|^{1+\beta}, \qquad s,t \in [0,T].
\end{equation}
Then $X$ admits a continuous modification $\tilde{X}$. Moreover, the paths of $\tilde{X}$ are 
almost surely locally $\gamma$-Hölder continuous for every $\gamma \in (0,\beta/\alpha)$. 
\hfill $\blacksquare$
\end{theorem}

Applied to Brownian motion the criterion is immediate. Since $W_t - W_s \sim \mathcal{N}(0,t-s)$, 
the Gaussian moments give $\mathbb{E}[|W_t-W_s|^{2m}] = c_m |t-s|^{m}$ for every $m \in \mathbb{N}$, 
so \eqref{eq:kolmogorovmoment} holds with $\alpha = 2m$ and $\beta = m-1$. A continuous 
modification therefore exists, which is what was needed; and the H\"older conclusion yields 
exponents $\gamma < (m-1)/2m$, which as $m \to \infty$ approach $1/2$ from below but never reach 
it. That the supremum is exactly $1/2$, and that it is not attained, is the subject of the next 
subsection.

Note the word \emph{modification}: what Theorem \ref{thm:kolmogorovcontinuity} produces is a 
process agreeing with $X$ at each fixed time almost surely, not one agreeing with $X$ along whole 
paths (Remark \ref{rem:modification}). This is not a defect of the theorem but a necessity, since 
\eqref{eq:kolmogorovmoment} constrains only two-dimensional distributions and cannot possibly 
determine path behaviour on its own. As is customary, we shall always work with the continuous 
modification and cease to mention the distinction.

\subsection{Path regularity}\label{subsec:pathregularity}

We now examine how irregular the trajectories of Brownian motion actually are. This subsection is 
short, and every statement in it is classical, but it is the most consequential material in the 
chapter: the construction of the Itô integral in Section \ref{sec:stochasticcalculus} is a 
\emph{response} to what follows, and the theory developed in Part \ref{part:part2} is a response 
to what happens when these estimates are pushed slightly further.

We begin with the observation that quadratic variation and total variation are in direct 
competition.

\begin{definition}[Total variation]\label{def:totalvariation}
The \textbf{total variation} of $f : [0,T]\to\mathbb{R}$ on $[0,t]$ is
\begin{equation}
V_1(f)_t := \sup_{\mathcal{P}} \sum_{k=1}^n \bigl| f(t_k) - f(t_{k-1}) \bigr|,
\end{equation}
the supremum being over all partitions $\mathcal{P} = \{0=t_0<\cdots<t_n=t\}$. If 
$V_1(f)_t < \infty$ we say $f$ has \textbf{bounded variation} on $[0,t]$.
\end{definition}

Bounded variation is exactly the condition under which the Riemann--Stieltjes integral 
$\int_0^t g\,df$ is defined for every continuous $g$; equivalently, it is the condition under 
which $f$ induces a signed measure $df$. It is therefore the natural first thing to ask of a 
prospective integrator.

\begin{proposition}\label{prop:bmunboundedvariation}
Let $W$ be a Brownian motion. Then, almost surely, $V_1(W)_t = \infty$ for every $t > 0$.
\end{proposition}

\textbf{\textit{Proof:}} Fix $t>0$ and a sequence of partitions $\mathcal{P}_n$ with 
$\lVert\mathcal{P}_n\rVert \to 0$. On any partition,
\begin{equation}\label{eq:qvtvbound}
\sum_{k} \bigl(W_{t_k} - W_{t_{k-1}}\bigr)^2 
\;\leq\; \Bigl(\max_{k} \bigl|W_{t_k} - W_{t_{k-1}}\bigr|\Bigr) \cdot \sum_{k}\bigl|W_{t_k}-W_{t_{k-1}}\bigr|
\;\leq\; \Bigl(\max_{k}\bigl|W_{t_k}-W_{t_{k-1}}\bigr|\Bigr)\, V_1(W)_t .
\end{equation}
Suppose, for contradiction, that $V_1(W)_t < \infty$ on an event of positive probability. Sample 
paths of $W$ are continuous, hence uniformly continuous on the compact $[0,t]$, so the maximal 
increment on the right of \eqref{eq:qvtvbound} tends to $0$ as $\lVert\mathcal{P}_n\rVert \to 0$. 
The whole right-hand side therefore tends to $0$, forcing $[W]_t = 0$ on that event. This 
contradicts \eqref{eq:bmqv}, which gives $[W]_t = t > 0$ in $L^2(\Omega)$ and hence almost surely 
along a subsequence. \hfill $\blacksquare$

This is the single most important fact in the chapter, and it is worth stating its consequence 
without any hedging. \textbf{The integral $\int_0^t f\,dW$ cannot be defined pathwise as a 
Riemann--Stieltjes integral}, for any non-trivial class of integrands, because $W$ does not induce 
a measure on $[0,t]$ along any path. Every difficulty in Section \ref{sec:stochasticcalculus} --- 
the restriction to adapted integrands, the insistence on evaluating at the left endpoint, the 
approximation in $L^2(\Omega)$ rather than pathwise, the appearance of a second-order term in the 
chain rule --- traces back to \eqref{eq:qvtvbound}. Itô's construction is not one possible approach 
among several; it is a way of exploiting probabilistic structure to recover an object which 
deterministic analysis cannot reach.

The same competition can be expressed more finely by asking, not whether the sum of increments 
converges, but for which exponent $p$ the sum of $p$-th powers does.

\begin{definition}[$p$-variation]\label{def:pvariation}
Let $p \geq 1$. The \textbf{$p$-variation} of $f : [0,T] \to \mathbb{R}$ on $[0,t]$ is
\begin{equation}
V_p(f)_t := \left( \sup_{\mathcal{P}} \sum_{k=1}^n \bigl|f(t_k)-f(t_{k-1})\bigr|^p \right)^{1/p},
\end{equation}
and $f$ is said to have \emph{finite $p$-variation} on $[0,t]$ if $V_p(f)_t < \infty$. We write 
$f \in \mathcal{C}^{p\text{-var}}([0,t])$.
\end{definition}

For $p = 1$ this is Definition \ref{def:totalvariation}, and for $p=2$ it is a supremum version of 
the quadratic variation of Definition \ref{def:quadratic-variation}. The scale is monotone: if 
$f$ is continuous and $V_p(f)_t < \infty$, then $V_q(f)_t < \infty$ for every $q > p$, since 
$|f(t_k)-f(t_{k-1})|^q \leq (\max_k |f(t_k)-f(t_{k-1})|)^{q-p} |f(t_k)-f(t_{k-1})|^p$. Each 
function thus has a critical exponent below which its $p$-variation is infinite and above which it 
is finite, and this exponent is a precise measure of roughness.

\begin{theorem}[Regularity of Brownian paths]\label{thm:bmregularity}
Let $W$ be a Brownian motion. Then, almost surely:
\begin{itemize}
    \item[i)] the sample path $t \mapsto W_t$ is locally $\gamma$-Hölder continuous for every 
    $\gamma < \tfrac12$, and is locally $\gamma$-Hölder for no $\gamma \geq \tfrac12$;
    \item[ii)] $V_p(W)_t < \infty$ for every $p > 2$ and every $t$, whereas $V_p(W)_t = \infty$ for 
    every $p \leq 2$;
    \item[iii)] the sample path is nowhere differentiable: for almost every $\omega$, there is no 
    $t$ at which $t \mapsto W_t(\omega)$ possesses a finite derivative. \hfill $\blacksquare$
\end{itemize}
\end{theorem}

Part i) was half-proved above via Kolmogorov's criterion; the negative half, like iii), is due to 
Paley, Wiener and Zygmund, and proofs of all three may be found in \cite{falconer} or 
\cite{karatzasshreve1991}. Part iii) is worth a moment's reflection. A continuous function 
which is differentiable nowhere is already a celebrated pathology in real analysis, exhibited with 
some effort by Weierstrass; the theorem says that if one picks a Brownian path at random one 
obtains such a function with probability one. The pathological case is the generic case.

It is part ii) which will matter most. Let us record it in the form we shall use.

\begin{remark}[The critical exponent]\label{rem:criticalexponent}
Brownian motion has finite $p$-variation exactly for $p > 2$. The number $2$ is a threshold in a 
strong sense, and the reason is the following. For a path of finite $p$-variation with $p < 2$, 
the integral $\int g\,df$ can be constructed \emph{pathwise}, by purely deterministic means, 
provided $g$ has finite $q$-variation with $1/p + 1/q > 1$; this is Young's integral, and no 
probability is involved. Brownian motion misses this regime --- but only just, and by an amount 
that no refinement of the estimate can recover, since $V_2(W)_t = \infty$.

The Itô integral of Section \ref{sec:stochasticcalculus} is precisely the tool that bridges this 
gap, and it does so by spending probabilistic structure: the integrand must be adapted, the 
approximation is in $L^2(\Omega)$ rather than pathwise, and the resulting object is defined only 
up to null sets. That is a real price. It buys us the entire theory of Part I, but it also means 
the theory is tied to a specific process and a specific filtration, and it is exactly this 
dependence which becomes an obstacle when the driving signal is rougher than Brownian --- when 
even $p = 2$ is out of reach, and neither Young nor Itô applies. The theory that operates in that 
regime is the subject of Part \ref{part:part2}.
\end{remark}

\figurehere{Simulated Brownian path with, alongside it, the partition sums for total variation 
(diverging) and quadratic variation (converging to $t$) as the mesh is refined. This is the visual 
statement of Proposition \ref{prop:bmunboundedvariation}.}

Finally, we record the relation between the pathwise quadratic variation of Definition 
\ref{def:quadratic-variation} and the predictable version produced by the Doob--Meyer 
decomposition of Subsection \ref{subsec:measurechanges}, since both notations occur in the 
literature and we shall use them interchangeably.

\begin{proposition}\label{prop:qvcoincide}
Let $M$ be a continuous square-integrable martingale. Then the pathwise quadratic variation 
$[M]_t$ of Definition \ref{def:quadratic-variation} exists as a limit in $L^2(\Omega)$ and 
coincides, up to indistinguishability, with the compensator $\langle M \rangle_t$ of Theorem 
\ref{thm:doobmeyer}. Furthermore, the polarised quantity
\begin{equation}\label{eq:polarisation}
\langle M, N \rangle_t := \tfrac{1}{2}\bigl( \langle M+N \rangle_t - \langle M \rangle_t - \langle N \rangle_t \bigr)
= \lim_{\lVert\mathcal{P}\rVert\to0} \sum_k \bigl(M_{t_k}-M_{t_{k-1}}\bigr)\bigl(N_{t_k}-N_{t_{k-1}}\bigr)
\end{equation}
defines the \textbf{quadratic covariation} of two continuous square-integrable martingales, is 
bilinear and symmetric, and satisfies the Kunita--Watanabe inequality
\begin{equation}
\int_0^t \bigl| d\langle M,N\rangle_s \bigr| \;\leq\; \sqrt{\langle M\rangle_t}\,\sqrt{\langle N\rangle_t}
\qquad \text{almost surely.} \hfill \blacksquare
\end{equation}
\end{proposition}

The distinction between $[M]$ and $\langle M\rangle$ is genuine for processes with jumps, where the 
two differ by the accumulated squared jumps; for the continuous processes of this dissertation it 
is empty, and we shall write $\langle M \rangle$ throughout. Quadratic covariation is what gives 
meaning to the off-diagonal entries $\rho_{ij}\,dt$ of the multiplication table in Subsection 
\ref{subsec:multidimensional}, and it is the object that the correlation parameter of a stochastic 
volatility model actually controls.







% Stochastic Calculus is a wide area in Mathematics that allows, for example, to incorporate some kind of randomness and uncertainty in differential equations. One of our main aims is to give meaning to expressions of the kind 

% $$
% dX_t = a(t,X_t) dt + b(t,X_T) dW_t \, \, \, \text{and} \, \, \, X(\omega,0) = X_0 \, \forall \omega \in \Omega.
% $$

% where $dW_t$ will be specified in the next subsection. The generalization aims to be in the sense that $X_t$ is the solution to the ODE $dY_t = a(t,Y_t) dt$ if $b = 0.$

% First, we will give sense to such expressions from an integral point of view, constructing concrete meaning for these. Then, we will show how a new and somewhat \textit{bizarre} Chain Rule arises in this context, and how to solve these new differential equations.

% \subsection{The Itô Integral}

% First, let's restate the problem as 

% $$
% X_t - X_0 = \int_0^t a(t,X_t) dt + \int_0^t b(t,X_T) dW_t 
% $$
% For each $\omega \in \Omega$, the first term in the right handside is nothing new. Regarding the second one, we may wonder if there is a sensible way of defining 
% $$
% M_t = \int_0^tf(s,\omega) dW(s,\omega), \, \, \, \, 0\leq t\leq T
% $$
% so that $M_t$ is a martingale. There is another witty consideration that leads to the definition of Itô Calculus, which is the necessity that the integrand is adapted in some sense to the filtration considered. These ideas lead to the following definition:

% \begin{definition}
    
% \begin{itemize}
%     \item[i)] The space of stochastic processes on $[0,T]$ adapted to a filtration $\mathcal{F}_t$ and such that $\int_0^T\mathbb{E}[|f(t)|^2]dt < \infty$ will be denoted by $L^2_{ad}([0,T]\times \Omega).$ It is a Hilbert space, the standard one that arises in the finite-measure space $[0,T]\times \Omega.$ 
%     \item[ii)] We say $f\in L^2_{ad}([,b]\times\Omega)$ is a \textbf{step stochastic process} if $$f(t,\omega) = \sum_{i= 1}^n \varphi_{i-1}(\omega)\mathds{1}_{[t_{i-1},t_i)}$$
%     where $\varphi_i$ is a collection of $\mathcal{F}_{t_i}$-measurable and of finite variance ($\mathbb{E}[\varphi_i^2] < \infty$).
% \end{itemize}
    
    
% \end{definition}

% The idea behind the following construction relies on working carefully with these step functions and approximating up to its closure in $L_{ad}^2([0,T]\times\Omega).$ For any such process $f$, take $t_0 = a$ and $t_n = b$ and \textbf{define}\footnote{One thing that must be noticed is why we take the left evaluation on each subinterval. It is the martingale property that motivates this. There are however variants of this kind of Calculus, where the evaluation point is different and different calculus rules apply. See Stratonovich Caclulus in (REFERENCE).} \textbf{its integral} as
% $$
% \int_0^T fdB_t = I(f) = \sum_{i=1}^n \varphi_{t_{i-1}}\left(  B(t_i)-B(t_{i-1}\right)(\omega).
% $$
% This transformation is clearly linear. Additionally, the following two lemmas are key for the construction for any $f\in L^2_{ad}(a,b\times\Omega).$

% \begin{lemma}

% Under these conditions, for $f$ step stochastic, $\mathbb{E}[I(f)]=0$ and $$\mathbb{E}[I(f)^2]= \int_0^T \mathbb{E}[f(t)^2]dt.$$
% \end{lemma}
% \textit{\textbf{Proof:}} (...) \hfill $\blacksquare$

% \begin{lemma}
% Suppose $f\in L^2_{ad}(a,b\times\Omega).$ There exists a sequence $\{f_n\}_{n= 1}^\infty$ of step stochastic processes in $L^2_{ad}([0,T]\times\Omega)$ such that
% $$
% \lim_{n\rightarrow \infty}\int_0^T \mathbb{E}[(f(t)-f_n(t))^2]dt =0
% $$    
% which amounts to say that these step processes are $dense$ in $L^2_{ad}([0,T]\times\Omega).$
% \end{lemma}

% \textbf{\textit{Proof:}} (...) \hfill $\blacksquare$

% We are know ready to define this to every process in $L^2_{ad}(...).$ Note that, in the conditions of the previous lemma, if $\{f_n\}_{n=1}^\infty$ is any such approximating sequence, then
% $$
% \mathbb{E}[(I(f_n)-I(f_m))^2] = \int_0^T \mathbb{E}[(f_n-f_m)^2] \rightarrow 0 \text{ as } n,m\rightarrow \infty
% $$
% since the righthandside is convergent, thus Cauchy condition holds. Then ${I(f_n)}_{n=1}^\infty$ is a Cauchy sequence in $L^2(\Omega)$ and, being a complete space, this leads to the following definition.

% \begin{definition}{Stochastic Integral}
%     Let $f \in L^2_{ad}(...)$. We define its \textit{stochastic integral} over $[0,T]$, namely $I(f)$ as the limit in $L^2(\Omega)$ of $I(f_n)$ for any sequence of step stocahstic processes such that $f_n \rightarrow_{L^2_{ad([0,T]\times\Omega)}} f.$
% \end{definition}

% It is linear. One of the previous lemmas must be stated now, together with an immediate consequece, as a legit Theorem.

% \begin{theorem}(Itô's isometry)
% \begin{itemize}
%     \item[i)]

% Suppose $f\in L^2_{ad}(...). $ Then $I(f)$ is a centered random variable such that

% $$
% \mathbb{E}[I(f)^2] = \int _a ^b \mathbb{E}[f(t)^2]dt.
% $$

% \item[ii)] For any $f,g \in L^2_{ad}(...)$, 
% $$
% \mathbb{E}\left[   \int_0^T f(t)dB(t) \int_0^Tg(t)dB(t)\right] = \int_0^T \mathbb{E}[f(t)g(t)]dt.
% $$
% \end{itemize}    
% \end{theorem}





% \textit{\textbf{Proof:}} (...) \hfill $\blacksquare$

% Also, note that whenever $f(t)$ does not show any dependence on $\omega$, then $I(f)$ is a centered normal variable with variance $\int_0^T f(t)^2 dt.$


\section{Stochastic Calculus}\label{sec:stochasticcalculus}

Stochastic Calculus is a wide area in Mathematics that allows, for example, the incorporation of randomness and uncertainty into differential equations. One of our main aims is to give meaning to expressions of the kind 
\begin{equation}\label{eq:sde_intro}
dX_t = a(t,X_t)\,dt + b(t,X_t)\,dW_t, 
\quad X_0 \in \mathbb{R}.
\end{equation}
Here $dW_t$ will be specified in the next subsection. Notice that this generalizes the deterministic case: if $b \equiv 0$, then $X_t$ reduces to the solution of the ODE $dY_t = a(t,Y_t)\,dt$.

We shall first construct a rigorous integral interpretation of such expressions, providing concrete meaning to the stochastic term. Afterwards, we will introduce the stochastic chain rule (Itô's Lemma), which differs from the classical one in a somewhat surprising way, and then proceed to study the existence and uniqueness of solutions, and some of their properties.

It must be highlighted that the following dissertation can be expanded up to the most general setting, however we will stick to the most classical presentation of the subject, leaving the most general case a worthy space at the end, but without letting the trees hide the forest.

\subsection{The Itô Integral}\label{subsec:itointegral}

The stochastic differential equation \eqref{eq:sde_intro} can be reformulated in integral form as
\begin{equation}\label{eq:sde_integral}
X_t - X_0 = \int_0^t a(s,X_s)\,ds \;+\; \int_0^t b(s,X_s)\,dW_s.
\end{equation}
For each $\omega \in \Omega$, the first integral on the right-hand side is nothing new. Regarding the second one, we may ask whether there is a sensible way of defining
\[
M_t = \int_0^t f(s,\omega)\,dW_s, \quad 0 \leq t \leq T,
\]
so that $\{M_t\}_{t\in[0,T]}$ is a martingale. Another important consideration is that the integrand must be \emph{adapted} to the filtration associated with $W_t$. These requirements naturally lead to the following definitions.

\begin{definition}[Adapted square-integrable processes]\label{def:L2_adapted}
Let ($\Omega, \mathcal F, \{\mathcal F_t\}, \mathbb P$) be a filtered probability space.
\begin{itemize}
    \item[i)] The space of stochastic processes on $[0,T]$, adapted to a filtration $(\mathcal{F}_t)_{t\geq 0}$ and such that 
    \[
    \int_0^T \mathbb{E}[|f(t)|^2]\,dt < \infty,
    \]
    will be denoted by $L^2_{ad}([0,T]\times \Omega)$. This is a Hilbert space with the inner product
    \[
    \langle f,g \rangle = \int_0^T \mathbb{E}[f(t)g(t)]\,dt.
    \]

    \item[ii)] We say $f \in L^2_{ad}([0,T]\times \Omega)$ is a \textbf{step stochastic process} if
    \[
    f(t,\omega) = \sum_{i=1}^n \varphi_{\omega}(t_{i-1})\,\mathds{1}_{[t_{i-1},t_i)}(t),
    \]
    where $\varphi(t_{i-1})$ is $\mathcal{F}_{t_{i-1}}$-measurable and has finite variance ($\mathbb{E}[\varphi(t_{i-1})^2] < \infty$).
\end{itemize}
\end{definition}

\medskip

The construction of the stochastic integral begins by defining it on step processes and extending by density. For such $f$, take $t_0=a$, $t_n=b$, and define\footnote{The choice of the left endpoint $t_{i-1}$ ensures that the stochastic integral preserves the martingale property. Variants with different evaluation points lead to different calculi, e.g. Stratonovich calculus, where the classical chain rule holds (see \cite{Kuo2006}).}
\begin{equation}\label{eq:integral_step}
I(f) = \int_0^T f(t)\,dW_t := \sum_{i=1}^n \varphi(t_{i-1})\,\big(W(t_i)-W(t_{i-1})\big).
\end{equation}
This mapping $f \mapsto I(f)$ is linear. The following lemmas are key.

\begin{lemma}\label{lem:isometry_step}
If $f$ is a step stochastic process, then
\[
\mathbb{E}[I(f)] = 0, 
\quad 
\mathbb{E}[I(f)^2] = \int_0^T \mathbb{E}[f(t)^2]\,dt.
\]
\end{lemma}
\textit{\textbf{Proof}:} (...)() \hfill $\blacksquare$

\begin{lemma}[Density of step processes]\label{lem:density}
For every $f \in L^2_{ad}([0,T]\times \Omega)$, there exists a sequence $\{f_n\}_{n=1}^\infty$ of step stochastic processes such that
\[
\lim_{n\to\infty}\int_0^T \mathbb{E}[(f(t)-f_n(t))^2]\,dt = 0.
\]
That is, step processes are dense in $L^2_{ad}([0,T]\times \Omega)$.
\end{lemma}
\textit{\textbf{Proof}:} (...) \hfill $\blacksquare$

\medskip

We are now ready to extend the definition.

\begin{definition}[Stochastic Integral]\label{def:stochastic_integral}
Let $f \in L^2_{ad}([0,T]\times\Omega)$. We define its \emph{stochastic integral} over $[0,T]$ as the limit in $L^2(\Omega)$ of $I(f_n)$ for any approximating sequence of step processes $\{f_n\}$ with $f_n \to f$ in $L^2_{ad}([0,T]\times\Omega)$.
\end{definition}

The following fundamental result follows.

\begin{theorem}[Itô's isometry]\label{thm:itoisometry}

Let $f \in L^2_{ad}([0,T]\times \Omega)$.
\begin{itemize}
    \item[i)] The stochastic integral $I(f)$ is a centered random variable and
    \[
    \mathbb{E}[I(f)^2] = \int_0^T \mathbb{E}[f(t)^2]\,dt.
    \]

    \item[ii)] For another $g \in L^2_{ad}([0,T]\times\Omega)$,
    \[
    \mathbb{E}\left[\left(\int_0^T f(t)\,dW_t\right)\left(\int_0^T g(t)\,dW_t\right)\right]
    = \int_0^T \mathbb{E}[f(t)g(t)]\,dt.
    \]
\end{itemize}
\end{theorem}
\textit{\textbf{Proof}:} (...) \hfill $\blacksquare$

It is worth noting that if $f$ is deterministic (i.e. does not depend on $\omega$), then
\[
I(f) \sim \mathcal{N}\!\left(0, \int_0^T f(t)^2\,dt\right).
\]


For $0\leq t \leq T$, we have $\int_0^t \mathbb{E}[f(s)^2]\,ds \leq \int_0^T \mathbb{E}[f(s)^2]\,ds$, so $f\in L^2_{ad}([0,t]\times \Omega)$ and the stochastic integral is well defined on $[0,t]$. Consider then the process $X_t = \int_0^t f(s) dW(s) .$ We know this process has finite first two moments, and so we can take the conditional expectation of $X_t$ with respect to the given filtration. 



\begin{theorem}[Martingale property]\label{theorem:martingale_property}\label{theorem:MartingalePropertyIntegration}

Take $f\in L ^2_{ad}([0,T]\times\Omega). $ Then the process given by $X_t = I(f \mathds{1}_{[0,t]}) $ is a martingale with respect to the   filtration $\mathcal{F}_t.$
    
\end{theorem}


\textbf{\textit{Proof:}} (...) \hfill $\blacksquare$

The Markov property also holds, but we will return to that when we see these processes as diffusions in Subsection \ref{section:itoDiff}. 

We shall now address the continuity of the process $X_t$. We must keep in mind that this integral is not defined for each $\omega$ in a Riemann sense. The following result is therefore not as elementary as in real analysis.

\begin{theorem}[Continuity property]\label{thm:itocontinuity}
    Suppose $f\in L ^2_{ad}([0,T]\times\Omega). $ Then, almost all paths of $X_t = \int_0^t f(s)dW(s)$ are continuous on $[0,T].$
\end{theorem}
\textbf{\textit{Proof:}} (...) \hfill $\blacksquare$


Finally, a result that actually links Riemann sums with stochastic integrals will be helpful in the next section.


\begin{theorem}\label{theorem:riemannian_limit}
    Suppose $f\in L ^2_{ad}([0,T]\times\Omega)$ and assume that $\mathbb{E}[f(t)f(s)]$ is a continuous function of $t$ and $s$. Then, for  partitions of $[0,T] $, given by $\mathcal{P}=\{0=t_0<t_1<\cdots<t_n=T\}$,

\begin{equation}
\int_0^T f(t)\, dW(t)
= \lim_{\|\mathcal P\|\to 0} \sum_{i=1}^n f(t_{i-1})\bigl(W(t_i)-W(t_{i-1})\bigr)
\end{equation}
in $L^2(\Omega),$ where the limit is taken as in Definition \ref{def:quadratic-variation}.
\end{theorem}
\textbf{\textit{Proof:}} (...) \hfill $\blacksquare$

More genearlly, we can consider a larger family of integrands. Let $\mathcal{L}(\Omega,L^2[0,T])$ be the family of processes $f$ such that
\begin{itemize}
    \item $f(t)$ is $\mathcal{F}_t-$adapted.
    \item $\int_0^T |f(t)| ^2 \, dt < \infty$ almost surely.
\end{itemize}

Notice how we dropped the condition of having finite expectation, so that we have indeed a larger class of processes. Consider for example the process $f(t) =e^{W(t)^2}.$ It follows that 

\begin{empheq}[left={\mathbb{E}[f(t)^2]= \empheqlbrace}]{equation}
  \begin{aligned}
      &\frac 1 {\sqrt{1-4t}} & & \text{if } 0\leq t< 1/4\\[1ex]
      & \infty & & \text{if } t \geq 1/4.
  \end{aligned}
\end{empheq}

Hence $\int_ 0 ^1 e^{W(t)^2 \, dW_t}$ does not make sense with the previous construction. It can be shown, however, that these processes are still pretty worth the attention. Beyond the scope of our purposes, it can be shown that:

\begin{theorem}(Local Martingale Property) Let $f\in \mathcal{L} (\Omega,L^2 [0,T])$ and $t \in [0,T]$. Then 
$$
X_ t = \int_0^t f(s) \,dW_s
$$
    is a local martingale. Furthermore, it has a continuous realization. \hfill $\blacksquare$
\end{theorem}



\subsection{Itô's Chain Rule}

One of the most shocking facts about this new integration is a sophisticated version of  the Chain Rule, that generalizes the familiar expression $d (f(g))(x) = f'(g(x))dg(x)$. The key underlying fact for this new integration rule is the non-zero quadratic variation of Brownian motion. This fantastic new calculus allows us to find some new relations and solve new equations that naturally generalize the already rich world of differential equations.

\begin{definition}(Itô process)

Given an initial random variable $X_0$, an Itô process is a stochastic process on $[0,T]$ of the form 

$$
X_t = X_0 + \int_0^t g(s) \,ds  +\int_0^t f(s) \, dW_s  \iff dX_t = g(s) \,ds + f(s) \, dW_s, \text{ } \text{ } X(0) = X_0
$$
    where we emphasize the equivalent formulations, since the latter is just a way of writing the former. We impose $X_0$ being $\mathcal F_0$-measurable, $g \in \mathcal{L}(\Omega,L^1 [0,T])$ and $f \in \mathcal{L}(\Omega,L^2 [0,T]).$ Note how the first integral is just a Lebesgue integral for each $\omega$, that is, for each path of $W$. The second is a proper Itô integral.
\end{definition}

We directly state the more general form of Itô's new calculus rule, in the convenient \textit{differential formulation}. A complete proof is beyond the scope of this Thesis. However, such a renowned Theorem has been addressed by several authors, for example in \cite{Kuo2006}.



\begin{theorem}[Itô's Rule]\label{theorem:itorule}
    Let $dX_t = \mu(t)\,dt + \sigma(t)\,dW_t$ be an Itô process, and let 
    $V(t,x)$ be a continuous function on $[0,T] \times \mathbb{R}$ with continuous partial derivatives 
    $\frac{\partial V}{\partial t}$, $\frac{\partial V}{\partial x}$ and $\frac{\partial^2 V}{\partial x^2}$.
    Then $V(t,X_t)$ is also an Itô process, and
    \begin{equation}\label{eq:ito_formula}
        dV(t,X_t) 
        = \left(
            \frac{\partial V}{\partial t}(t,X_t)
            + \mu(t)\frac{\partial V}{\partial x}(t,X_t)
            + \frac{1}{2}\sigma^2(t)\frac{\partial^2 V}{\partial x^2}(t,X_t)
        \right) dt
        + \sigma(t)\frac{\partial V}{\partial x}(t,X_t)\,dW_t.
    \end{equation}
\end{theorem}

\noindent
A convenient way to derive Equation~\eqref{eq:ito_formula} is through a symbolic
derivation using the Taylor expansion to first order in $dt$ and second order in $dX_t$,
together with the \textit{Itô table}:
\[
\begin{array}{c|cc}\label{ItoTable}
\times & dW_t & dt \\ \hline
dW_t & dt & 0 \\
dt & 0 & 0
\end{array}
\]


\noindent
Applying the Taylor expansion to $V(t,X_t)$ gives
\[
dV(t,X_t)
= \frac{\partial V}{\partial t}(t,X_t)\,dt
+ \frac{\partial V}{\partial x}(t,X_t)\,dX_t
+ \frac{1}{2}\frac{\partial^2 V}{\partial x^2}(t,X_t)\,(dX_t)^2.
\]
Using the Itô table, we have $(dX_t)^2 = \sigma(t)^2\,dt$, so using
$dX_t = \mu(t)\,dt + \sigma(t)\,dW_t$ yields
\[
dV(t,X_t)
= \frac{\partial V}{\partial x}(t,X_t)\sigma(t)\,dW_t
+ \left[
\frac{\partial V}{\partial t}(t,X_t)
+ \mu(t)\frac{\partial V}{\partial x}(t,X_t)
+ \frac{1}{2}\sigma(t)^2\frac{\partial^2 V}{\partial x^2}(t,X_t)
\right] dt,
\]
which is precisely Equation~\eqref{eq:ito_formula}.

We stress that the computation just performed is a \emph{mnemonic}, not a proof: the symbols $dW_t$ and $(dW_t)^2$ have been manipulated as though they were infinitesimals, which they are not. What makes the manipulation legitimate is the content of \eqref{eq:bmqv} --- the second-order term survives precisely because $W$ has non-vanishing quadratic variation --- and a genuine proof proceeds by Taylor-expanding along a partition and controlling the remainder in $L^2(\Omega)$. We refer to \cite{Kuo2006} for the details. The value of the Itô table is that it makes the bookkeeping automatic once one accepts that $\langle W \rangle_t = t$, and we shall use it in that spirit throughout.

Just like ordinary Chain Rule, this has opened plenty of possibilities in order to calculate some stochastic integrals in terms of others.

\begin{example}
 We can use this result to easily obtain a simpler form for the process   $W_t\,dW_t$. We can simply consider $g(x) = x^2$ and define $H_t = g(W_t)$ to obtain
 \begin{equation}
     dH_t = 2W_t\,dW_t + \frac 2 2 \,(dW_t)^2 = dt + 2 W_t\,dW_t
 \end{equation}
 so that $\int_0^t W_s\,dW_s = \frac 1 2 \left( W_t^2 - t\right)$, which is a much simpler expression, and clearly a martingale.
\end{example}

\begin{example}\label{example:blackscholessde}
    In Finance, the Black-Scholes model relies on the \textit{log-normal} equation, given by $dS_t = rS_t \, dt + \sigma S_t \, dW_t$, where $r, \sigma \in \mathbb R$ and $\sigma > 0$. It can be solved by taking $H_t = \ln (S_t),$ so that
    \begin{equation}
        dH_t = \frac 1 S_t dS_t -\frac{1}{2S_t^2}(dS_t)^2 = r\,dt + \sigma\, dW_t -  \frac 1 2 \sigma^2 \, dt
    \end{equation}

    obtaining $H_t = H_0 + (r-\sigma^2 / 2) \, t + \sigma W_t \implies S_t = S_0e^{(r-\sigma^2/ 2)\, t + \sigma \, W_t} $.
\end{example}



\begin{example}(Lamperti's Transform) \hl{(NOTA PARA CARLOS: Estoy usando esta transformación en el BBVA :)}


     Consider a general process  $dX_t = \mu(t,X_t)\, dt + \sigma(t,X_t) \, dW_t$ and assume $\sigma$ is differentiable once with respect to time and twice with respect to $x$. Define the function 
    \begin{equation}
        \Gamma(t,x) = \int_{X_0}^x \frac{1}{\sigma(t,s)}\, ds
    \end{equation}
and consider the process $H_t = \Gamma(t,X_t).$ After a few calculations, we arrive to
\begin{equation}
    \begin{aligned}
        dH_t &= \left( \frac{\partial \Gamma}{\partial t} + \mu(t,X_t)\frac{\partial \Gamma}{\partial x} + \frac{1}{2}\sigma^2(t,X_t)\frac{\partial^2 \Gamma}{\partial x^2} \right) dt + \underbrace{\sigma(t,X_t)\frac{\partial \Gamma}{\partial x}}_{=1} \, dW_t \\
        &= \left( \frac{\partial \Gamma}{\partial t} + \frac{\mu(t,X_t)}{\sigma(t,X_t)} - \frac{1}{2}\frac{\partial \sigma}{\partial x}(t,X_t) \right) dt + dW_t
    \end{aligned}
\end{equation}
so that all the complexity of the original process is contained in the drift. This transformation is convenient because it creates a process of constant volatility, which enables much faster and accurate numerical methods, such as derivative pricing on recombining trees. \todonote{cross-ref to Ch. 3 numerics}
\end{example}


% \textbf{\textit{Proof:}} (...)

%  The proof follows from two technical lemmas. First, note that if $f\in\mathcal{C}^2([0,T])$ and consider a partition $\{a=t_0,t_1,\dots,t_n=t\}$ to obtain, due to Taylor Theorem, an expression of the form

% \begin{equation}
% f(W_t) - f(W_a = 
% \sum_{i=1}^{n} f'(W_{t_{i-1}})\,\bigl(W_{t_i} - W_{t_{i-1}}\bigr)  + \frac{1}{2} \sum_{i=1}^{n} 
% f''\!\left(\chi_i\right) 
% \bigl(W_{t_i} - W_{t_{i-1}}\bigr)^{2}
% \end{equation}
% where $\chi_i = W_{t_{i-1}} + \lambda_i \bigl(W_{t_i} - W_{t_{i-1}}\bigr)$ for some $\lambda_i\in(0,1).$

% The first summation has already been addressed in $REFERENCE,.5.3.3;7.1.3;4.7.1$

% For the second summation term, it seems reasonable to assess the limit should be 
% \begin{equation}
% %\lim_{\lVert \mathcal{P} \rVert \rightarrow 0}
% \sum_{i=1}^{n} 
% f''\!\left(\chi_i\right) 
% \bigl(W_{t_i} - W_{t_{i-1}}\bigr)^{2} \rightarrow \int_0^t f''(W(s)) \, ds
% \end{equation}
% in probability as ${\lVert \mathcal{P} \rVert \rightarrow 0}$, that is, naming $Z_\mathcal{P}$ the summation associated to $\mathcal{P}$, $\forall \varepsilon >0$ there exists $\delta >0$ such that if $\lVert \mathcal{P} \rVert <\delta$ then $\mathbb{P}(|Z_\mathcal{P}-\int_0^t f''(W_s) ds|>\varepsilon)<\varepsilon$. Additionally, this condition is equivalent to the property of having, for any sequence of partitions $\{\mathcal{P}_{n}\}_{n=1}^\infty$ such that $\lim _ {n \rightarrow \infty} \lVert \mathcal{P}_n \rVert = 0$, a subsequence that makes $Z_{\mathcal{P}_{n_k}}$ converge almost surely. In order to obtain (REFERENCE EQUATION), note that

% \begin{equation}
% \begin{aligned}
% \sum_{i=1}^{n} f''(\chi_i) (W_{t_i} - W_{t_{i-1}})^2 &= \sum_{i=1}^{n} \Big( f''(\chi_i) - f''(W_{t_{i-1}}) \Big) (W_{t_i} - W_{t_{i-1}})^2 \\
% &\quad + \sum_{i=1}^{n} f''(W_{t_{i-1}}) \Big( (W_{t_i} - W_{t_{i-1}})^2 - (t_i - t_{i-1}) \Big) \\
% &\quad + \sum_{i=1}^{n} f''(W_{t_{i-1}}) (t_i - t_{i-1}).
% \end{aligned}
% \end{equation}

% The two following lemmas address the first two summations.

% \begin{lemma}

% Let $g(x)$ be a continuous function on $\mathbb R$. For each $n\geq 1$, let $\Delta_n = \{t_0, t_1, \dots, t_n\}$ be a partition of $[a,t]$ and let $0< \lambda_i < 1$ for $1\leq i\leq n.$ There exists a subsequence of 

% \begin{equation}
%     \sum_{i=1}^n \left( g(B_{t_{i-1}}+\lambda_i (B_{t_i}-B-{t_{i-1}}) - g(B_{t_{i-1}})  \right) (B_{t_i}-B_{t_{i-1}})^2
% \end{equation}
% that converges to 0 almost surely as $\lVert \Delta_n\rVert \rightarrow 0.$
% \end{lemma}

% \textbf{\textit{Proof:}}  (...)

% \begin{lemma}
% Let $g(x)$ be a continuous function on $\mathbb{R}.$ For each $n\geq 1$, let $\Delta_n$ be again a partition of $[a,t]$. Then the sequence 

% \begin{equation}
% \sum_{i=1}^{n} g(B_{t_{i-1}})\,\bigl(B_{t_i} - B_{t_{i-1}}\bigr)^{2} - (t_i - t_{i-1})
% \;\;\xrightarrow[\Delta_n \to 0]{}\; 0 \quad \text{in probability}.
% \end{equation}


% INCLUDE HERE THE RELATION WITH THE PREVIOUS RESULT; GENERAL DIFFERNETIAL FORMULA


% \end{lemma}

% \textbf{\textit{Proof:}}  






\subsection{Stochastic Differential Equations}\label{subsec:sdes}

The main aim of this section is to provide necessary and sufficient conditions for the existence of Itô processes. We will follow a direct approach and state the existence and uniqueness theorem. In addition, notice how we will be working with the differential form of the equations, however realising the integral sense they have been provided with. We need a precise definition for \textit{stochastic differential equation} or $SDE$ and thus some preliminary results.

Given a $\mathcal{F}_a-$measurable random variable $\Theta$ of finite variance and consider the following formal equation on the interval $[0,t]$: 

\begin{equation}\label{sde_template}
    dX_t = \mu(t,X_t) \, dt + \sigma ( t, X_t) \, dW_t \text{ }\text{ }\text{ }\text{ }\text{ }\text{ } X(0)= \Theta.
\end{equation}

\begin{definition}
    A stochastic process $X_t$, $0\leq t \leq T$, is called a solution of \ref{sde_template} if it satisfies:

    \begin{itemize}
        \item $\sigma(t,X_t) \in \mathcal L (\Omega,L^2[0,T])$ so its Itô integral makes sense for each $t\in [0,T].$
        \item Almost all sample paths of $\mu(t,X_t)$ belong to $L^1[0,T]$, that is, its expected value is finite.
        \item Equation \ref{sde_template} holds almost surely.
     \end{itemize}
\end{definition}

We need to establish some growth conditions on  $\mu$ and $\sigma$ for the to make sense on, at least, an open interval, just like in ordinary differential equations. Let us define the following conditions:

\begin{definition}\label{def:lipschitzgrowth}
    Let $g(t,x)$  a measurable function on $[0,T] \times \mathbb R$.

    \begin{itemize}
        \item It is said to satisfy a Lipschitz condition in $x$ if there exists $K>0$ such that $|g(t,x)-g(t,y) |\leq K |x-y| $ for all $0 \leq t \leq T,$ $x,y\in\mathbb{R}.$
        \item It is said to satisfy a linear growth condition in x if there exists $K>0$ such that $|g(t,x)|\leq K(1+|x|)$ for all $0 \leq t \leq T,$ $x\in\mathbb{R}.$
    \end{itemize}
\end{definition}

Now, notice that for $x\geq 0$, it holds that $1+x^2 \leq (1+x) ^2 \leq 2(1+x^2)$  so that the linear growth condition can be restated in terms of a new constant $M$ such that $|g(t,x)| \leq M \sqrt{1+x^2}.$ 

The proof of existence and uniqueness rests on a Picard iteration, exactly as for ordinary differential equations, and the comparison step is furnished by Grönwall's inequality (Lemma \ref{lem:gronwall} in Appendix \ref{app:auxiliary}); since the argument is deterministic and entirely standard we do not reproduce it here.

\begin{theorem}\label{thm:sdeexistence}
    Let $\mu$ and $\sigma$ be measurable functions on $[0,T] \times \mathbb{R}$ satisfying both the Lipschitz and the linear growth conditions in $x$ of Definition \ref{def:lipschitzgrowth}, and let the initial condition $\Theta$ be $\mathcal{F}_0$-measurable with $\mathbb{E}[\Theta^2] < \infty$. Then the stochastic differential equation \ref{sde_template} admits a solution, which is unique in the following sense: any two solutions are indistinguishable. Moreover the solution has continuous sample paths.
\end{theorem}
\textbf{\textit{Proof:}}
We shall only prove the uniqueness part. For that purpose, (...) $\blacksquare$

Three comments on the hypotheses, none of which is pedantry.

First, both conditions are needed, and they do different jobs: the Lipschitz condition delivers uniqueness, whereas the linear growth condition prevents the solution from exploding in finite time --- exactly as $\dot{y}=y^2$ does in the deterministic theory. If $\mu$ and $\sigma$ are globally Lipschitz in $x$ uniformly in $t$, and $|\mu(t,0)|$ and $|\sigma(t,0)|$ are bounded on $[0,T]$, then linear growth follows automatically, so the second hypothesis is often not stated separately; we state it because the multidimensional version, Theorem \ref{theorem:solutions:ndimensional_sde}, does.

Second, the square-integrability of the initial condition is not decoration. The whole construction of the Itô integral took place in $L^2$, and the Picard iterates are estimated there; without $\mathbb{E}[\Theta^2]<\infty$ the very first iterate need not lie in the space where the argument is conducted.

Third, and most consequentially for what follows, the word ``unique'' conceals a distinction that this formulation is too coarse to see. We have asserted \emph{pathwise} uniqueness: given the Brownian motion, the solution is determined. One may instead ask only that the \emph{law} of the solution be determined, which is a weaker requirement and corresponds to the notion of a \emph{weak} solution --- one where the driving Brownian motion is part of the output rather than the input. The two notions genuinely differ, and the gap between them is not a curiosity: the Lipschitz hypothesis above is violated by one of the most widely used models in this dissertation, the square-root variance process of Subsection \ref{subsec:stochvol}, whose diffusion coefficient $\xi\sqrt{v}$ is only Hölder continuous of order $\tfrac12$ at the origin. Theorem \ref{thm:sdeexistence} as stated therefore does \emph{not} cover the Heston model. We now repair this.

\subsubsection*{Beyond Lipschitz: strong and weak solutions}

Let us first make the distinction precise, since it is systematically blurred in the applied 
literature and the blurring causes real confusion.

\begin{definition}[Strong and weak solutions]\label{def:strongweak}
Consider the equation $dX_t = \mu(t,X_t)\,dt + \sigma(t,X_t)\,dW_t$ with $X_0 = \Theta$.
\begin{itemize}
    \item[i)] A \textbf{strong solution} is a process $X$, adapted to the augmented filtration 
    generated by a \emph{given} Brownian motion $W$ and the initial condition, satisfying the 
    equation almost surely. \emph{Pathwise uniqueness} holds if any two strong solutions on the same 
    space, driven by the same $W$, are indistinguishable.
    \item[ii)] A \textbf{weak solution} is a pair $(X,W)$ together with a filtered probability space, 
    such that $W$ is a Brownian motion for that filtration and the equation holds. 
    \emph{Uniqueness in law} holds if any two weak solutions induce the same law on $C[0,T]$.
\end{itemize}
\end{definition}

The distinction is one of information. For a strong solution the noise is handed to us and the 
trajectory is a measurable functional of it --- one may, at least in principle, simulate the 
solution from a given stream of random numbers. For a weak solution we are permitted to construct 
the noise ourselves, and all we control is the distribution of the result. For pricing this is often 
enough, since prices are expectations and expectations depend only on the law; for hedging, and for 
any pathwise statement, it is not.

The relation between the four notions is settled by the Yamada--Watanabe theorem: pathwise 
uniqueness together with existence of a weak solution implies existence of a strong solution and 
uniqueness in law. What we require is the companion criterion, which relaxes the Lipschitz condition 
in exactly the direction needed.

\begin{theorem}[Yamada--Watanabe]\label{thm:yamadawatanabe}
Consider the one-dimensional equation $dX_t = \mu(t,X_t)\,dt + \sigma(t,X_t)\,dW_t$. Suppose there 
exist a strictly increasing function $h:[0,\infty)\to[0,\infty)$ with $h(0)=0$ and
\begin{equation}\label{eq:yamadacondition}
\int_{0^+} \frac{du}{h(u)^2} = \infty,
\end{equation}
and a strictly increasing concave $\kappa$ with $\kappa(0)=0$ and $\int_{0^+}du/\kappa(u)=\infty$, 
such that for all $t$ and all $x,y$
\[
\bigl|\sigma(t,x)-\sigma(t,y)\bigr| \leq h\bigl(|x-y|\bigr), \qquad 
\bigl|\mu(t,x)-\mu(t,y)\bigr| \leq \kappa\bigl(|x-y|\bigr).
\]
Then pathwise uniqueness holds for the equation \cite{yamada1971}. \hfill $\blacksquare$
\end{theorem}

The point of \eqref{eq:yamadacondition} is that it admits $h(u) = u^{\gamma}$ for every 
$\gamma \geq \tfrac12$, since $\int_{0^+} u^{-2\gamma}du$ diverges precisely when 
$2\gamma \geq 1$. Hölder continuity of order $\tfrac12$ is thus permitted for the diffusion 
coefficient --- and only just, the exponent $\tfrac12$ being critical. The drift is treated 
separately and more leniently, needing only an Osgood-type modulus.

\begin{example}[The square-root process]\label{example:cirwellposed}
Consider the process which will serve as the variance in the Heston model of Subsection 
\ref{subsec:stochvol},
\begin{equation}\label{eq:cir}
dv_t = \kappa(\theta - v_t)\,dt + \xi\sqrt{v_t}\,dW_t, \qquad v_0 > 0,
\end{equation}
with $\kappa,\theta,\xi>0$. The drift is affine, hence Lipschitz. The diffusion coefficient 
$\sigma(v)=\xi\sqrt{v}$ is not Lipschitz at the origin, but it satisfies 
$|\sqrt{x}-\sqrt{y}| \leq \sqrt{|x-y|}$, so Theorem \ref{thm:yamadawatanabe} applies with 
$h(u)=\xi\sqrt{u}$, and pathwise uniqueness holds. A non-negative solution exists, and 
\eqref{eq:cir} is therefore well posed --- but by Yamada--Watanabe, not by Theorem 
\ref{thm:sdeexistence}.
\end{example}

Well-posedness is not the end of the story for \eqref{eq:cir}, because we also need the solution to 
stay strictly positive: it is to be a variance, and $\sqrt{v_t}$ must make sense. Note that the 
diffusion coefficient vanishes at $v=0$ while the drift there equals $\kappa\theta > 0$, so the 
origin is a point where the noise switches off and the drift pushes upward. Whether that upward push 
is strong enough is a quantitative question with a sharp answer.

\begin{proposition}[Feller condition]\label{prop:feller}
Let $v$ solve \eqref{eq:cir} with $v_0>0$. If
\begin{equation}\label{eq:fellercondition}
2\kappa\theta \geq \xi^2,
\end{equation}
then $\mathbb{P}(v_t > 0 \text{ for all } t \in [0,T]) = 1$. If $2\kappa\theta < \xi^2$, the process 
reaches $0$ with positive probability, though it is instantaneously reflecting and does not become 
negative \cite{feller1951}. \hfill $\blacksquare$
\end{proposition}

Condition \eqref{eq:fellercondition} has a transparent reading: mean reversion must be strong 
enough, and the volatility of volatility small enough, that the drift dominates the noise near the 
origin. It is worth knowing that in practice it is very often \emph{violated} by calibrated 
parameters --- fitting a steep short-maturity skew tends to demand a large $\xi$ --- and that this 
is a well-known source of difficulty in numerical work, since a naive Euler discretisation of 
\eqref{eq:cir} will produce negative values and then fail on $\sqrt{v_t}$. Practitioners resort to 
truncation or reflection schemes, or to exact simulation of the non-central $\chi^2$ transition law. 
We flag the point here because it is a genuine instance of the theory in this section constraining 
what can be done in Chapter \ref{ch:cap6}, rather than a technicality to be waved through.

\subsection{The multidimensional case}\label{subsec:multidimensional}

For the time being, we have only addressed processes in dimension 1. Almost every result already covered can be restated for multidimensional processes or has a direct analogue.

By a multidimensional Brownian motion on $\mathbb{R}^n$ we mean a stochastic process of the form  $\mathbf{W}(t) = (W_1(t),W_2(t),\dots,W_n(t))$ formed by $n$ independent Brownian motions on $\mathbb{R}$. Similarly, we can set them to start at $x_0\in \mathbb{R}^n$ just by a simple translation. It is a martingale and a Markov process that, starting at $\vec{0}$, has a transition density function given by $p_t(x) = (2\pi t)^{-n/2}e^{-\lVert x \rVert^2 / 2t}.$

There are certain properties of these processes that dramatically change with dimension. For example, for $n= 1$, it returns to 0 infinitely often, and for $n= 2$ enters arbitrarily small neighborhoods of $0$ infinitely often with probability 1, however the probability that it reaches $\vec 0$ is 0. For $n\geq 3$, the process is much more sparse, satisfying $\lim_{t\rightarrow \infty } \lVert \mathbf{W}(t) \rVert = \infty$ almost surely. 

In terms of fractal dimensions, the \emph{range} $\mathbf{W}([0,t])$ has Hausdorff dimension $2$ for $n \geq 2$, and dimension $1$ for $n=1$ (where it is simply an interval). The \emph{graph} $\{(s,\mathbf{W}(s)) : s\in [0,t]\} \subset \mathbb{R}^{1+n}$ behaves differently: it has dimension $3/2$ for $n = 1$, a classical result of Taylor \cite{taylor1953}, and dimension $2$ for $n \geq 2$. These fascinating properties of Brownian motion can be found in \cite{falconer}.

The exponent $3/2$ in the one-dimensional graph is not an isolated curiosity. A graph of dimension $3/2$ is exactly what one obtains from a function which is $\alpha$-Hölder for every $\alpha < 1/2$ and for no $\alpha \geq 1/2$, and that critical exponent $1/2$ --- together with what happens on either side of it --- will return in Part \ref{part:part2} as the organising theme of the whole dissertation.

We can therefore define a multidimensional stochastic integral (expressed in differential form) as
\begin{equation}
    dX_t = \mu \, dt + \sum_ {i= 1} ^n \sigma_ i \, dW_t^i
\end{equation}
where each integral must be considered separately. Let \(B_1(t), B_2(t), \ldots, B_m(t)\) be independent Brownian motions. Consider \(n\) Itô processes \(X_t^{(1)}, X_t^{(2)}, \ldots, X_t^{(n)}\) given by
\begin{equation}\label{equation:general_multidmensionalProcess}
X_t^{(i)} = X_0^{(i)} +  \int_0^t \mu_i(s) \, ds +\sum_{j=1}^m \int_0^t \sigma_{ij}(s) \, dW_j(s)  , \quad 1 \leq i \leq n,
\end{equation}
where \(\mu_{ij} \in \mathcal{L}(\Omega, L^2[0,T])\) and \(\sigma_i \in \mathcal{L}(\Omega, L^1[0,T])\) for all \(1 \leq i \leq n\) and \(1 \leq j \leq m\). If we introduce the matrices
\[
\vec W(t) = \begin{bmatrix} W_1(t) \\ \vdots \\ W_m(t) \end{bmatrix}, \quad
\vec X_t = \begin{bmatrix} X_t^{(1)} \\ \vdots \\ X_t^{(n)} \end{bmatrix},
\]
\[
\vec{\vec \sigma}(t) = \begin{bmatrix} \sigma_{11}(t) \cdots \sigma_{1m}(t) \\ \vdots \ddots \vdots \\ \sigma_{n1}(t) \cdots \sigma_{nm}(t) \end{bmatrix}, \quad
\vec \mu(t) = \begin{bmatrix} \mu_1(t) \\ \vdots \\ \mu_n(t) \end{bmatrix},
\]
then Equation \ref{equation:general_multidmensionalProcess} can be written as a matrix equation:
\[
\vec{X}_t = \vec{X}_0 +  \int_0^t \vec\mu(s) \, ds+\int_0^t  \vec{\vec\sigma}(s) \, d\vec W (s) , \quad 0 \leq t \leq T.
\]

We will usually drop the arrows if the dimension is clear. Additionally, we must impose that $\mu_i$ and $\sigma_ {i,j}$ satisfy certain growth conditions so that the solution makes sense. Itô's formula holds similarly for $\theta(\vec{X}_t)$, now accounting for the crossed second derivatives $\partial^2 \theta / \partial x_l \partial x_h$.

Before writing the corresponding multiplication table we must introduce an object which the table 
presupposes and which the definition above does not provide. So far, a multidimensional Brownian 
motion has \emph{independent} components. In modelling one almost always wants the components to be 
correlated --- a falling equity price accompanied by rising volatility, two rates moving together, 
an exchange rate and the foreign interest rate --- and the following makes that precise.

\begin{definition}[Correlated Brownian motion]\label{def:correlatedbm}
Let $P = (\rho_{ij}) \in \mathbb{R}^{n\times n}$ be symmetric and positive semi-definite with 
$\rho_{ii} = 1$ for all $i$. An $n$-dimensional process 
$\widetilde{\mathbf{W}} = (\widetilde W^1, \dots, \widetilde W^n)$ is a \textbf{Brownian motion with 
correlation matrix $P$} if each $\widetilde W^i$ is a one-dimensional Brownian motion and
\begin{equation}\label{eq:correlatedbm}
\bigl\langle \widetilde W^i, \widetilde W^j \bigr\rangle_t = \rho_{ij}\,t, \qquad 1 \leq i,j \leq n, \quad t \in [0,T].
\end{equation}
\end{definition}

Such a process always exists, and is obtained from an independent one by a linear map. Since $P$ is 
symmetric positive semi-definite it admits a factorisation $P = LL\transp$ --- the Cholesky 
factorisation when $P$ is positive definite, and a square root more generally. Setting 
$\widetilde{\mathbf{W}} := L\,\mathbf{W}$ for an independent $\mathbf{W}$, bilinearity of the 
quadratic covariation \eqref{eq:polarisation} gives
\[
\bigl\langle \widetilde W^i, \widetilde W^j \bigr\rangle_t 
= \sum_{k,l} L_{ik}L_{jl} \bigl\langle W^k, W^l \bigr\rangle_t 
= \sum_{k} L_{ik}L_{jk}\,t = (LL\transp)_{ij}\,t = \rho_{ij}\,t,
\]
and each component is a continuous local martingale with $\langle \widetilde W^i\rangle_t = t$, hence 
a Brownian motion by Theorem \ref{thm:levycharacterisation}.

The positive semi-definiteness of $P$ is a genuine constraint and not a formality. It is what makes 
the factorisation possible, and it is a real restriction as soon as $n \geq 3$: one cannot 
independently choose $\rho_{12}$, $\rho_{13}$ and $\rho_{23}$, since a strongly positive correlation 
between the first two pairs forbids a strongly negative one between the third. This is a familiar 
nuisance in practice, where correlation matrices estimated entrywise from market data frequently 
fail to be positive semi-definite and must be repaired before they can be used.

With this in hand, the Itô multiplication table extends as follows.

\[
\begin{array}{c|ccccc}
\times & dt & dW_t^1 & dW_t^2 & \cdots & dW_t^n \\ \hline
dt & 0 & 0 & 0 & \cdots & 0 \\
dW_t^1 & 0 & dt & \rho_{12}\,dt & \cdots & \rho_{1n}\,dt \\
dW_t^2 & 0 & \rho_{21}\,dt & dt & \cdots & \rho_{2n}\,dt \\
\vdots & \vdots & \vdots & \vdots & \ddots & \vdots \\
dW_t^n & 0 & \rho_{n1}\,dt & \rho_{n2}\,dt & \cdots & dt
\end{array}
\]
\captionof{table}{Itô multiplication table for correlated $n$-dimensional Brownian motion.}
\label{tab:ItoTable}

The entries are not a new convention but a restatement of \eqref{eq:correlatedbm}: the symbol 
$dW^i_t\,dW^j_t$ abbreviates $d\langle W^i,W^j\rangle_t = \rho_{ij}\,dt$. The independent case is 
recovered by taking $P = \mathbb{I}_{n\times n}$.

\begin{example}
    
Let $f : \mathbb{ R}^n \rightarrow \mathbb{ R} $ a  $\mathcal{C}^2$ function and let  $\nabla f = (\frac{\partial  f}{ \partial x_1}, \frac{\partial  f}{ \partial x_2}, \dots, \frac{\partial  f}{ \partial x_n})$ and $\Delta f = \sum_{i= 1}^n \frac {\partial^2 f}{\partial x_ i ^2}  $ to obtain 

$$
f(\mathbf{W}(t)) = f(\mathbf{W}(0)) + \int_0^t \nabla f (\mathbf{W}(s)) \cdot d\mathbf{W}_s + \frac{1}{2}\int_0 ^t \Delta f (\mathbf{W}(s)) \, ds.
$$
This equation exhibits the deep connection between Brownian motion and diffusions, showing how \emph{one half of the Laplace operator} is the generator of Brownian motion; see Subsection \ref{section:itoDiff}, where the same $\tfrac12$ reappears from the semigroup side. The factor is not decorative: it is the quadratic variation $\langle W^i \rangle_t = t$ entering through the second-order term of Itô's rule, and it is the reason the heat equation associated with Brownian motion is $\partial_t u = \tfrac12 \Delta u$ rather than $\partial_t u = \Delta u$.
\end{example}
\begin{example}
    An important example is given by $\theta(x,y) = xy.$ Clearly, $\partial \theta / \partial x = y$ and also $\partial \theta / \partial y = x$. Therefore, if we have processes $X_t$ and $Z_t$, we obtain
    \begin{equation}
        d(X_t\,Z_t) = Z_t\,dX_t + X_t\,dZ_t + \frac 1 2 dX_t\, dZ_t+ \frac 1 2 dX_t\, dZ_t=Z_t\,dX_t + X_t\,dZ_t + dX_t\, dZ_t
    \end{equation}
    what is usually known as Itô product formula, understanding $dX_t\,dZ_t$ as in table \ref{tab:ItoTable}.
\end{example}

The following generalization holds.

\begin{theorem} \label{theorem:solutions:ndimensional_sde}
    Let $\sigma(t,x)$ be an $n\times m$ matrix-valued function and $\mu(t,x)$ an $\mathbb{R}^n$-valued function. Assume there exists $C>0$ such that for all $t\in[0,T]$ and $x,y \in \mathbb R^n$, 
    $$
    \lVert \sigma (t,x) - \sigma(t,y) \rVert \leq C | x - y| \text{ } \text{ , } \text{ } |\mu(t,x) - \mu(t,y)| \leq C|x-y|
    $$
    and 
    $$
    \lVert \sigma(t,x)\rVert ^2 \leq C(1+|x|^2 ) \text{ } \text{ , } \text{ } |\mu(t,x)| ^2 \leq C(1+|x|^2 ) 
    $$ 
    where $\lVert \sigma \rVert^2 = \sum_i \sum_j \sigma_{ij}^2$. For any $\mathcal{F}_ a$-measurable random vector $\Theta$ of finite variance, the SDE $dX_t = \mu(t,X_t) \, dt +  \sigma(t,X_t) \, dW_t  $, $X_0 = \Theta$ has a unique continuous solution. \hfill $\blacksquare$
\end{theorem}


\subsection{Probability measure changes}\label{subsec:measurechanges}

Another important topic we must cover has to do with probability changes. Brownian motion depends crucially on the underlying probability $\mathbb{P}$. Given a process $X_t$, it seems legit to ask whether there is a probability $\mathbb{Q}$ so that $X_t$ is a Brownian motion with respect to $\mathbb{Q}$. We assume the $\sigma$-algebra and the filtration remain unchanged.

We say the filtration is \textit{right-continuous} provided
\begin{equation}\label{rightcontinuousFiltration}
    \cap_{n=1}^\infty \mathcal{F}_{t+1/n} = \mathcal{F}_t.
\end{equation}

\textit{Technical note:} In the following, we shall work with complete, right-continuous filtrations.

\begin{definition}(Predictable Process)
    Let $\{\mathcal F_t\}$ be a right continuous filtration and let $\mathbb{L}$ the collection  of stochastic processes adapted to $\{\mathcal F_t\}$ and such that almost all their sample paths are left continuous. Let $\Sigma$ the smallest $\sigma-$algebra of $[0,T]\times \Omega$ such that all processes in $\mathbb L$ are measurable. The family of processes that are $\Sigma-$measurable are called \textbf{predictable.}
\end{definition}

\begin{definition}(\textit{Càdlàg} Process)\label{def:cadlag}
We say that a stochastic process is \textbf{\textit{càdlàg}} provided its sample paths are right-continuous and have left limits (\textit{continue à droite, limite à gauche}).
\end{definition}

A word is owed on why predictability appears at all, since for the processes studied in this dissertation it will make no difference whatsoever.

When the integrator has continuous sample paths --- as Brownian motion does, and as every model in Part I does --- the distinction between predictable and merely adapted integrands is immaterial: the two classes give rise to the same stochastic integrals, and one may work with left-continuous adapted processes throughout without ever noticing. This is why the construction of Subsection \ref{subsec:itointegral} could proceed without the notion.

The distinction becomes essential the moment the integrator is allowed to jump, and it does so for a reason that is easy to state and worth stating, because it is genuinely about information rather than about technique. If the integrator jumps at a random time and the integrand is permitted to ``see'' that jump as it happens --- that is, to be adapted but not predictable --- then one may construct a strategy that bets on the jump at the instant of the jump, and the resulting integral fails to be a martingale. Predictability is precisely the requirement that the position be chosen strictly \emph{before} the jump is observed, using only information available up to, but not including, the present instant. In a financial reading this is not a technicality at all: it is the statement that one may not trade on information one does not yet have.

This matters in practice for underlyings whose price processes are genuinely discontinuous. Electricity is the standard example: spot power prices exhibit violent, short-lived spikes driven by demand peaks and generation outages, and no diffusion model reproduces them; the same is true, in milder form, for several commodities and for credit-sensitive instruments, where default is by nature a jump. Modelling such markets requires jump processes --- Poisson processes and their compensated versions being the elementary building blocks, as in Example \ref{example:poissonprocess} --- and there the machinery of predictable integrands is unavoidable. We record the definitions for completeness, and because the Doob-Meyer theorem below is stated in that language, but we shall not need them again.

These technical definitions play a role in the following fundamental theorem, that we state without proof. 



\begin{theorem}[Doob--Meyer Decomposition Theorem]\label{thm:doobmeyer}
    Let $M(t)$ , $t\in[0,T]$ be a \textit{càdlàg}, square integrable martingale. Then there exists a unique decomposition of the form
    \begin{equation}
        M_t ^2 = R(t) + \langle M_t\rangle
    \end{equation}
    where $R(t)$ is \textit{càdlàg} martingale and $\left<M_t\right>$ is a predictable, right continuous and increasing process such that $\left<M_0\right> = 0$ and $\mathbb E[\left<M_t\right>] < \infty$ for all $t\in[0,T].$ We call $\left<M_t\right>$ the \textbf{compensator} of $M_t^2$ and we say it is the predictable quadratic variation of $M_t$. \hfill $\blacksquare$
\end{theorem}

The theorem holds in a way more general form for submartingales, but if $M_t$ is a square integrable martingale, then, by Jensen's inequality, $M_t^2$ is a submartingale and it suffices for our purposes. For a proof of the Doob--Meyer decomposition theorem, see \cite{Kuo2006} or \cite{Shreve2004}.


\begin{theorem}[Lévy's Characterisation Theorem]\label{thm:levycharacterisation}
    Let $X$ be a continuous local martingale on a filtered probability space $\FPSpace$ with $X_0 = 0$ and $\langle X \rangle_t = t$ almost surely for every $t \in [0,T]$. Then $X$ is an $\Filt$-Brownian motion under $\mathbb{P}$.

    More generally, let $X = (X^1,\dots,X^n)$ be a vector of continuous local martingales with $X_0 = 0$ and
    \[
    \langle X^i, X^j \rangle_t = \delta_{ij}\,t \qquad \text{almost surely, for all } i,j \text{ and all } t,
    \]
    where $\delta_{ij}$ is the Kronecker delta. Then $X$ is an $n$-dimensional $\Filt$-Brownian motion with independent components.
\end{theorem}
\textbf{\textit{Proof:}} (...) $\blacksquare$

The statement deserves to be read slowly, because it is considerably stronger than it looks. Being a Brownian motion is, by Definition \ref{def:Brownian-motion}, a statement about the \emph{joint law} of the whole family $\{X_t\}$: independent, stationary, Gaussian increments. Lévy's theorem says that all of this is already implied by two facts about second moments --- that $X$ is a local martingale and that its quadratic variation is exactly $t$. Gaussianity is not assumed anywhere; it is a conclusion.

The multidimensional form, with the condition $\langle X^i,X^j\rangle_t = \delta_{ij}t$, is the one we shall actually need. It is what will allow us, in the proof of Girsanov's theorem, to recognise a process as a Brownian motion under a \emph{new} measure without recomputing any distributions --- one need only check that quadratic variation is preserved, which it is, since quadratic variation is a pathwise notion and equivalent measures agree on which paths occur.

\begin{definition}(Exponential Process)
    Let $\alpha \in \mathcal L (\Omega, L^2[0,T])$. We define its exponential process by
    \begin{equation}\label{exponential_process_definition}
        \mathcal E _\alpha(t) = e^{  \int_0^t \alpha(s) dW(s)-\frac 1 2 \int_0^t \alpha^2(s)ds}
    \end{equation}
    which is precisely the solution of $d\mathcal E_\alpha(t) =\alpha(t)\mathcal{E}_\alpha (t) dW(t)$, $\mathcal E_\alpha(0) = 1$ so that $\mathcal{E}_\alpha(t)$ is a local martingale.
\end{definition}

This exponential process is important because it leads to the construction of a new probability $\mathbb{Q}$ for which $\mathcal{E}_\alpha(t)$ is actually a martingale.

\begin{theorem}\label{CHUCHIEXPONENTIALESPERANZA1}
    Let $\alpha \in \mathcal L (\Omega, L^2[0,T])$. Then $\mathcal E_\alpha(t)$ is a martingale if, and only if, $\mathbb E[\mathcal E_\alpha (t)]= 1 $ for all $t\in [0,T].$
\end{theorem}
\textbf{\textit{Proof:}} (...) $\blacksquare$

The previous result is tautological and hard to check in practice, so we need a more practical condition. The following result is known as Novikov's condition.

\begin{theorem}(Novikov's Condition)
    Let $\alpha \in \mathcal L (\Omega, L^2[0,T])$. If $\mathbb E\left[ e^{\frac 1 2 \int_0^T \alpha^2(s) ds} \right] < \infty$, then $\mathcal E_\alpha(t)$ is a martingale.

\end{theorem}

This is a much more practical way to check whether $\mathcal E_\alpha(t)$ is a martingale, and therefore whether we can construct a new probability $\mathbb{Q}$ such that $d\mathbb{Q} = \mathcal E_\alpha(T) d\mathbb P$. For the following Theorem, we need to make use of two auxiliary lemmas.

\begin{lemma}
    Let $\gamma$ be a nonnegative integrable random variable, so that $d\mathbb{Q} = \gamma d \mathbb{P}$ defines a new probability measure (see \ref{eq:radonnikodym}). Then, for any $\sigma$-algebra $\mathcal{G} \subset \mathcal{F}$ and $\Theta$ such that $\mathbb E _ \mathbb Q [\Theta] < \infty$, it follows that 
    \begin{equation}
        \mathbb{E}_ \mathbb Q [\Theta \,|\, \mathcal G] = \frac{\mathbb{E}_ \mathbb P [\gamma\Theta  \,|\, \mathcal G] }{\mathbb{E}_ \mathbb P [\gamma \,|\, \mathcal G] }, \, \, \mathbb{Q}-\text{almost surely.}
    \end{equation}
\end{lemma}

\textbf{\textit{Proof:}} (...) $\blacksquare$

Then, if $\alpha \in \mathcal L (\Omega, L^2[0,T])$ satisfies the condition on Theorem \ref{CHUCHIEXPONENTIALESPERANZA1}, it follows that $\mathcal E_\alpha(t)$ is a $\mathbb P$-martingale. Let $\mathbb Q$ be defined by $d\mathbb Q = \mathcal E_\alpha(T)d\mathbb P$, that is, $\mathbb Q (A) = \int _A \mathcal E_\alpha (T) \, d\mathbb P \text{ for every  } A\in \mathcal F$. 


\begin{lemma}
Let $\alpha \in \mathcal L (\Omega, L^2[0,T])$.   Then a $\{\mathcal F_t \}$-adapted process $X_t$, $t\in[0,T]$ is a $\mathbb Q$-martingale if, and only if, $X(t) \mathcal E_\alpha(t)$  is a $\mathbb{P}$-martingale.
\end{lemma}
\textbf{\textit{Proof:}} (...) $\blacksquare$

Finally, the following key result holds from the previous lemmas.

\begin{theorem}(Girsanov's Theorem)\label{girsanovtheorem}
Let $\alpha \in \mathcal L (\Omega, L^2[0,T])$ such that $\mathcal E_\alpha$ is a martingale. Then the process 

\begin{equation}
    Z(t) = W(t) - \int_0^t \alpha(s) \, ds, \text{ }\text{ } t\in[0,T]
\end{equation}
is a Brownian motion with respect to the measure $d\mathbb Q  = \mathcal E_\alpha (T) d\mathbb P$.
\end{theorem}
\textbf{\textit{Proof:}} (...) $\blacksquare$

Girsanov's theorem is the formal statement of something we shall use constantly: \emph{changing the 
measure changes the drift and leaves the volatility alone}. The Brownian motion $Z$ under 
$\mathbb{Q}$ is the old $W$ with a drift subtracted, and nothing else has moved. The reason is 
already contained in Theorem \ref{thm:levycharacterisation}: to verify that $Z$ is a $\mathbb{Q}$-Brownian 
motion one need only check that it is a continuous local martingale with $\langle Z\rangle_t = t$, 
and quadratic variation is a pathwise limit, so it is unaffected by passing to an equivalent 
measure. Two equivalent measures disagree about how likely a path is, never about which paths 
exist. This observation is the mathematical content of the statement --- which we shall meet again 
in Section \ref{sec:blackscholes} --- that an option's price does not depend on the expected return 
of the underlying, while it depends critically on its volatility.

For the multi-factor models of Section \ref{sec:othermodels} we require the vector-valued form.

\begin{theorem}[Girsanov's Theorem, multidimensional]\label{thm:girsanovmulti}
Let $\mathbf{W} = (W^1,\dots,W^n)$ be an $n$-dimensional Brownian motion and let 
$\alpha = (\alpha^1,\dots,\alpha^n)$ be an adapted process with 
$\int_0^T \lVert\alpha(s)\rVert^2\,ds < \infty$ almost surely. Define
\begin{equation}
\mathcal{E}_\alpha(t) = \exp\left( \int_0^t \alpha(s)\transp\,d\mathbf{W}(s) - \frac{1}{2}\int_0^t \lVert \alpha(s)\rVert^2\,ds \right),
\end{equation}
and suppose Novikov's condition holds in the form 
$\mathbb{E}\bigl[\exp\bigl(\tfrac12\int_0^T \lVert\alpha(s)\rVert^2 ds\bigr)\bigr] < \infty$. Then 
$\mathcal{E}_\alpha$ is a martingale, $d\mathbb{Q} = \mathcal{E}_\alpha(T)\,d\mathbb{P}$ defines a 
probability measure equivalent to $\mathbb{P}$, and
\begin{equation}
\mathbf{Z}(t) = \mathbf{W}(t) - \int_0^t \alpha(s)\,ds, \qquad t \in [0,T],
\end{equation}
is an $n$-dimensional $\mathbb{Q}$-Brownian motion with independent components. \hfill $\blacksquare$
\end{theorem}

That the components of $\mathbf{Z}$ remain independent under $\mathbb{Q}$ deserves emphasis, since 
it is not obvious and is exactly what the multidimensional Lévy characterisation delivers: the 
cross-variations $\langle Z^i, Z^j\rangle_t = \delta_{ij}t$ are pathwise quantities and survive the 
change of measure untouched. Note also that this is a statement about \emph{independent} driving 
factors; when a model is specified with correlated Brownian motions, as the Heston model of 
Subsection \ref{subsec:stochvol} is, one applies the theorem after decorrelating, in the manner of 
Definition \ref{def:correlatedbm} below.

\subsubsection*{The Martingale Representation Theorem}

Girsanov's theorem tells us which measures are available. The following tells us which \emph{payoffs} 
are attainable, and it is the deepest result of this subsection. Together the two constitute the 
entire mathematical content of the fundamental theorems of asset pricing in Section 
\ref{sec:derivatives}.

\begin{theorem}[Itô's Martingale Representation Theorem]\label{thm:mrt}
Let $\mathbf{W}$ be an $n$-dimensional Brownian motion on $\PSpace$ and let 
$\FiltW$ be its augmented natural filtration. Let $M$ be a square-integrable 
$\mathbb{P}$-martingale with respect to that filtration. Then there exists a unique adapted process 
$\varphi = (\varphi^1,\dots,\varphi^n) \in L^2_{\mathrm{ad}}([0,T]\times\Omega)$ such that
\begin{equation}\label{eq:mrt}
M_t = M_0 + \int_0^t \varphi(s)\transp \, d\mathbf{W}(s), \qquad t \in [0,T],
\end{equation}
almost surely. If $M$ is only a local martingale, the same holds with $\varphi$ in the larger class 
$\mathcal{L}(\Omega,L^2[0,T])$. \hfill $\blacksquare$
\end{theorem}

The theorem should be read as a converse to Theorem \ref{theorem:martingale_property}. There we 
showed that every Itô integral is a martingale; here we learn that, provided the filtration is the 
one generated by $\mathbf{W}$ and nothing more, \emph{every} martingale is an Itô integral. In the 
language of Proposition \ref{def:conditional-expectation}, item 9, the stochastic integrals do not 
merely sit inside the space of martingales: they exhaust it.

Two aspects of the hypothesis deserve to be underlined, because both will be load-bearing.

The first is that the filtration must be generated by the Brownian motion itself. If 
$\Filt$ carries additional information --- another independent Brownian motion, a 
Poisson process, an independent random variable revealed at time $T$ --- the conclusion fails 
outright, and it fails for an obvious reason: a martingale driven by that extra randomness cannot be 
written as an integral against $\mathbf{W}$, because the integral would be measurable with respect 
to the wrong $\sigma$-algebra. The theorem is a statement about there being \emph{no more randomness 
in the model than the driving Brownian motion supplies}.

The second is that $\varphi$ is not constructed. The proof is an orthogonality argument: one shows 
that the stochastic integrals form a closed subspace of $L^2(\Omega,\mathcal{F}_T)$ whose orthogonal 
complement contains only constants, using the fact that exponential martingales 
$\mathcal{E}_\alpha(T)$ span a dense set. Existence is therefore obtained without any formula for 
$\varphi$, which is worth knowing in advance: the theorem will guarantee that a replicating strategy 
exists while saying nothing whatever about how to compute it. Recovering $\varphi$ in concrete cases 
is precisely the delta-hedging computation of Section \ref{sec:blackscholes}, and in models where no 
closed form is available it is a numerical problem.

We can already state the consequence for which we shall need it, though the vocabulary is only 
introduced in the next section.

\begin{remark}[Why this matters for pricing]\label{rem:mrtpricing}
Let $\Payoff_T$ be a square-integrable contingent claim and let $\mathbb{Q}$ be an equivalent 
martingale measure. The discounted price process 
$M_t := \mathbb{E}^{\mathbb{Q}}[\Payoff_T/S^0_T \mid \mathcal{F}_t]$ is by construction a 
$\mathbb{Q}$-martingale --- this is the tower property, nothing more. If the filtration is generated 
by the Brownian motion driving the traded assets, Theorem \ref{thm:mrt} applied under $\mathbb{Q}$ 
produces $\varphi$ with $M_t = M_0 + \int_0^t \varphi\transp\,d\mathbf{Z}$, and $\varphi$ is 
--- after rewriting the integral in terms of the discounted asset prices --- a self-financing 
strategy replicating $\Payoff_T$ exactly.

Every claim is thus attainable, so the market is \emph{complete}, and the price of the claim is not 
merely a plausible expectation but the unique cost of a portfolio that reproduces it in every state 
of the world. This is the substance behind the Second Fundamental Theorem 
(Theorem \ref{thm:secondFTAP}), and it is the reason the Black--Scholes model is complete 
(Section \ref{sec:blackscholes}) while the stochastic volatility models of Subsection 
\ref{subsec:stochvol} are not: in the latter, the filtration is generated by \emph{two} Brownian 
motions while only \emph{one} asset is traded, the hypothesis of Theorem \ref{thm:mrt} fails, and 
the volatility risk cannot be hedged away.
\end{remark}




\subsection{ Itô Diffusions}\label{section:itoDiff}

Recall that, for a Markov process, the \textit{transition probability} $\mathbb P_{s,x}(t,A) = \mathbb P (X_t \in A \, | \, X_s = x)$, for $A$ measurable and $a\leq s\leq t \leq b$ and $x\in \mathbb{R}^n$. The Chapman-Kolmogorov equation states that 
$$
\mathbb P _ {s,x}(t,A) = \int _ {\mathbb R ^n} \mathbb P _ {u,z}(t,A) \, \mathbb{P}_ {s,x}(u,dz).
$$

\begin{definition}(Diffusion Process)\label{def:diffusion}
    A $\mathbb R ^n$-valued process $X_t$, $t\in [0,T]$ is called a \textbf{diffusion process} if its transition probabilities $\{\mathbb P_ {s,x}(t, \cdot )\}$ satisfy the following three conditions for every $t\in [0,T]$, every $x\in\mathbb R ^n$ and every $\delta > 0$:

     \begin{equation*}
                 \text{i) }   \lim_ {\varepsilon \rightarrow 0^+ }\frac{1}{\varepsilon }  \int_ {\{\|y-x\|>\delta\}} \mathbb P_ {t,x}(t+\varepsilon,dy) = 0             \end{equation*} 

                  \begin{equation*}
                 \text{ii) }   \lim_ {\varepsilon \rightarrow 0^+ }\frac{1}{\varepsilon }  \int_ {\{\|y-x\|\leq\delta\}} (y-x) \, \mathbb P_ {t,x}(t+\varepsilon,dy) = \mu(t,x) \in \mathbb R ^n \text{ exists,}\end{equation*}

                        \begin{equation*}
                 \text{iii) }   \lim_ {\varepsilon \rightarrow 0^+ }\frac{1}{\varepsilon }  \int_ {\{\|y-x\|\leq\delta\}}  (y-x)(y-x)\transp  \, \mathbb P_ {t,x}(t+\varepsilon,dy) = \Sigma(t,x) \in \mathbb{R}^{n\times n} \text{ exists,}\end{equation*}

\noindent where the vectors in $\mathbb{R}^n$ are implicitly taken as columns, so that $(y-x)(y-x)\transp$ yields a symmetric positive semi-definite matrix. $\mu$ and $\Sigma$ are called the \textit{drift} and the \textit{diffusion} coefficient, respectively.

The restriction of the domain of integration to a ball around $x$ is essential and not a technicality: the integral of a probability measure over the whole of $\mathbb{R}^n$ is $1$, so dropping the constraint in i) would make the limit $+\infty$ rather than $0$. Read correctly, the three conditions say that in a short time $\varepsilon$ the process moves a distance of order $\sqrt{\varepsilon}$ with overwhelming probability, that the mean of that displacement is of order $\varepsilon$, and that its covariance is of order $\varepsilon$ as well. When conditions ii) and iii) hold, the truncation at level $\delta$ turns out not to affect the limits, so we may and shall write them informally as
\[
\mu(t,x) = \lim_{\varepsilon \to 0^+}\frac{1}{\varepsilon}\mathbb{E}\left[X_{t+\varepsilon}-x \mid X_t = x\right], \qquad
\Sigma(t,x) = \lim_{\varepsilon \to 0^+}\frac{1}{\varepsilon}\mathbb{E}\left[(X_{t+\varepsilon}-x)(X_{t+\varepsilon}-x)\transp \mid X_t = x\right].
\]

Note how condition i) forbids the process from having jumps: it is precisely the statement that a displacement of size $\delta$ in time $\varepsilon$ has probability $o(\varepsilon)$. Therefore the Poisson process, for which such a displacement has probability $\lambda\varepsilon + o(\varepsilon)$, is not a diffusion.
 
\end{definition}

\begin{example}
    Take a $\mathbb{R}^n$-valued Brownian motion $W_t$. Note that $W_t^2$ is a sum $n$ independent squares of normal variables, so that for any $c>0$, 
    $$
    \mathbb P (|W_{t+\varepsilon}-x|\geq c \, | \, W_t=x) = \mathbb{P}|W_{t+\varepsilon}-W_t| \geq c) \leq c^{-4}\mathbb{E}(W_\varepsilon^4) = c^{-4}n(n+2) \varepsilon^2
    $$
    so that the first condition in the previous definition is satisfied, $\Sigma(t,x)= \mathbb{I} _{n\times n}$ and $\mu(t,x) = 0. $
\end{example}

The following result finally connects the previous concepts and gives an unambiguous meaning to $\mu.$

\begin{theorem}
    Let $\sigma(t,x)$ and $\mu(t,x)$ be functions as in Theorem \ref{theorem:solutions:ndimensional_sde}. Assume they are also continuous on $[0,T] \times \mathbb R^n$. The solution to the SDE there defined is a diffusion process where $\mu(t,x)$ is identified and $\Sigma(t,x)= \sigma(t,x)\cdot\sigma(t,x)^t$.
\end{theorem}

The proof of this Theorem is beyond the scope of our purposes and can be found in \cite{Kuo2006}.


We have reached a critical point where diffusions, SDE and semigroups interesect. These seemingly unrelated mathematical objects give rise to a set of (parabolic) partial differential equations (\textit{PDEs}) that are crucial for computational purposes and pricing financial derivatives. In this context, powerful mathematical tools can be applied.

The first of the approaches, where we consider a semigroup of operators, is included for completeness purposes, showing how its infinitesimal generators may be used to construct diffusion processes. It is also included to illustrate how Functional Analysis arises in this context.
\vspace{0.3cm}


\textbf{\underline{1. Semigroup approach}}
\vspace{0.2cm}

In this section, let's stick to the \textbf{stationary case}, that is, when both $\mu$ and $\Sigma$ just depend on $x$. We will need to make use of some basic concepts of Functional Analysis such as what a \textit{Banach space} is. All the definitions and standard results can be found in \cite{Brezis2011}.


\begin{definition}
    A family of linear operators $\{T_t, t\geq 0 \}$ on a Banach space $\mathbb{B}$ (with norm $\lVert \cdot \rVert)$ is called a contraction semigroup if it satisfies the following three properties:

    \begin{itemize}
    \item[a)] \textbf{Semigroup property:} $T_0 = \text{I}_\mathbb B$ and $T_{t+s}=T_tT_s$, for all $t,s\geq 0$.
        \item[b)] \textbf{Strong continuity:} $\lim_{t\downarrow0}T_t(v) = v\, \text{ for every } v\in\mathbb{B}.$
        \item[c)] \textbf{Contraction:} $\lVert T_t(v) \rVert \leq \lVert v\rVert \, \forall t\geq 0 $ and $v\in\mathbb B$.
        
    \end{itemize}

    If $\{T_t, t\geq 0\}$ is a contraction semigroup, we define its \textbf{infinitesimal generator} $\mathcal A$ by
    \begin{equation}
        \mathcal A (v) = \lim_{t\downarrow 0} \frac 1 t \left( T_t (v) - v \right),
    \end{equation}
    whose domain consists in all vectors $v$ for which the limit exists (strongly) in $\mathbb{ B}.$
\end{definition}

\begin{example}
    We can take the space of continuous and bounded functions $C_b[0,\infty)$ with the supremum norm $\lVert \cdot \rVert _ \infty$ and the translation transformations $T_t(f)(x)=f(x+t)$, so that $\mathcal A f = f'$, for those differentiable $f$ such that $f'\in C_b [0,\infty)$. In this case, $\mathcal A$ is an unbounded operator.
\end{example}

Connecting now with the stationary diffusion process $X_t$ driven by $\mu$ and $ \Sigma$, let's assume that its transition probabilities $\mathbb P_t$ satisfy the following three properties:
\begin{itemize}
    \item[1.)] For any $t>0$ and $c>0$, $\lim_{<\lVert x\rVert\rightarrow \infty} \mathbb P_t (x,\{y \, : \, \lVert y \rVert < c\}) = 0.$
    \item[2.)] For any $c>0$, $\lim_{t\downarrow 0}\mathbb P_t (x, \{y\, : \, \lVert y-x \rVert \geq c \}) = 0$ uniformly on $x\in \mathbb R^n.$
    \item[3.)] For any bounded and continuous function $f$ and $t >0$, 
    \begin{equation}\label{pre_hilleYoshida}
        T_t(f)(x) = \int_ {\mathbb{ R}^n} f(y) \, \mathbb{ P}_ t (x,dy), \, x \in \mathbb{ R}^n\, \, \text{is continuous.}
    \end{equation}
\end{itemize}
\begin{theorem}
    In the previous conditions, the family of linear operators given in condition 3.) is a contraction semigroup on the Banach space $C_0(\mathbb{R}^n)$ (the space of continuous functions vanishing at infinity), whose infinitesimal generator $\mathcal A$ is given by
    \begin{equation}
        \mathcal A (f) (x)= \frac 1 2 \sum_{i,j= 1}^n \Sigma_{ij}(x) \frac{\partial ^2 f}{\partial x_i \partial x_j} + \sum_{i= 1}^n \mu_i(x) \frac{\partial f}{ \partial x_i}
    \end{equation}
\end{theorem}
\begin{example}
    In the case of Brownian motion, since $\Sigma=\text{I}_{n\times n} $ and $\mu = 0$, we have $\mathcal A (f) = \frac 1 2 \Delta f$, where $\Delta$ is the Laplacian operator.
\end{example}

We shall not develop the semigroup picture further, having obtained from it what we needed: the identification of the generator, and the confirmation that the $\tfrac12$ appearing in Itô's rule is the same $\tfrac12$ appearing in $\tfrac12\Delta$. It is worth recording, however, what the semigroup property actually encodes, because it is the first of the three approaches to fail in Part \ref{part:part2}. The identity $P_{t+s} = P_t \circ P_s$ is a restatement of the Markov property: the evolution over $t+s$ can be decomposed into evolution over $s$ followed by evolution over $t$, because at the intermediate instant the present state is a sufficient statistic for the future. When the driving noise has memory --- as fractional Brownian motion does --- this decomposition is unavailable, the family $(P_t)$ is no longer a semigroup, and the generator ceases to exist as a local operator. Neither the PDE approach nor the semigroup approach survives intact; only the pathwise, integral formulation does, which is a large part of why the theory of Part \ref{part:part2} is built the way it is.

\vspace{0.3cm}

% \textbf{\underline{2. Partial Differential Equations Approach}}
% \vspace{0.2cm}

% Let's informally obtain one of the main results involving the relation between stochastic processes and PDEs. Let $\mathbb{P}_{s,x}(t, \cdot)$ for $0\leq s \leq t \leq T$ and $x\in\mathbb R^n$ be the trasition probabilities of a diffusion process $X_t$ on $[0,T]$ driven by $\Sigma(t,x)$ and $\rho(t,x).$ Let also $F_{s,x}$ be the distribution function of the random variable $X_t$ such that $X_0=x$. The Chapman-Kolmogorov equation adopts the form

% \begin{equation}
%     F_{s,x} (t,y)= \int_ {\mathbb{R}^n} F_{u,z}(t,y)dF_{s,x}(u,z).
% \end{equation}

% Now, fixing $t_0$ and $y_0$ and taking $u = s+ \varepsilon$, we can approximate (all evaluated in $(t_0,y_0))$
% \begin{equation}
%     F_{s+\varepsilon,z} \approx  F_{s+\varepsilon,x} + \frac {\partial  F_{s+\varepsilon,x} }{\partial x}(z-x) + \frac{1}{2}\frac{\partial ^2  F_{s+\varepsilon,x} }{\partial x ^2} (z-x)^2
% \end{equation}

% what in the multidimensional case leads to the approximation 


% \begin{equation}
%      F_{s,x} \approx F_{s+\varepsilon,x}+\frac{\partial F_{s+\varepsilon,x} }{\partial x} \varepsilon \mu(s,x) + \frac{1}{2} \frac{\partial^2 F_{s+\varepsilon,x}}{\partial x ^2} \varepsilon \Sigma(s,x)
% \end{equation}

% what, after rearraging, dividing by $\varepsilon $ and computing the limit $\varepsilon \rightarrow 0$ yields
% \begin{equation}\label{Kolmogorov_backward}
%     -\frac{\partial F_{s,x}}{\partial s} = \mu \frac{\partial F_{s,x}}{\partial x} + \frac 1 2 \Sigma(s,x) \frac{\partial ^2 F_{s,x}}{\partial x^2},
% \end{equation}
% together with the terminal condition 
% \begin{equation}
%     \lim_{s\uparrow t_0} F_{s,x}(t_0,y_0) = \biggl\{ \begin{array}{cc}
%          1& \text{ if } x<y_0 , \\
%          0 & \text{ if } x>y_0.
%     \end{array}    
% \end{equation}

% Equation \ref{Kolmogorov_backward} is called \textbf{Kolmogorov backward equation}, because the variable $s$ moves backwards from $t_0.$ However, we aim for another formula that enables \textit{forward} calculations.

% We will stay in the one dimensional case. Assume the transition probabilities of a Markov process $X_t$ possess density functions, that is,

% \begin{equation}
%     P_{s,x}(t,dy) = p_{s,x}(t,y)\, dy,
% \end{equation}
% so that the Chapman-Kolmogorov equation thakes the form 
% \begin{equation}
%     p_{s,x}(t,y) = \int_ {\mathbb{R}}p_{u,z}(t,y)p_{s,x}(u,z)\, dz.
% \end{equation}
% By using this identity and some auxiliary test functions, it can be proven that

% \begin{equation}
%     \frac{\partial p_{s_0,x_0}(t,y)}{\partial t}  = -\frac{\partial (\mu(t,y) p_{s_0,x_0}(t,y))}{\partial y} +\frac 1 2 \frac{\partial^2 (\Sigma(t,y)p_{s_0,x_0}(t,y))}{\partial y^2} 
% \end{equation}
% what, together with 
% \begin{equation}
%     \lim_{t\downarrow s_0}p_{s_0,x_0}(t,y) = \delta_{x_0}(y)
% \end{equation}
% is known as either the \textbf{Kolmogorov forward equation} or the \textbf{Fokker-Plank equation.} This equation is much more useful in Finance context because it allows calculations \textit{forward} in time, that is, from known market variables.

% \begin{theorem}
%     Let $\Sigma$ and $\mu$ satisfy the linear and growth conditions in $x.$ Assume, additionally, that there exists $c>0$ such that $w^t\Sigma(t,y)w \geq c\lVert w \rVert ^2 $
% for all $w\in \mathbb R^n$. Then:

% \begin{itemize}
%     \item[a)] The Kolmogorov backward equation has a unique solution. Moreover, there exists a continuous Markov process $X_t$ whose transition probabilities satisfy $\mathbb P_{s,x}(t_0,dy_0) = dF_{s,x}(t_0,y_0).$
%     \item[b)] Assume additionally that $\partial  \mu/ \partial x$, $\partial \Sigma/ \partial x$ and $\partial \Sigma / \partial x^2$ satisfy the Lipschitz and linear growth conditions in $x.$ Then the Kolmogorov forward equation has a unique solution. $\blacksquare$
% \end{itemize}

% \end{theorem}


% In terms of the infinitesimal generator, we can summarize the previous discussion in the following table:

% \begin{center}
% \begin{tabular}{|c|c|}
% \hline
% \textbf{Forward Equation} & \textbf{Backward Equation} \\
% \hline
% $\displaystyle \frac{\partial \rho}{\partial t}(t,x) = \mathcal A^* \rho(t,x)$ & 
% $\displaystyle \frac{\partial u}{\partial t}(t,x) = \mathcal A u(t,x), \quad u(0,x) = f(x)$ \\
% \hline
% \end{tabular}
% \end{center}
% where $\mathcal A^*$ represents the adjoint of the generator $\mathcal A.$


% \vspace{0.3cm}

% \textbf{\underline{3. Stochastic Differential Equations Approach}}
% \vspace{0.2cm}

% In this case, take $\mu$, $\Sigma$ and $\sigma$ such that $\sigma \sigma^T = \Sigma$. Notice that $\sigma U$ is also a valid solution for any $U$ such that $UU^T = \textbf{I}_{n\times n}.$ Therefore, the election of $\sigma$ is unique up the multiplication by an orthogonal matrix, since $UW(t)$ is also a Brownian motion, so the transition probabilities of $X_t$ do not depend on the particular choice of $\sigma$ within its class over the group of orthogonal matrices. Therefore, we call this class of matrices the \textit{square root} of $\Sigma.$

% Now, suppose $\Sigma$ and $\mu$ satisfy the linear and growth conditions. It is clear that any square root of $\Sigma$ satisfies them. Moreover, any such $\sigma$ also satifies the Lipschitz conditions given that 

% \begin{equation}
%     \lVert\sigma(t,x)-\sigma(t,y)\rVert = \frac{\lVert\Sigma(t,x)-\Sigma(t,y)\rVert}{\lVert\sigma(t,x)\rVert+\lVert\sigma(t,y)\rVert}\leq \frac{\lVert\Sigma(t,x)-\Sigma(t,y)\rVert}{2\sqrt{c}}.
% \end{equation}

% Having this, we can directly state the main theorem corresponding to the SDE approach.

% \begin{theorem}
%     The unique solution to 
%     \begin{equation}
%         dX_t = \mu(t,X_t) \, dt + \sigma(t,X_t)\, dW_t, \text{ } \text{ } X_0 = \Theta
%     \end{equation}
%     is a diffusion process, whose transition distribution function $F_{s,x}(t,y)$ is the Kolmogorov backward equation. If the additional conditions on part $b)$ of the previous Theorem are satisfied, then $X_t$ has density functions $p_{s,x}(t,y)$, satisfying the corresponding Kolmogorov forward equation.
% \end{theorem}

% \begin{example}
%     We illustrate the previous discussion with an example. Consider the infinitesimal generator (...).
% \end{example}

\textbf{\underline{2. Partial Differential Equations Approach}}
\vspace{0.2cm}

Let's informally obtain one of the main results involving the relation between stochastic processes and PDEs. Let $\mathbb{P}_{s,x}(t, \cdot)$ for $0\leq s \leq t \leq T$ and $x\in\mathbb R^n$ be the transition probabilities of a diffusion process $X_t$ on $[0,T]$ driven by $\mu(t,x)$ and $\Sigma(t,x).$ Let also $F_{s,x}$ be the distribution function of the random variable $X_t$ such that $X_0=x$. The Chapman-Kolmogorov equation adopts the form

\begin{equation}
    F_{s,x} (t,y)= \int_ {\mathbb{R}^n} F_{u,z}(t,y)\, dF_{s,x}(u,z).
\end{equation}

Now, fixing $t$ and $y$ and taking $u = s+ \varepsilon$, we can approximate the integrand around $x$ (for $z$ close to $x$):
\begin{equation}
    F_{s+\varepsilon,z}(t,y) \approx  F_{s+\varepsilon,x}(t,y) + \nabla_x F_{s+\varepsilon,x}(t,y) \cdot (z-x) + \frac{1}{2} (z-x)^T H_x(F_{s+\varepsilon,x})(z-x)
\end{equation}
where $\nabla_x$ and $H_x$ denote the gradient and Hessian with respect to the spatial variable. Substituting this back into the integral and recalling that $\int (z-x) dF_{s,x}(s+\varepsilon, z) = \mathbb{E}[X_{s+\varepsilon}-X_s | X_s=x] \approx \mu(s,x)\varepsilon$, we obtain:

\begin{equation}
     F_{s,x} \approx F_{s+\varepsilon,x} + \nabla_x F_{s+\varepsilon,x} \cdot (\mu(s,x) \varepsilon) + \frac{1}{2} \text{Tr}\left( H_x(F_{s+\varepsilon,x}) \Sigma(s,x) \right) \varepsilon
\end{equation}

Rearranging, dividing by $\varepsilon$, and computing the limit $\varepsilon \rightarrow 0$ yields
\begin{equation}\label{Kolmogorov_backward}
    -\frac{\partial F_{s,x}}{\partial s} = \sum_{i=1}^n \mu_i \frac{\partial F_{s,x}}{\partial x_i} + \frac{1}{2} \sum_{i,j=1}^n \Sigma_{ij} \frac{\partial^2 F_{s,x}}{\partial x_i \partial x_j},
\end{equation}
together with the terminal condition 
\begin{equation}
    \lim_{s\uparrow t} F_{s,x}(t,y) = \mathbb{I}_{\{x \le y\}}.
\end{equation}

Equation (\ref{Kolmogorov_backward}) is called the \textbf{Kolmogorov backward equation}, because the variable $s$ is the starting time, effectively moving backwards from the terminal time $t.$ However, we often aim for a formula that enables \textit{forward} calculations. Directly linked to this equation, the Feynman-Kac formula states that the solution to the Kolmogorov backward equation can be represented as an expected value of a functional of the diffusion process $X_t$. This will be stated in the next section.



\vspace{0.2cm}
Notice that the Kolmogorov backward equation derived in Equation (\ref{Kolmogorov_backward}) is effectively a special case of the Feynman-Kac formula where the interest rate $r(t,x) = 0$ and the terminal condition is the indicator function $h(x) = \mathbb{I}_{\{x \le y\}}$. In the context of quantitative finance, the Feynman-Kac theorem is a cornerstone result: it establishes the fundamental bridge between the risk-neutral pricing of a derivative (computed as a discounted expected payoff under a martingale measure) and the deterministic PDE that the derivative's price must satisfy to preclude arbitrage.



We shall stay in the one-dimensional case for simplicity now. Assume the transition probabilities of a Markov process $X_t$ possess density functions, that is,

\begin{equation}
    P_{s,x}(t,dy) = p_{s,x}(t,y)\, dy,
\end{equation}
so that the Chapman-Kolmogorov equation takes the form 
\begin{equation}
    p_{s,x}(t,y) = \int_ {\mathbb{R}}p_{u,z}(t,y)p_{s,x}(u,z)\, dz.
\end{equation}
By using this identity and some auxiliary test functions (integration by parts), it can be proven that

\begin{equation}
    \frac{\partial p_{s_0,x_0}(t,y)}{\partial t}  = -\frac{\partial (\mu(t,y) p_{s_0,x_0}(t,y))}{\partial y} +\frac 1 2 \frac{\partial^2 (\Sigma(t,y)p_{s_0,x_0}(t,y))}{\partial y^2} \label{eq:fokkerplanck}
\end{equation}
which, together with the initial condition
\begin{equation}
    \lim_{t\downarrow s_0}p_{s_0,x_0}(t,y) = \delta_{x_0}(y)
\end{equation}
is known as either the \textbf{Kolmogorov forward equation} or the \textbf{Fokker-Planck equation.} This equation is much more useful in a financial context because it allows calculations \textit{forward} in time, i.e., evolving from known current market variables.

\begin{theorem}
    Let $\Sigma$ and $\mu$ satisfy the linear growth and Lipschitz conditions in $x.$ Assume, additionally, that there exists $c>0$ such that $w^T\Sigma(t,y)w \geq c\lVert w \rVert ^2 $
for all $w\in \mathbb R^n$ (uniform ellipticity). Then:

\begin{itemize}
    \item[a)] The Kolmogorov backward equation has a unique solution. Moreover, there exists a continuous Markov process $X_t$ whose transition probabilities satisfy $\mathbb P_{s,x}(t_0,dy_0) = dF_{s,x}(t_0,y_0).$
    \item[b)] Assume additionally that the derivatives of $\mu$ and $\Sigma$ up to second order satisfy the Lipschitz and linear growth conditions. Then the Kolmogorov forward equation has a unique solution.
\end{itemize}

\end{theorem}


In terms of the infinitesimal generator, we can summarize the previous discussion in the following table:

\begin{center}
\begin{tabular}{|c|c|}
\hline
\textbf{Forward Equation} & \textbf{Backward Equation} \\
\hline
$\displaystyle \partial \rho/\partial t = \mathcal A^* \rho$ & 
$\displaystyle -\partial u/\partial s = \mathcal A u$ \\
\hline
\end{tabular}
\end{center}
where $\mathcal A^*$ represents the adjoint of the generator $\mathcal A$; see \cite{karatzasshreve1991} for the functional-analytic details. The duality is worth reading off the table: the backward equation propagates a \emph{payoff} backwards from maturity and is the natural object for pricing a single contract, whereas the forward equation propagates a \emph{density} forwards from today and is the natural object for calibration, since it evolves from known current market data. Dupire's formula in Subsection \ref{subsec:localvol} is precisely an application of the second.


\vspace{0.3cm}

\textbf{\underline{3. Stochastic Differential Equations Approach}}
\vspace{0.2cm}

In this case, take $\mu$, $\Sigma$ and $\sigma$ such that $\sigma \sigma^T = \Sigma$. Notice that $\sigma U$ is also a valid solution for any matrix $U$ such that $UU^T = \textbf{I}_{n\times n}$ (orthogonal matrix). Therefore, the choice of $\sigma$ is unique up to multiplication by an orthogonal matrix. Since $U W(t)$ is also a Brownian motion, the transition probabilities of $X_t$ do not depend on the particular choice of $\sigma$ within its class. We call this class of matrices the \textit{square root} of $\Sigma.$

Now, suppose $\Sigma$ and $\mu$ satisfy the linear growth and Lipschitz conditions. It is clear that any square root of $\Sigma$ satisfies the growth condition. Moreover, if $\Sigma$ is uniformly elliptic (bounded away from zero), $\sigma$ also satisfies the Lipschitz conditions. In the 1D case, this can be seen via the inequality:

\begin{equation}
    |\sigma(t,x)-\sigma(t,y)| = \frac{|\Sigma(t,x)-\Sigma(t,y)|}{|\sigma(t,x)|+|\sigma(t,y)|}\leq \frac{|\Sigma(t,x)-\Sigma(t,y)|}{2\sqrt{c}}.
\end{equation}

Having this, we can directly state the main theorem corresponding to the SDE approach.

\begin{theorem}
    The unique solution to 
    \begin{equation}
        dX_t = \mu(t,X_t) \, dt + \sigma(t,X_t)\, dW_t, \quad X_{t_0} = \Theta
    \end{equation}
    is a diffusion process, whose transition distribution function $F_{s,x}(t,y)$ satisfies the Kolmogorov backward equation. If the additional conditions on part $b)$ of the previous Theorem are satisfied, then $X_t$ has density functions $p_{s,x}(t,y)$, satisfying the corresponding Kolmogorov forward equation.
\end{theorem}

\begin{example}
    We illustrate the previous discussion with an example. Consider the infinitesimal generator (...).
\end{example}



\subsection{References to Section 1.2}
The theoretical framework presented in this chapter draws primarily from two authoritative sources. The general Probability Theory concepts follow the classic exposition of P. Billingsley (\cite{Biggy}), while the construction of the Stochastic Integral is based on the comprehensive treatment by H. Kuo (\cite{Kuo2006}). Additionally, references for the underlying Functional Analysis theory, implicit in the PDE treatment of diffusions, can be found in the seminal work of H. Brezis (\cite{Brezis2011}).

Certain advanced topics, such as the general formulation of Stochastic Integration with respect to semimartingales or deeper pathwise properties of Brownian motion, have been omitted as they lie beyond the scope of this work. However, the selected contents provide a rigorous and sufficient foundation for the financial applications discussed in subsequent sections and chapters.





\section{Financial Derivatives}\label{sec:derivatives}

In this section, we introduce the main ideas behind Mathematical Finance and the role played by Mathematics in this field. We first informally introduce the concept of \textit{derivative} and give some examples of the main quoted derivatives. The rigorous mathematical treatment of hedging and market completeness is addressed afterwards.

\subsection{The Concept of a Derivative}

In the context of modern mathematical finance, a financial derivative is formally defined as a \textit{contingent claim}. Its value is derived from---or ``contingent'' upon---the state of some other underlying variables.

Let us fix a filtered probability space $(\Omega, \mathcal{F}, (\mathcal{F}_t)_{t \in [0, T]}, \mathbb{P})$, where $\mathbb{P}$ represents the physical (or real-world) probability measure, and the filtration $(\mathcal{F}_t)_{t \geq 0}$ models the flow of information over time up to a finite time horizon $T$. Mathematically, a financial derivative with maturity $T$ is an $\mathcal{F}_T$-measurable random variable, denoted by $H_T$. There are certain derivatives, such as perpetual bonds, which have an infinite time horizon ($T= \infty$), but these fall outside the scope of our finite-horizon models, as they lack standard liquidity and do not constitute a core part of the derivatives market.

This measurability condition implies that the value of the derivative at maturity is fully determined by the information available at time $T$. The function determining this value is often referred to as the \textit{payoff function}, denoted by $\Phi_H$. The key idea behind the subsequent theory is being able to assign \textit{fair prices} to those contingent claims, deduced from quoted (and liquid) market data.

In order to motivate the following definitions, let us start by thinking of a linear expression of the form $ax+by=c$, given $a,b,c\in\mathbb{R}$ and $ab\neq 0$. It is clear that, if we know the value of either $x$ or $y$ determines the value of the other. Something similar happens for market derivatives. Entangled by more sophisticated relationships, the prices of certain claims are not independent of others. Modern markets are composed of a cohort of derivatives, bought and sold by banks, hedge funds, and other institutions. However, the core of these derivatives are standardized contracts that are highly liquid (bought and sold in large volumes with high frequency), and the language institutions and contracts share is rooted in the mathematical machinery of Probability Theory and Stochastic Processes.

These standardized contracts provide enough information to infer prices of more custom contracts, ensuring that these new derivatives are priced consistently with the rest of the market data. The foundational principle behind this consistency is that markets must (or at least, theoretically tend to) lack \textbf{arbitrage.} An informal definition of arbitrage is the possibility of making a profit in the market without carrying any risk and with zero net initial investment. 

For example, consider a non-dividend-paying stock trading at $S_0$ and a forward contract on that stock with delivery price $K > S_0 e^{rT}$ (where $r$ is the continuous risk-free rate). An arbitrageur could borrow $S_0$ from the bank, buy the stock, and simultaneously short the forward contract. At maturity $T$, the arbitrageur delivers the stock, receives $K$, and repays the bank $S_0 e^{rT}$, locking in a strictly positive, risk-free profit of $K - S_0 e^{rT}$.

There are, of course, practical criticisms of such idealized examples. Factors such as transaction costs, bid-ask spreads, differential borrowing and lending rates, and taxes all create friction. However, since modern markets deal with highly fungible assets and increasingly efficient electronic trading, these arbitrage bounds are remarkably tight. Consequently, it became clear that Mathematics is essential here—not only to price assets consistently but to quantify underlying risks and engineer dynamic strategies to \textbf{hedge} them.

\subsubsection*{The Underlying Asset}
The stochastic process driving the derivative's value is known as the \textit{underlying asset}, denoted by $S = (S_t)_{0 \le t \le T}$. The underlying is defined as an $(\mathcal{F}_t)$-adapted stochastic process, and a derivative may even depend on several underlyings. In financial markets, $S_t$ typically represents the price of:
\begin{itemize}
    \item \textbf{Equities:} Individual stocks or market indices.
    \item \textbf{Fixed Income:} Bonds or interest rates (often modeled via a short-rate process $r_t$).
    \item \textbf{Commodities:} Physical goods such as oil, gold, or electricity.
    \item \textbf{Currencies (FX):} Exchange rates between two economies.
\end{itemize}

Derivatives are generally categorized by the functional form of their terminal payoff $H_T$.

\paragraph{1. Linear Derivatives (Forwards and Futures)}
These instruments possess a payoff that is a linear function of the underlying asset price at maturity. A prime example is the \textbf{Forward Contract}, which is an agreement to buy or sell the asset at time $T$ for a pre-specified delivery price $K$. The payoff at time $T$ for the holder of a long position is:
\begin{equation}
    H_T = S_T - K
\end{equation}
Note that this payoff ranges over $\mathbb{R}$, implying a potential liability for the holder if the asset price falls below $K$.

\paragraph{2. Non-Linear Derivatives (Options)}
Options provide the \textit{right}, but not the \textit{obligation}, to transact. This asymmetry introduces convexity (non-linearity) into the payoff structure, truncating downside risk.
\begin{itemize}
    \item \textbf{European Call Option:} The right to buy $S_T$ at strike price $K$. The payoff is defined as:
    \begin{equation}
        H_T = (S_T - K)^+ = \max(S_T - K, 0)
    \end{equation}
    \item \textbf{European Put Option:} The right to sell $S_T$ at strike price $K$. The payoff is defined as:
    \begin{equation}
        H_T = (K - S_T)^+ = \max(K - S_T, 0)
    \end{equation}
        \item \textbf{American Options:} American call or put options are similar to their European counterparts, with the key difference being that the option can be \textit{exercised} at any time up to maturity $T$.
\end{itemize}




From a theoretical standpoint, derivatives serve two primary functions in the market:
\begin{enumerate}
    \item \textbf{Hedging:} They allow agents to transfer specific risks (e.g., volatility or tail risk) to those willing to bear them.
    \item \textbf{Market Completeness:} In an idealized complete market, derivatives allow for the replication of any payoff profile, effectively spanning the state space of future uncertainties.
\end{enumerate}

These concepts will be rigorously defined in the next subsection. However, the main question remains: what price should these contracts hold? Before entering the mathematical machinery, we can establish bounds based on physical intuition, much like conservation laws:
\begin{itemize}
    \item If $T=0$, meaning we are pricing a claim that expires immediately (and only depends on the underlying at maturity, not on its trajectory), the price must equal the payoff. For a European Call Option with strike $K$ and underlying $S_0$, its price must be $(S_0-K)^+$.
    \item If an asset gives its buyer strictly more optionality compared to another derivative, its price must be greater than or equal to the latter. For example, the price of an American Call Option must be bounded below by the price of a European Call Option with same strike and maturity.
\end{itemize}



\subsection{Derivative Pricing}\label{subsec:derivativepricing}

We now turn to the rigorous problem of pricing contingent claims. While elementary treatments often assume asset prices follow geometric Brownian motions and interest rates are constant, these assumptions are insufficient for a general mathematical theory.

We work in a filtered probability space $(\Omega, \mathcal{F}, (\mathcal{F}_t)_{0 \le t \le T}, \mathbb{P})$ satisfying the usual conditions (right-continuity and completeness). The time horizon $T$ is finite.

Let the market consist of $n+1$ assets denoted by the vector process $S = (S^0, S^1, \dots, S^n)$. We assume that each component $S^i$ is an \textbf{Itô process}. While the most abstract fundamental theorems of asset pricing rely on the broader class of general semimartingales, Itô processes provide the precise mathematical framework required for the diffusion models prevalent in this thesis, seamlessly connecting with the stochastic calculus developed in Section 1.2.

We distinguish the $0$-th asset, $S^0$, as the \textbf{numeraire}.
\begin{definition}[Numeraire]
A numeraire is a price process $S^0$ that is strictly positive ($S^0_t > 0$ almost surely for all $t$) and non-dividend paying. It serves as the unit of account for the market.
\end{definition}
\noindent While $S^0$ is often chosen to be a bank account $B_t = \exp(\int_0^t r_s ds)$, the theory holds for \textit{any} valid Itô process numeraire (e.g., a foreign currency or a stock). The bank account is merely a convenient particular case because it explicitly models the time value of money. 

We define the \textbf{discounted price process} $X$ as the prices of the primary assets expressed in units of the numeraire:
\begin{equation}
    X_t = \left( 1, \frac{S^1_t}{S^0_t}, \dots, \frac{S^n_t}{S^0_t} \right)
\end{equation}
Because the denominator $S^0_t$ is strictly positive almost surely, the transformation $f(x, y) = x/y$ is smooth. By a direct application of the multidimensional Itô's formula, the discounted price processes $X_t^i$ are themselves Itô processes.

A trading strategy is a predictable process $H = (H^0, H^1, \dots, H^n)$, where $H^i_t$ represents the number of units of asset $i$ held at time $t$. The value of the portfolio at time $t$ is the scalar process:
\begin{equation}
    V_t(H) = \sum_{i=0}^n H^i_t S^i_t
\end{equation}

For the strategy to be economically meaningful, it must be \textbf{self-financing}. In the language of stochastic integration, this implies that changes in wealth are generated solely by the fluctuations of the asset prices, with no external injection or extraction of capital:
\begin{equation}
    dV_t(H) = \sum_{i=0}^n H^i_t \, dS^i_t \label{eq:selffinancing}
\end{equation}
When working with discounted prices $X$, the self-financing condition simplifies. Applying Itô's product rule, it can be shown that a strategy is self-financing if and only if the discounted wealth process $V_t/S^0_t$ satisfies:
\begin{equation}
    \frac{V_t(H)}{S^0_t} = \frac{V_0(H)}{S^0_0} + \int_0^t \sum_{i=1}^n H^i_u \, dX^i_u = \frac{V_0}{S^0_0} + (H \cdot X)_t
\end{equation}
where $(H \cdot X)$ denotes the vector stochastic integral.

In continuous time, the standard definition of arbitrage is insufficient because "doubling strategies" (integrands with extremely large negative values) can technically generate profit from nothing in local time intervals. To preclude this, we restrict ourselves to certain strategies.
\begin{definition}[Admissible Strategy]\label{def:admissible}
A self-financing strategy $H$ is called $a$-admissible if its discounted value process is uniformly bounded from below by a constant $-a$ almost surely (i.e., $(H \cdot X)_t \ge -a$ for all $t$). A strategy is admissible if it is $a$-admissible for some $a \ge 0$.
\end{definition}


The classical definition of "No Arbitrage" fails to be closed in the topology of $L^\infty$. Delbaen and Schachermayer (1994) proved that the appropriate condition for the existence of an equivalent martingale measure is \textit{No Free Lunch with Vanishing Risk}.

Let $K$ be the set of random variables representing the terminal wealth of all admissible strategies starting with zero capital:
\begin{equation}
    K = \{ (H \cdot X)_T \mid H \text{ is admissible and } H_0 = 0 \}
\end{equation}
Let $C$ be the cone of claims dominated by elements in $K$, i.e., $C = \{ g \in L^\infty(\mathcal{F}_T) \mid \exists f \in K, g \le f \}$.
The market  satisfies (at least tries to) \textbf{NFLVR} if:
\begin{equation}
    \overline{C} \cap L_+^\infty = \{0\}
\end{equation}
where $\overline{C}$ is the closure of $C$ in the $L^\infty$ norm. Intuitively, this condition prevents not just exact arbitrage, but also sequences of strategies that approximate an arbitrage opportunity with risk converging to zero.

\vspace{0.2cm}

We are now equipped to state the core result linking market topology to measure theory.

\begin{definition}[Equivalent Local Martingale Measure]
A probability measure $\mathbb{Q}$ is an Equivalent Local Martingale Measure (ELMM) for the numeraire $S^0$ if:
\begin{enumerate}
    \item $\mathbb{Q} \sim \mathbb{P}$ (they share the same null sets).
    \item The discounted price process $X_t = S_t / S^0_t$ is a \textbf{local martingale} under $\mathbb{Q}$.
\end{enumerate}
\end{definition}

\begin{theorem}[First Fundamental Theorem of Asset Pricing]\label{thm:firstFTAP}
Assume the process $X$ is locally bounded. The market satisfies the NFLVR condition if and only if there exists at least one Equivalent Local Martingale Measure $\mathbb{Q}$.
\end{theorem}



The other reasonable question to ask ourselves is whether every contingent claim can be replicated by a self-financing strategy. Before stating the second theorem, we must rigorously define market completeness.

\begin{definition}[Market Completeness]
A market is \textbf{complete} if every contingent claim $\xi \in L^\infty(\mathcal{F}_T)$ is \textit{attainable}. That is, for every bounded random variable $\xi$, there exists an admissible self-financing strategy $H$ and an initial capital $x \in \mathbb{R}$ such that:
\begin{equation}
    \frac{\xi}{S^0_T} = x + (H \cdot X)_T \quad \mathbb{P}\text{-almost surely.}
\end{equation}
Note that the statement is made under the physical measure $\mathbb{P}$. This is deliberate: at this stage $\mathbb{Q}$ has not been shown to be unique --- indeed its uniqueness is exactly what the next theorem characterises --- so quantifying over ``$\mathbb{Q}$-almost surely'' would be ambiguous. Since every equivalent local martingale measure shares its null sets with $\mathbb{P}$, the choice is harmless, but it is the honest one.
\end{definition}

\begin{theorem}[Second Fundamental Theorem of Asset Pricing]\label{thm:secondFTAP}
Assume that the set of equivalent local martingale measures $\mathcal{M}^e$ is non-empty (i.e., the market is arbitrage-free). The market is complete if and only if $\mathcal{M}^e$ is a singleton set (i.e., $\mathbb{Q}$ is unique).
\end{theorem}

If the market is arbitrage-free and complete (i.e., there exists a unique $\mathbb{Q}$), the fair price at time $t$ of a generic derivative paying $H_T$ at maturity is given by the expectation of the discounted payoff:
\begin{equation}
    \frac{V_t}{S^0_t} = \mathbb{E}^\mathbb{Q} \left[ \frac{H_T}{S^0_T} \bigg| \mathcal{F}_t \right] \implies V_t = S^0_t \, \mathbb{E}^\mathbb{Q} \left[ \frac{H_T}{S^0_T} \bigg| \mathcal{F}_t \right]
\end{equation}

This formulation reveals that "risk-neutral probabilities" are not absolute; they are  linked to the choice of numeraire.


Let $N$ be a different numeraire (strictly positive semimartingale). Let $\mathbb{Q}^S$ be the measure associated with numeraire $S^0$, and $\mathbb{Q}^N$ be the measure associated with numeraire $N$.
Prices must be consistent regardless of the unit of account. Therefore, for any asset $Z$:
\begin{equation}
    \frac{Z_t}{S^0_t} = \mathbb{E}^{\mathbb{Q}^S}\left[ \frac{Z_T}{S^0_T} \mid \mathcal{F}_t \right] \quad \text{and} \quad \frac{Z_t}{N_t} = \mathbb{E}^{\mathbb{Q}^N}\left[ \frac{Z_T}{N_T} \mid \mathcal{F}_t \right]
\end{equation}

The change of measure from $\mathbb{Q}^S$ to $\mathbb{Q}^N$ is defined by the Radon-Nikodym derivative, which turns out to be exactly the normalized ratio of the numeraires:
\begin{equation}
    \frac{d\mathbb{Q}^N}{d\mathbb{Q}^S}\bigg|_{\mathcal{F}_T} = \frac{N_T / S^0_T}{N_0 / S^0_0} \label{eq:numerairechange}
\end{equation}
This elegantly shows that changing the numeraire is mathematically equivalent to performing a change of probability measure. This technique is particularly powerful in interest rate theory (e.g., switching to the $T$-forward measure to efficiently price bond options and swaptions).

\vspace{0.2cm}

Finally, we can now state the Feynman-Kac formula, which provides a powerful link between the PDE approach and the probabilistic representation of derivative prices.

\begin{theorem}[Feynman--Kac Formula]\label{thm:feynmankac}
    Let $X_\tau$ be an $n$-dimensional Itô diffusion process on the interval $[t, T]$ driven by the stochastic differential equation
    \begin{equation}
        dX_\tau = \mu(\tau, X_\tau) d\tau + \sigma(\tau, X_\tau) dW_\tau, \quad \tau \in [t, T],
    \end{equation}
    with initial condition $X_t = x$, where $W_\tau$ is a standard $n$-dimensional Brownian motion, and the diffusion matrix is given by $\Sigma(\tau, x) = \sigma(\tau,x)\sigma(\tau,x)^\top$. Let $h: \mathbb{R}^n \rightarrow \mathbb{R}$ be a Borel-measurable terminal payoff function, and $r: [0,T] \times \mathbb{R}^n \rightarrow \mathbb{R}$ be a continuous function representing the continuously compounded interest rate.
    
   Assume $u(t,x)$ is a sufficiently regular function (typically $u \in C^{1,2}([0,T) \times \mathbb{R}^n) \cap C([0,T] \times \mathbb{R}^n)$) that satisfies the following linear parabolic partial differential equation:
    \begin{equation}\label{FK_PDE}
        \frac{\partial u}{\partial t}(t,x) + \sum_{i=1}^n \mu_i(t,x) \frac{\partial u}{\partial x_i}(t,x) + \frac{1}{2} \sum_{i,j=1}^n \Sigma_{ij}(t,x) \frac{\partial^2 u}{\partial x_i \partial x_j}(t,x) - r(t,x)u(t,x) = 0,
    \end{equation}
    for all $(t,x) \in [0,T) \times \mathbb{R}^n$, subject to the terminal condition:
    \begin{equation}
        u(T,x) = h(x).
    \end{equation}
    
    If $u(t,x)$ satisfies a suitable polynomial growth condition preventing the existence of explosive solutions, then $u(t,x)$ admits the stochastic representation:
    \begin{equation}\label{FK_Expectation}
        u(t,x) = \mathbb{E} \left[ \left. e^{-\int_t^T r(\tau, X_\tau) d\tau} h(X_T) \right| X_t = x \right].
    \end{equation}
\end{theorem}


\section{The Black-Scholes Framework}\label{sec:blackscholes}
\label{sec:black_scholes_framework}

Having established the general theory of arbitrage-free pricing and martingale measures, we now specialize our market model to the most famous and widely used framework in financial history: the Black-Scholes-Merton model.

\subsection{The Model}
\label{subsec:bs_model}

We assume a frictionless market: no transaction costs, no bid-ask spread, unrestricted short 
selling, and a constant risk-free rate $\rf$ at which one may borrow and lend freely. Under the 
\emph{physical} measure $\PP$ --- the measure describing what actually happens --- the underlying is 
modelled by a geometric Brownian motion,
\begin{equation} \label{eq:BS_SDE_P}
    d\Spot_t = \mu \Spot_t \, dt + \sigma \Spot_t \, d\Wt_t, \qquad \Spot_0 > 0,
\end{equation}
where $\mu \in \R$ is the expected rate of return, $\sigma>0$ the volatility, and $\Wt$ a 
$\PP$-Brownian motion. By Example \ref{example:blackscholessde} the solution is 
$\Spot_t = \Spot_0\exp\bigl((\mu-\tfrac12\sigma^2)t + \sigma \Wt_t\bigr)$, which is positive for all 
$t$ --- as a limited-liability equity should be.

It is essential to begin under $\PP$ and not under a risk-neutral measure. Doing so makes the 
central result of the theory visible rather than assumed, and that result is the following: 
\textbf{the price of the derivative does not depend on $\mu$.} Two investors who disagree completely 
about the expected return of the stock, but agree on its volatility, must agree on the price of 
every option written on it. This is not intuitive, it is not a modelling convention, and it is worth 
seeing emerge from the computation rather than being handed the risk-neutral dynamics with $\mu$ 
already replaced by $\rf$.

\subsubsection*{The hedging argument}

Consider a derivative whose value is a smooth function $V(t,\Spot)$ of time and the current spot. 
Form the portfolio
\[
\Pi_t = -V(t,\Spot_t) + \Delta_t \Spot_t,
\]
consisting of a short position in the derivative and $\Delta_t$ units of the underlying. Assuming 
the position is self-financing in the sense of \eqref{eq:selffinancing}, its value evolves as
\begin{equation}\label{eq:portfoliodynamics}
    d\Pi_t = -dV_t + \Delta_t \, d\Spot_t .
\end{equation}
Itô's rule (Theorem \ref{theorem:itorule}) applied to $V(t,\Spot_t)$, using 
$(d\Spot_t)^2 = \sigma^2\Spot_t^2\,dt$ from Table \ref{tab:ItoTable}, gives
\begin{equation}\label{eq:dV}
    dV_t = \left( \frac{\partial V}{\partial t} + \mu \Spot_t\frac{\partial V}{\partial \Spot} 
    + \frac{1}{2} \sigma^2 \Spot_t^2 \frac{\partial^2 V}{\partial \Spot^2} \right) dt 
    + \sigma \Spot_t \frac{\partial V}{\partial \Spot}\, d\Wt_t .
\end{equation}
Substituting \eqref{eq:BS_SDE_P} and \eqref{eq:dV} into \eqref{eq:portfoliodynamics} and collecting 
terms,
\begin{equation}\label{eq:portfoliogrouped}
    d\Pi_t = \underbrace{\left( -\frac{\partial V}{\partial t} 
    - \frac{1}{2}\sigma^2 \Spot_t^2 \frac{\partial^2 V}{\partial \Spot^2} 
    + \mu \Spot_t\Bigl[\Delta_t - \frac{\partial V}{\partial \Spot}\Bigr] \right)}_{\text{drift}} dt 
    + \underbrace{\sigma \Spot_t\left( \Delta_t - \frac{\partial V}{\partial \Spot} \right)}_{\text{diffusion}} d\Wt_t .
\end{equation}
The portfolio is exposed to randomness only through the second term. Choosing
\begin{equation}\label{eq:delta}
    \Delta_t = \frac{\partial V}{\partial \Spot}(t,\Spot_t)
\end{equation}
annihilates it. This is the \emph{delta} of the derivative, and \eqref{eq:delta} is the hedge.

Now observe what \eqref{eq:portfoliogrouped} has become. The same choice \eqref{eq:delta} that kills 
the diffusion term also kills the bracket multiplying $\mu$ in the drift. \textbf{The expected return 
of the stock has cancelled}, and it has done so for a structural reason rather than by accident: 
$\mu$ enters $d\Spot_t$ and $dV_t$ through the identical channel, so a position that is neutral to 
the random part of $\Spot$ is automatically neutral to its drift as well. What remains,
\[
d\Pi_t = \left( -\frac{\partial V}{\partial t} - \frac{1}{2}\sigma^2 \Spot_t^2\frac{\partial^2 V}{\partial \Spot^2}\right) dt,
\]
is riskless. A riskless self-financing portfolio must earn the risk-free rate --- otherwise, by 
borrowing or lending against it, one manufactures an arbitrage in the sense of Subsection 
\ref{subsec:derivativepricing} --- so $d\Pi_t = \rf \Pi_t\,dt$ with 
$\Pi_t = -V + \Spot_t \partial V/\partial \Spot$. Equating the two expressions and rearranging yields 
the Black--Scholes equation.

\begin{theorem}[Black--Scholes--Merton]\label{thm:blackscholespde}
Under the dynamics \eqref{eq:BS_SDE_P}, the value $V(t,\Spot)$ of a European derivative with payoff 
$V(T,\Spot) = \PayoffFn(\Spot)$ satisfies the linear parabolic partial differential equation
\begin{equation} \label{eq:BS_PDE}
    \frac{\partial V}{\partial t} + \rf \Spot \frac{\partial V}{\partial \Spot} 
    + \frac{1}{2} \sigma^2 \Spot^2 \frac{\partial^2 V}{\partial \Spot^2} - \rf V = 0, 
    \qquad (t,\Spot)\in[0,T)\times(0,\infty),
\end{equation}
together with the terminal condition $V(T,\Spot)=\PayoffFn(\Spot)$. \hfill $\blacksquare$
\end{theorem}

Equation \eqref{eq:BS_PDE} contains $\rf$ but not $\mu$, which is the claim made above, now proved. 
It is also exactly the Feynman--Kac equation of Theorem \ref{thm:feynmankac} for a diffusion with 
drift $\rf \Spot$ and diffusion coefficient $\sigma \Spot$, whence the probabilistic representation
\begin{equation}\label{eq:bsexpectation}
V(t,\Spot) = e^{-\rf(T-t)}\,\Ep{\QQ}\bigl[\PayoffFn(\Spot_T)\,\big|\,\Spot_t = \Spot\bigr],
\qquad d\Spot_t = \rf \Spot_t\,dt + \sigma \Spot_t\,d\Wt^{\QQ}_t .
\end{equation}
The passage from \eqref{eq:BS_SDE_P} to the dynamics in \eqref{eq:bsexpectation} is precisely 
Girsanov's theorem with $\alpha = (\mu-\rf)/\sigma$, the \emph{market price of risk}: the measure 
change subtracts a drift and leaves $\sigma$ untouched, exactly as discussed after Theorem 
\ref{girsanovtheorem}. That the hedging argument and the martingale argument give the same answer is 
not a coincidence but the content of Feynman--Kac, and the two viewpoints will alternate throughout 
what follows --- the PDE when we want to compute, the expectation when we want to reason.

Finally, the market is complete: by Theorem \ref{thm:mrt} and Remark \ref{rem:mrtpricing}, the 
filtration is generated by the single Brownian motion driving the single traded asset, so every 
square-integrable claim is attainable and $\QQ$ is unique. The delta \eqref{eq:delta} is the 
integrand $\varphi$ whose abstract existence the representation theorem guarantees; here we have it 
explicitly.

\subsubsection*{The closed-form solution}

For a European call, $\PayoffFn(\Spot)=(\Spot-\Strike)^+$, evaluating \eqref{eq:bsexpectation} 
against the log-normal density gives
\begin{equation}\label{eq:bscall}
    C_{BS}(t,\Spot) = \Spot \, N(d_+) - \Strike e^{-\rf(T-t)} N(d_-), 
    \qquad d_{\pm} = \frac{\ln(\Spot/\Strike) + (\rf \pm \tfrac12\sigma^2)\tau}{\sigma\sqrt{\tau}},
\end{equation}
with $\tau = T-t$ and $N$ the standard normal distribution function. The corresponding put price 
follows from put--call parity, Proposition \ref{prop:putcallparity} below.

One might instead assume \textbf{normal} (Bachelier) dynamics $d\Spot_t = \mu\,dt + \sigma\,d\Wt_t$, 
as in the original 1900 thesis. This permits negative prices, which is why it was criticised for 
equities, but it is standard for interest rates and for spread options where negative values are 
economically meaningful. For equities the log-normal convention prevails, on account of the limited 
liability of shareholders.

\subsubsection{Dividends, and why the forward is the right variable}\label{subsec:dividends}

The derivation above quietly assumed that holding the stock produces no cash flows. Real equities pay 
dividends, and the way they do so is awkward enough to force a change of variable that we shall keep 
for the rest of the dissertation.

Suppose first that dividends are paid \emph{continuously} at a constant proportional rate $\divy$. 
The holder of the stock receives $\divy \Spot_t\,dt$ over $[t,t+dt]$, so the self-financing condition 
\eqref{eq:portfoliodynamics} acquires the extra term $\Delta_t \divy \Spot_t\,dt$; repeating the 
argument verbatim replaces the coefficient $\rf$ of $\partial V/\partial \Spot$ in \eqref{eq:BS_PDE} 
by $\rf - \divy$, and $\rf$ by $\rf-\divy$ in $d_\pm$. Nothing structural changes. This is the case 
usually presented, and it is the case that does not occur.

Actual equity dividends are \emph{discrete cash amounts}, announced in advance and paid on known 
dates. If an amount $D_i$ is paid at time $t_i$, then no-arbitrage forces the spot to drop by that 
amount across the payment:
\begin{equation}\label{eq:dividendjump}
\Spot_{t_i} = \Spot_{t_i^-} - D_i .
\end{equation}
This is genuinely destructive of the framework above, and it is worth being precise about why. 
Equation \eqref{eq:dividendjump} says $\Spot$ has a \emph{predictable jump}: the process is not 
continuous, so Itô's rule as stated does not apply across $t_i$. Worse, the jump is \emph{additive} 
while geometric Brownian motion is multiplicative, so the two are structurally incompatible --- one 
cannot write $\Spot_t = \Spot_0 \mathcal{E}(\cdot)$ for a positive $\mathcal{E}$ and reproduce 
\eqref{eq:dividendjump}. Worse still, a fixed cash dividend can in principle exceed the spot price, 
so a model treating $D_i$ as deterministic does not even guarantee positivity. And since 
$\sigma$ multiplies $\Spot_t$ in \eqref{eq:BS_SDE_P}, a discontinuous drop in $\Spot$ produces a 
discontinuous drop in the absolute volatility, which is not what is observed.

The resolution is to stop modelling the spot and model the \emph{forward} instead.

\begin{definition}[Forward price]\label{def:forwardprice}
The \textbf{forward price} $\Fwd(t,T)$ is the delivery price agreed at $t$ for a transaction in the 
underlying at $T$ which has zero value at inception. Equivalently, writing $\Bond{t}{T}$ for the 
price at $t$ of a zero-coupon bond maturing at $T$,
\begin{equation}\label{eq:forwarddef}
\Fwd(t,T) = \frac{\Ep{\QT}\bigl[\Spot_T \mid \Ft\bigr] \cdot \Bond{t}{T}}{\Bond{t}{T}} = \Ep{\QT}\bigl[\Spot_T \mid \Ft\bigr],
\end{equation}
where $\QT$ is the $T$-forward measure associated with the numeraire $\Bond{\cdot}{T}$.
\end{definition}

Identity \eqref{eq:forwarddef} says that the forward price is a $\QT$-martingale --- it is a 
conditional expectation --- and this is the whole point. It holds by construction, with no assumption 
whatever on interest rates, on dividends, or on their nature. In particular we deliberately do 
\emph{not} write $\Fwd(t,T) = \Spot_t e^{(\rf-\divy)(T-t)}$, which is valid only for constant $\rf$, 
constant $\divy$, and purely proportional dividends. Instead we record the relation abstractly:
\begin{equation}\label{eq:spotforwardfactor}
\Spot_t = \Carry_t \, \Fwd(t,T),
\end{equation}
where $\Carry_t$ is a strictly positive, adapted, model-independent conversion factor assembled from 
the discount curve, the repo or borrow cost, and the announced cash and proportional dividend 
schedule. Its precise construction is a market-data question, not a modelling one, and it is 
addressed in Chapter \ref{ch:cap6}. What matters here is that \eqref{eq:spotforwardfactor} remains 
valid for time-dependent $\rf(t)$ and $\divy(t)$, for discrete cash dividends, for proportional 
dividends, and for any mixture of them.

Modelling the forward as a driftless log-normal martingale,
\begin{equation}\label{eq:blackdynamics}
d\Fwd_t = \sigma \Fwd_t \, d\Wt^{T}_t,
\end{equation}
therefore has three advantages simultaneously. The dividend discontinuities are absorbed into 
$\Carry_t$ and never touch the dynamics, which stay diffusive. The drift vanishes, because a forward 
is already a martingale under its own measure --- so there is no drift to specify and none to get 
wrong. And the resulting formula is calibrated directly against quantities the market quotes.

\subsubsection{The Forward Measure and Black-76}
Changing numeraire from the bank account to the zero-coupon bond $\Bond{\cdot}{T}$, in the manner of 
\eqref{eq:numerairechange}, converts \eqref{eq:bsexpectation} into an expectation of the undiscounted 
payoff under $\QT$, multiplied by a deterministic discount factor. With the dynamics 
\eqref{eq:blackdynamics} this gives Black's 1976 formula.

\begin{proposition}[Black-76 \cite{black1976}]\label{prop:black76}
Under \eqref{eq:blackdynamics}, the time-$t$ value of a European call with strike $\Strike$ and 
maturity $T$ is
\begin{equation}\label{eq:black76}
C(t) = \Bond{t}{T}\Bigl[\Fwd_t\,N(d_+) - \Strike\,N(d_-)\Bigr], \qquad 
d_\pm = \frac{\ln(\Fwd_t/\Strike) \pm \tfrac12 \sigma^2\tau}{\sigma\sqrt{\tau}}, \quad \tau = T-t.
\end{equation}
\hfill $\blacksquare$
\end{proposition}

Formula \eqref{eq:black76} is formally identical to \eqref{eq:bscall} with $\Spot$ replaced by 
$\Fwd_t$ and the interest rate removed from $d_\pm$ --- but note that the discount factor 
$\Bond{t}{T}$ remains, multiplying the whole bracket. It has not disappeared; it has merely been 
collected outside. We shall use \eqref{eq:black76} as the primitive throughout, treating 
\eqref{eq:bscall} as the special case obtained when $\Bond{t}{T} = e^{-\rf\tau}$ and 
$\Fwd_t = \Spot_t e^{\rf\tau}$.

\subsubsection*{Put--call parity}

Before turning to implied volatility we record the one relation in this subject that holds 
independently of any model, and which we shall invoke repeatedly.

\begin{proposition}[Put--call parity]\label{prop:putcallparity}
Let $C(t)$ and $P(t)$ be the time-$t$ prices of European call and put options on the same underlying, 
with the same strike $\Strike$ and maturity $T$. Then, in any arbitrage-free market,
\begin{equation}\label{eq:putcallparity}
C(t) - P(t) = \Bond{t}{T}\bigl(\Fwd(t,T) - \Strike\bigr).
\end{equation}
\end{proposition}

\textbf{\textit{Proof:}} The payoffs satisfy the pointwise identity 
$(\Spot_T-\Strike)^+ - (\Strike-\Spot_T)^+ = \Spot_T - \Strike$, valid for every value of $\Spot_T$. 
Taking $\Ep{\QT}[\,\cdot\mid\Ft]$ and multiplying by $\Bond{t}{T}$, the right-hand side becomes 
$\Bond{t}{T}(\Fwd(t,T)-\Strike)$ by \eqref{eq:forwarddef}. \hfill $\blacksquare$

The proof used no dynamics --- not \eqref{eq:blackdynamics}, not even continuity of paths --- only 
the linearity of expectation and an algebraic identity between payoffs. Consequently 
\eqref{eq:putcallparity} must hold in every model in this dissertation, and it is the standard 
sanity check on any implementation: a pricer that violates put--call parity is wrong, and no appeal 
to model choice can rescue it.

Two consequences are worth noting now. First, since $C-P$ is model-free, calls and puts of the same 
strike carry \emph{identical} information about volatility; a call and a put struck at $\Strike$ must 
therefore have the same implied volatility, and the smile is a single function of strike rather than 
two. Second, \eqref{eq:putcallparity} is how the forward is extracted from quoted option prices in 
practice: the call-minus-put spread is linear in $\Strike$ with slope $-\Bond{t}{T}$ and root 
$\Strike = \Fwd(t,T)$, so a regression across strikes recovers both the discount factor and the 
forward without any assumption about dividends. This is exactly the construction of $\Carry_t$ in 
\eqref{eq:spotforwardfactor}, obtained from market data rather than postulated.

\subsubsection{The Inversion Problem and Implied Volatility}
In professional markets, participants do not "price" options using Black-Scholes to find the dollar value. The market price $C^{mkt}$ is determined by supply and demand.

Instead, Black-Scholes is used as a translation mechanism. Since the pricing function $C_{BS}(\sigma)$ is strictly monotonic in $\sigma$, we can invert it to find the unique parameter $\ivol$ that matches the market price:
\begin{equation}
    C_{BS}(S, K, T, r, q, \ivol) = C^{mkt}
\end{equation}
This 
implied volatility is the true "price" quoted in the market. It represents the market's consensus on the average future volatility of the underlying asset over the life of the option.

\subsection{Smile and Skewness}
\label{subsec:smile_skew}

If the assumptions of Black-Scholes were satisfied physically (i.e., if returns were truly log-normal with constant variance), the implied volatility $\ivol$ would be constant across all strike prices $K$ and maturities $T$.

Empirical data reveals this is not the case. When we plot $\ivol(K)$ against $K$, we observe distinct shapes:
\begin{itemize}
    \item \textbf{The Skew (Smirk):} Typical in equity markets. $\ivol$ decreases as $K$ increases. This reflects "Crash Phobia"---investors bid up the price of low-strike Puts (insurance against market drops), implying higher volatility for downside strikes.
    \item \textbf{The Smile:} Typical in Forex markets. $\ivol$ is lowest at the money ($K \approx S$) and increases for both deep ITM and OTM options. This reflects "Fat Tails" (leptokurtosis) in the return distribution---extreme moves are more probable than the Gaussian model predicts.
\end{itemize}

The two shapes are not merely descriptive. Both say that the market's risk-neutral distribution 
is not log-normal, and the following classical result makes the connection exact.

\begin{proposition}[Breeden--Litzenberger \cite{breeden1978}]\label{prop:breedenlitzenberger}
Let $C(\Strike,T)$ denote the price of a European call of strike $\Strike$ and maturity $T$, and 
suppose the risk-neutral law of $\Spot_T$ admits a density $\psi_T$. Then
\begin{equation}\label{eq:breedenlitzenberger}
    \psi_T(\Strike) = \frac{1}{\Bond{0}{T}}\,\frac{\partial^2 C}{\partial \Strike^2}(\Strike,T).
\end{equation}
\end{proposition}

\textbf{\textit{Proof:}} Write $C(\Strike,T) = \Bond{0}{T}\int_{\Strike}^{\infty}(x-\Strike)\psi_T(x)\,dx$. 
Differentiating once under the integral sign, the boundary term vanishes because the integrand is 
zero at $x=\Strike$, leaving 
$\partial_\Strike C = -\Bond{0}{T}\int_\Strike^\infty \psi_T(x)\,dx$. Differentiating a second time 
gives \eqref{eq:breedenlitzenberger}. \hfill $\blacksquare$

The result is remarkable for what it does not assume: no model, no dynamics, only that prices are 
discounted expectations of the payoff. A continuum of option quotes at a fixed maturity therefore 
determines the entire marginal distribution of $\Spot_T$ --- the market does not merely express an 
opinion about volatility, it expresses a complete probability law, and the smile is the visible 
trace of that law failing to be log-normal.

Two consequences of \eqref{eq:breedenlitzenberger} will be used constantly. First, since 
$\psi_T \geq 0$, arbitrage-free call prices must be \emph{convex} in strike; a violation is a 
butterfly arbitrage, executable by buying two wing options and selling two at the middle strike. 
Second, the formula involves a \emph{second} derivative of market data, which is numerically 
delicate in a way that Subsection \ref{subsec:localvol} will have to confront directly.

\figurehere{Empirical implied volatility smile for a liquid index at two maturities, showing the 
flattening of the skew with increasing $T$. To be produced from the market data of Chapter 
\ref{ch:cap6}.}

\section{Other Important Models}\label{sec:othermodels}
\label{sec:other_models}

To account for the volatility smile and skew, we must extend the Black-Scholes framework. There are two primary approaches: making volatility a deterministic function of the state (Local Volatility) or making it a stochastic process (Stochastic Volatility).

\subsection{Local Volatility Models}\label{subsec:localvol}
\label{subsec:local_vol}

The Local Volatility (LV) framework, introduced by Dupire (1994), posits that volatility is a deterministic function of the current asset price and time: $\sigma = \sigma(S_t, t)$. The dynamics under the risk-neutral measure are:
\begin{equation}
    dS_t = r S_t dt + \sigma(S_t, t) S_t dW_t
\end{equation}
This model preserves the completeness of the market (there is only one source of randomness $W_t$) while allowing the model to fit \textit{any} observed arbitrage-free implied volatility surface exactly.

\subsubsection{The Dupire Formula and Gatheral's Transformation}

Dupire's insight was to link the unknown function $\lvol(\Strike,T)$ directly to observed market 
option prices. Applying the Fokker--Planck equation \eqref{eq:fokkerplanck} to the transition density 
of the diffusion and combining it with Breeden--Litzenberger yields the following.

\begin{theorem}[Dupire]\label{thm:dupire}
Suppose call prices $C(\Strike,T)$ are known for a continuum of strikes and maturities and are 
consistent with a diffusion of the form $d\Spot_t = (\rf-\divy)\Spot_t\,dt + \lvol(\Spot_t,t)\Spot_t\,d\Wt_t$. 
Then the local volatility function is uniquely determined by
\begin{equation} \label{eq:dupire}
    \lvol^2(\Strike, T) = \frac{\dfrac{\partial C}{\partial T} + (\rf-\divy)\,\Strike \dfrac{\partial C}{\partial \Strike} + \divy\,C}{\tfrac{1}{2} \Strike^2 \dfrac{\partial^2 C}{\partial \Strike^2}} .
\end{equation}
\hfill $\blacksquare$
\end{theorem}

The content of \eqref{eq:dupire} is striking and deserves to be stated plainly: \emph{there is 
exactly one} local volatility diffusion consistent with a given arbitrage-free surface. The 
calibration problem has a unique solution, not a family of them, and the model can therefore 
reproduce every European quote simultaneously and exactly. Note also that the denominator is 
proportional to the risk-neutral density by Proposition \ref{prop:breedenlitzenberger}, so 
\eqref{eq:dupire} is well posed precisely when that density is strictly positive.

\subsubsection*{Static arbitrage and the implied-variance parameterisation}

Formula \eqref{eq:dupire} is, however, close to unusable as written. It requires one derivative in 
maturity and two in strike of raw market data, and market data are quoted on a sparse grid, carry 
bid-ask spreads, and are contaminated by microstructure noise. Numerical differentiation amplifies 
all three, and the practical symptom is a negative denominator --- that is, a negative local 
variance, which is meaningless.

The remedy is not better differentiation but a change of variables, together with an interpolation 
that is guaranteed arbitrage-free by construction. Define the \textbf{log-moneyness} and the 
\textbf{total implied variance}
\begin{equation}\label{eq:totalvariance}
\lmon := \log\frac{\Strike}{\Fwd(t,T)}, \qquad \tvar(\lmon,T) := \ivol^2(\lmon,T)\,T .
\end{equation}
Total variance is the natural coordinate because it, rather than implied volatility itself, is what 
accumulates linearly in time under Black--Scholes. We now state the conditions under which a surface 
$\tvar$ is free of \emph{static} arbitrage --- that is, admits no arbitrage exploitable by a 
buy-and-hold portfolio of European options, with no trading in between.

\begin{proposition}[Absence of static arbitrage]\label{prop:staticarbitrage}
A total implied variance surface $\tvar(\lmon,T)$ is free of static arbitrage if the following hold.
\begin{itemize}
    \item[i)] (\textbf{Calendar spread}) For each fixed $\lmon$, the map $T \mapsto \tvar(\lmon,T)$ 
    is non-decreasing:
    \begin{equation}\label{eq:calendarspread}
    \frac{\partial \tvar}{\partial T}(\lmon,T) \geq 0 .
    \end{equation}
    \item[ii)] (\textbf{Butterfly}) For each fixed $T$, the function
    \begin{equation}\label{eq:butterfly}
    g(\lmon) := \left(1 - \frac{\lmon\,\partial_\lmon \tvar}{2\tvar}\right)^{\!2} 
    - \frac{(\partial_\lmon \tvar)^2}{4}\left(\frac{1}{\tvar} + \frac{1}{4}\right) 
    + \frac{\partial^2_{\lmon\lmon} \tvar}{2}
    \end{equation}
    satisfies $g(\lmon) \geq 0$ for all $\lmon$.
    \item[iii)] (\textbf{Wings}) $\tvar(\lmon,T)$ grows at most linearly in $|\lmon|$ as 
    $|\lmon|\to\infty$; more precisely, the asymptotic slopes are bounded by $2$, which is Lee's 
    moment formula \cite{lee2004}.
\end{itemize}
\hfill $\blacksquare$
\end{proposition}

Each condition has a direct financial meaning. Condition i) says a longer-dated option cannot be 
cheaper than a shorter-dated one of the same moneyness --- otherwise one buys the long and sells the 
short. Condition ii) is convexity of the call price in strike, expressed in the new coordinates, and 
its violation is a butterfly spread with negative cost and non-negative payoff. Condition iii) 
prevents the extrapolated wings from implying infinite moments of $\Spot_T$.

In practice one first fits a parametric form guaranteeing i)--iii) --- Gatheral's SVI 
parameterisation \cite{gatheral2014} being the standard choice, with SABR common in rates markets --- and only then 
differentiates. Since the fitted surface is smooth and arbitrage-free by construction, the 
differentiation is stable and the local variance is automatically non-negative. This yields 
Gatheral's form of the Dupire formula, which is the version actually implemented.

\begin{theorem}[Dupire's formula in implied-variance coordinates]\label{thm:gatheral}
With $\lmon$ and $\tvar$ as in \eqref{eq:totalvariance},
\begin{equation}\label{eq:gatheral}
\lvol^2(\lmon,T) 
= \frac{\dfrac{\partial \tvar}{\partial T}}
{\;1 - \dfrac{\lmon}{\tvar}\dfrac{\partial \tvar}{\partial \lmon} 
+ \dfrac{1}{4}\left(-\dfrac{1}{4} - \dfrac{1}{\tvar} + \dfrac{\lmon^2}{\tvar^2}\right)\!\left(\dfrac{\partial \tvar}{\partial \lmon}\right)^{\!2} 
+ \dfrac{1}{2}\dfrac{\partial^2 \tvar}{\partial \lmon^2}\;} .
\end{equation}
\hfill $\blacksquare$
\end{theorem}

Formula \eqref{eq:gatheral} is worth a closer look, because the two no-arbitrage conditions of 
Proposition \ref{prop:staticarbitrage} are visible in it. The numerator is exactly the calendar 
spread quantity \eqref{eq:calendarspread}, and the denominator is precisely the butterfly quantity 
$g(\lmon)$ of \eqref{eq:butterfly}. Hence
\[
\lvol^2 \geq 0 \iff \frac{\partial \tvar}{\partial T} \geq 0 \ \text{ and } \ g(\lmon) \geq 0,
\]
so that \textbf{the local variance is non-negative if and only if the implied volatility surface is 
free of static arbitrage}. The two statements are the same statement. This is the cleanest 
justification for the two-step procedure above: arbitrage-freeness is not an aesthetic requirement 
imposed on the interpolation, it is exactly the condition under which \eqref{eq:gatheral} returns a 
usable answer.

\subsubsection*{The limitation}

Local volatility fits the surface perfectly, and by Theorem \ref{thm:dupire} it is the unique 
diffusion that does. It also preserves market completeness, since a single Brownian motion drives a 
single traded asset and Theorem \ref{thm:mrt} applies unchanged. These are substantial virtues.

The defect is equally fundamental, and it concerns \emph{dynamics} rather than fit. In a local 
volatility model the instantaneous volatility is a deterministic function of the spot, so the future 
smile is entirely determined by where the spot goes. The consequence, observed universally in 
practice, is that the model predicts the smile flattening as the spot moves, whereas real smiles 
translate with the spot and retain their shape. A model may therefore reproduce every European price 
today and still be wrong about every European price tomorrow. This matters not at all for pricing 
European options --- their value depends only on the terminal marginal, which is fitted exactly --- 
and a great deal for anything path-dependent: barriers, cliquets, forward-starting options and 
autocallables all depend on the joint law across dates, which local volatility gets wrong. That 
observation is what motivates the next subsection.

\subsection{Stochastic Volatility Models}\label{subsec:stochvol}
\label{subsec:stochastic_vol}

To capture realistic dynamics of the volatility surface itself, we introduce a second source of randomness. In Stochastic Volatility (SV) models, the variance $v_t = \sigma_t^2$ follows its own Stochastic Differential Equation. Under the risk-neutral measure $\mathbb{Q}$, the system is described by:
\begin{equation}
    \begin{cases}
        dS_t = (r-q) S_t dt + \sqrt{v_t} S_t dW_t^S \\
        dv_t = \alpha(v_t, t) dt + \beta(v_t, t) dW_t^v
    \end{cases}
\end{equation}
A critical parameter is the correlation $\rho$ between the two Brownian motions: $d\langle W^S, W^v \rangle_t = \rho dt$.
\begin{itemize}
    \item If $\rho < 0$, a drop in price ($S \downarrow$) is associated with an increase in volatility ($v \uparrow$), naturally generating the \textbf{skew} observed in equity markets.
    \item The "vol of vol" parameter (in the diffusion of $v_t$) controls the convexity or \textbf{curvature} of the smile.
\end{itemize}

\subsubsection{The Heston Model}

The most prominent SV model is that of Heston (1993), in which the variance follows the 
Cox--Ingersoll--Ross square-root process already analysed in Example \ref{example:cirwellposed}:
\begin{equation}\label{eq:heston}
    \begin{cases}
        d\Spot_t = (\rf-\divy)\,\Spot_t\,dt + \sqrt{v_t}\,\Spot_t\,d\Wt_t^{\Spot},\\[2pt]
        dv_t = \kappa (\theta - v_t)\,dt + \xi \sqrt{v_t}\,d\Wt_t^{v},
    \end{cases}
    \qquad d\bigl\langle \Wt^{\Spot}, \Wt^{v}\bigr\rangle_t = \rho\,dt,
\end{equation}
where $\kappa$ is the speed of mean reversion, $\theta$ the long-run variance, $\xi$ the volatility 
of volatility, and $\rho$ the correlation, understood in the sense of Definition 
\ref{def:correlatedbm}. We recall from Subsection \ref{subsec:sdes} that \eqref{eq:heston} is well 
posed by Yamada--Watanabe rather than by the Lipschitz theory, and that strict positivity of $v$ 
requires the Feller condition $2\kappa\theta \geq \xi^2$ of Proposition \ref{prop:feller} --- a 
condition frequently violated by calibrated parameters, with the numerical consequences noted there.

Unlike local volatility, stochastic volatility renders the market \textbf{incomplete}, and we are 
now in a position to say exactly why rather than merely asserting it. The filtration is generated by 
two Brownian motions while only one asset is traded, so the hypothesis of Theorem \ref{thm:mrt} 
fails: a claim whose value depends on $\Wt^v$ cannot be written as a stochastic integral against the 
traded asset alone. By Theorem \ref{thm:secondFTAP} the equivalent martingale measure is therefore 
not unique, and pricing requires either specifying a market price of volatility risk --- already 
embedded in the $\QQ$-dynamics \eqref{eq:heston} --- or, as is universal in practice, calibrating 
$(\kappa,\theta,\xi,\rho,v_0)$ directly to liquid quotes. Volatility risk cannot be hedged away with 
the underlying; it can only be hedged with other options.

The reason Heston's model, rather than some other stochastic volatility model, became the industry 
standard is computational, and worth stating explicitly because Part \ref{part:part2} will build 
directly on it. The characteristic function of the log-price is available in closed form.

\begin{theorem}[Heston's characteristic function]\label{thm:hestoncharfn}
Let $X_t = \log \Spot_t$ under \eqref{eq:heston} and write $\tau = T-t$. Then
\begin{equation}\label{eq:hestoncf}
\charf(u,\tau) := \Ep{\QQ}\bigl[e^{iuX_T}\,\big|\,\Ft\bigr] 
= \exp\Bigl( iu\bigl(X_t + (\rf-\divy)\tau\bigr) + C(u,\tau)\,\theta + D(u,\tau)\,v_t \Bigr),
\end{equation}
where $C$ and $D$ solve the system of Riccati ordinary differential equations
\begin{equation}\label{eq:hestonriccati}
\frac{\partial D}{\partial \tau} = \frac{\xi^2}{2}D^2 + (\rho\xi i u - \kappa)D - \frac{u^2+iu}{2}, 
\qquad
\frac{\partial C}{\partial \tau} = \kappa D,
\end{equation}
with $C(u,0)=D(u,0)=0$. The system \eqref{eq:hestonriccati} admits an explicit solution in terms of 
elementary functions. \hfill $\blacksquare$
\end{theorem}

Given \eqref{eq:hestoncf}, European option prices follow by Fourier inversion: the call price is 
recovered from the characteristic function through an integral of Lewis--Carr--Madan type \cite{carrmadan1999},
\begin{equation}\label{eq:fourierinversion}
C(\Strike,T) = \Bond{t}{T}\left[ \Fwd_t - \frac{\sqrt{\Fwd_t \Strike}}{\pi}\int_0^{\infty} 
\mathrm{Re}\left( e^{iu\lmon}\,\charf\bigl(u - \tfrac{i}{2},\tau\bigr)\right)\frac{du}{u^2 + \tfrac14} \right],
\end{equation}
a one-dimensional integral evaluated by quadrature in microseconds. Calibration --- which requires 
pricing hundreds of options at each step of an optimiser --- is thereby feasible, whereas a model 
requiring Monte Carlo for each European quote would not be.

That the tractability of the model reduces to a system of Riccati equations is a point worth 
holding onto. The affine structure of \eqref{eq:heston} --- drift and variance both affine functions 
of the state --- is exactly what produces \eqref{eq:hestonriccati}, and it is a structural property 
rather than a happy accident. When we come to rough volatility in Part \ref{part:part2}, the model 
will retain this affine character but lose the Markov property, and \eqref{eq:hestonriccati} will be 
replaced by a \emph{fractional} Riccati equation. The pricing machinery \eqref{eq:fourierinversion} 
survives essentially unchanged; only the equation for $D$ becomes non-local in time. Recording the 
classical case here is what will make that generalisation a short step rather than a fresh start.

\subsubsection*{What stochastic volatility does and does not fix}

Stochastic volatility repairs the principal defect of local volatility: because the volatility has 
its own noise, the smile moves with the spot instead of flattening, and forward-starting products 
are priced far more sensibly. The correlation $\rho$ generates skew and the vol-of-vol $\xi$ 
generates curvature, both dynamically rather than by construction.

It does not, however, fit the surface exactly --- five parameters cannot match hundreds of quotes --- 
and it fails in one specific and instructive place. Empirically, the implied volatility skew of index 
options behaves like $T^{-H'}$ for small $T$ with $H'$ well below $\tfrac12$, becoming very steep at 
short maturities. Any model driven by standard Brownian motion, Heston included, produces a 
short-maturity skew that flattens as $T\to0$, because over short horizons a diffusive variance has 
had no time to move. No choice of $(\kappa,\theta,\xi,\rho,v_0)$ repairs this: the failure is not one 
of calibration but of the driving noise. This is the empirical observation from which Part 
\ref{part:part2} begins.

\subsection{Path-Dependence and American Options}\label{subsec:american}
\label{subsec:american_options}

We conclude this section by briefly discussing the challenges posed by path-dependent derivatives, focusing on the most liquid and prominent example: American Options. 

Because American options hold the right to be exercised at any time up to maturity, merely knowing a single implied volatility is insufficient to price them; a single volatility parameter is compatible with an infinite set of potential asset dynamics. This distinction is critically important because options on \textit{single names} (i.e., individual corporate stocks) are almost exclusively American-style. Investors in single equities are reluctant to relinquish the flexibility of early exercise, particularly to capture discrete dividends or navigate corporate actions. In contrast, broad market indices are generally considered more stable and are traded via extremely liquid European options.

We aim to show how a Local Volatility model can be adapted to price American options. However, we must first note a fundamental limitation: Local Volatility models do not perform perfectly for all types of path-dependent derivatives. In particular, because LV models assume volatility is a deterministic function perfectly correlated with the spot price, they often misprice derivatives with discontinuous payoffs (which often lead to ill-posed pricing problems) or derivatives whose strikes are set at a future date, such as forward-starting options and cliquets. 

Nevertheless, to price American options accurately we must model the entire volatility surface prior 
to maturity, not merely the terminal marginal. The model we use, expressed under the $T$-forward 
measure of Subsection \ref{subsec:dividends}, is
\begin{equation}\label{eq:localvolforward}
    d\Fwd_t = \lvol(\Fwd_t, t)\,\Fwd_t \, d\Wt^{T}_t .
\end{equation}
Working in the forward has the same three advantages set out in Subsection \ref{subsec:dividends}: 
the drift vanishes because a forward is a martingale under its own measure, the dynamics stay 
purely diffusive, and --- crucially here --- the discontinuities introduced by discrete cash 
dividends are absorbed into the spot-to-forward factor $\Carry_t$ of \eqref{eq:spotforwardfactor} 
and never enter \eqref{eq:localvolforward} at all.

\subsubsection*{Optimal stopping and the Snell envelope}

Before writing any partial differential equation, let us state what an American option \emph{is} in 
probabilistic terms, since the PDE formulation is derived from it and is opaque without it.

The holder may exercise at any stopping time of her choosing, and will choose the one maximising the 
value of the resulting cash flow. Writing $\Stop{[t,T]}$ for the set of stopping times valued in 
$[t,T]$ and $\PayoffFn$ for the exercise payoff, the value at time $t$ is
\begin{equation}\label{eq:optimalstopping}
V_t = \operatorname*{ess\,sup}_{\tau \in \Stop{[t,T]}} 
\Ep{\QQ}\left[ e^{-\int_t^{\tau}\rf(s)\,ds}\,\PayoffFn(\tau, \Spot_\tau) \,\Big|\, \Ft \right].
\end{equation}
The essential supremum is required because the family is uncountable; this is an optimal stopping 
problem, and its solution is classical.

\begin{theorem}[Snell envelope]\label{thm:snell}
Let $Y$ be a right-continuous adapted process of class (D) representing the discounted payoff. Then 
the value process $\Snell$ defined by \eqref{eq:optimalstopping} is the \textbf{Snell envelope} of 
$Y$: that is, $\Snell$ is the smallest càdlàg supermartingale dominating $Y$. It satisfies:
\begin{itemize}
    \item[i)] $\Snell_t \geq Y_t$ for all $t$ (domination);
    \item[ii)] $\Snell$ is a $\QQ$-supermartingale;
    \item[iii)] $\Snell$ is a $\QQ$-martingale on the \emph{continuation region} 
    $\mathcal{C} = \{(t,\omega) : \Snell_t > Y_t\}$;
\end{itemize}
and the stopping time $\tau^{*} = \inf\{ s \geq t : \Snell_s = Y_s \}$ is optimal in 
\eqref{eq:optimalstopping}. A full treatment is given in \cite{peskirshiryaev2006}. \hfill $\blacksquare$
\end{theorem}

The three properties have an immediate financial reading, and it is worth pausing on it because the 
LCP below is nothing more than their restatement. Property i) says the option is never worth less 
than exercising it --- otherwise one would buy the option, exercise immediately, and pocket the 
difference. Property ii) says that holding the option is, on average and after discounting, a losing 
proposition: waiting has a cost, which is what makes early exercise ever attractive. Property iii) 
says that this cost is zero exactly where continuation is optimal, so that in the region where one 
does not exercise, the option is a fair game. The optimal rule is then transparent: continue while 
$\Snell > Y$, and stop the instant they meet.

\subsubsection{The Linear Complementarity Problem (LCP)}

Translating Theorem \ref{thm:snell} into the language of partial differential equations gives a free 
boundary problem, most rigorously expressed as a linear complementarity problem. Write 
$\mathcal{L}$ for the pricing operator associated with \eqref{eq:localvolforward},
\[
\mathcal{L}V := \frac{\partial V}{\partial t} + \frac{1}{2}\lvol^2(\Fwd,t)\,\Fwd^2 \frac{\partial^2 V}{\partial \Fwd^2} - \rf(t)\,V .
\]

A point of care is required with the obstacle, and it is the reason this subsection is not a routine 
transcription. Exercise delivers the \emph{spot}, not the forward: the holder of an American put 
receives $\Strike - \Spot_\tau$, and by \eqref{eq:spotforwardfactor} we have 
$\Spot_t = \Carry_t\,\Fwd_t$. Expressed in the forward variable, the obstacle is therefore
\begin{equation}\label{eq:americanobstacle}
\PayoffFn(t,\Fwd) = \bigl(\Strike - \Carry_t\,\Fwd\bigr)^{+} \quad \text{for a put}, \qquad 
\PayoffFn(t,\Fwd) = \bigl(\Carry_t\,\Fwd - \Strike\bigr)^{+} \quad \text{for a call},
\end{equation}
which is \textbf{explicitly time-dependent through $\Carry_t$}. Writing $(\Strike-\Fwd)^+$ instead, 
as is sometimes done, silently assumes $\Carry_t \equiv 1$ and is simply wrong: it misplaces the 
exercise boundary, and it does so by the widest margin in a name paying large discrete dividends --- precisely 
the case where early exercise matters most. Keeping $\Carry_t$ explicit costs nothing and preserves 
validity for time-dependent $\rf(t)$ and $\divy(t)$, for cash and proportional dividends, and for 
any mixture of them.

\begin{theorem}[American pricing as an LCP]\label{thm:LCP}
The value $V(t,\Fwd)$ of the American option satisfies, almost everywhere on 
$[0,T)\times(0,\infty)$, the three simultaneous conditions
\begin{equation}\label{eq:LCP_American}
    \begin{cases}
        \mathcal{L}V \le 0 & \text{(supermartingale property, Theorem \ref{thm:snell}(ii))} \\[6pt]
        V(t, \Fwd) \ge \PayoffFn(t,\Fwd) & \text{(domination, Theorem \ref{thm:snell}(i))} \\[6pt]
        \bigl(\mathcal{L}V\bigr) \cdot \bigl(V(t, \Fwd) - \PayoffFn(t,\Fwd)\bigr) = 0 & \text{(complementarity, Theorem \ref{thm:snell}(iii))}
    \end{cases}
\end{equation}
together with the terminal condition $V(T,\Fwd) = \PayoffFn(T,\Fwd)$ and $\PayoffFn$ as in 
\eqref{eq:americanobstacle}. \hfill $\blacksquare$
\end{theorem}

The correspondence with Theorem \ref{thm:snell} is exact, line for line, and this is the point of 
having stated the probabilistic version first. The complementarity condition encodes the dichotomy: 
at each point either $V > \PayoffFn$, in which case $\mathcal{L}V=0$ and the ordinary pricing 
equation holds, or $V = \PayoffFn$ and the option is exercised. The boundary separating the two 
regions is not known in advance --- it is part of the solution --- which is what makes the problem 
free-boundary rather than merely parabolic, and what rules out any closed-form solution.

\subsubsection{Calibration to Single Names: De-Americanisation}

A major practical hurdle arises during calibration: Dupire's formula (Theorem \ref{thm:dupire}) requires the partial derivatives of \emph{European} option prices, but single-name market quotes are American. We cannot apply Dupire directly to American prices. 

Therefore, calibrating the local volatility surface to American options requires a short, iterative mapping procedure known as \textbf{De-Americanisation}:
\begin{enumerate}
    \item \textbf{Quote Extraction:} Obtain the market prices of the liquid American options and calculate their Black-Scholes implied volatilities using an American pricer (such as a Binomial Tree or a PDE solver).
    \item \textbf{Pseudo-European Mapping:} Assume these implied volatilities approximate the volatilities of European options with the exact same strikes and maturities. Use these volatilities to generate a grid of "pseudo-European" option prices.
    \item \textbf{Local Volatility Construction:} Apply a smoothing parameterization (e.g., SVI) to the pseudo-European surface to ensure it is free of static arbitrage, and then apply Gatheral's form \eqref{eq:gatheral} of the Dupire formula to construct the local volatility surface $\lvol(\Fwd,t)$.
    \item \textbf{Repricing and Verification:} Substitute this $\lvol(F, t)$ back into the LCP (Equation \eqref{eq:LCP_American}) and solve it backwards using a finite difference method (such as Crank-Nicolson combined with a Projected Successive Over-Relaxation (PSOR) algorithm) to verify that the model accurately recovers the original American market quotes.
\end{enumerate}

This iterative procedure effectively bridges the theoretical gap between the European foundations of Local Volatility and the practical reality of American single-name markets, ensuring that our models reflect true market dynamics without introducing arbitrage.


However, the de-Americanisation process can be inverted, and this must be of no suprise at this stage, by the so called PINNs (\textit{Physics-Informed Neural Networks}). Both in the construction of the Local Volatility and in pricing, we can use neural networks to get approximate yet extremely fast solutions. The constant and volatile value of the market parameters suggests that approximate and untractable solutions may suffice the market necessities. The also retain a key property, which is allowing for smooth derivatives, what is key for the hedge and portfolio management.

We will dive further into this topic in the final chapter of the present dissertation.
\subsection{References to Sections 1.3 to 1.5}

The financial framework and pricing methodologies developed in the latter half of this chapter are built upon cornerstone texts in quantitative finance and stochastic analysis. The rigorous transition from abstract probability to continuous-time asset pricing, including the formulation of equivalent martingale measures and the \textit{No Free Lunch with Vanishing Risk} (NFLVR) condition, relies heavily on the definitive work by Delbaen and Schachermayer (\cite{DelbaenSchachermayer1994}). This abstract measure-theoretic approach is  bridged with continuous-time applications in the classic textbooks by Shreve (\cite{Shreve2004}), Øksendal (\cite{Oksendal2003}), and Björk (\cite{Bjork2009}).

The explicit formulation of diffusion models is rooted in the foundational breakthroughs of Black and Scholes (\cite{BlackScholes1973}) and Merton (\cite{Merton1973}). To address the empirical shortcomings of this framework, Local Volatility was formalized by Dupire (\cite{Dupire1994}), with its practical translation into the implied volatility domain extensively covered by Gatheral (\cite{Gatheral2006}). For an exhaustive, practitioner-oriented treatment of volatility smile dynamics and Stochastic Volatility models (such as the Heston model \cite{Heston1993}), Bergomi (\cite{Bergomi2016}) serves as an indispensable reference. 

